\documentclass[a4paper]{article}
\usepackage[T2A]{fontenc}
\usepackage[utf8]{inputenc}
\usepackage[russian]{babel}
\usepackage{amssymb}
\usepackage{amsmath}
\usepackage{fullpage}
\usepackage{rotating}
\usepackage{stmaryrd}
\usepackage{mathtools}
\usepackage{bbm}
\usepackage[unicode,colorlinks=true,pagebackref=true]{hyperref}
\usepackage[all]{xy}
\usepackage{microtype}
\usepackage{amsthm}

\usepackage{tikz}
\usetikzlibrary{arrows}
\usetikzlibrary{cd}
\usetikzlibrary{calc}
\usetikzlibrary{through}
\usetikzlibrary{tikzmark}
\usetikzlibrary{positioning}
\usetikzlibrary{decorations.pathreplacing}

\DeclareFontFamily{OT1}{pzc}{}
\DeclareFontShape{OT1}{pzc}{m}{it}{<-> s * [1.1] pzcmi7t}{}
\DeclareMathAlphabet{\mathpzc}{OT1}{pzc}{m}{it}

\DeclareMathOperator{\PGSp}{PGSp}
\DeclareMathOperator{\PGO}{PGO}
\DeclareMathOperator{\ind}{ind}
\DeclareMathOperator{\End}{End}
\DeclareMathOperator{\Gr}{Gr}
\DeclareMathOperator{\Cent}{Cent}
\DeclareMathOperator{\Spec}{Spec}
\DeclareMathOperator{\Map}{Map}
\DeclareMathOperator{\GW}{GW}
\DeclareMathOperator{\rk}{rk}
\DeclareMathOperator{\Br}{Br}
\DeclareMathOperator{\Aut}{Aut}
\DeclareMathOperator{\Hom}{Hom}
\DeclareMathOperator{\Lie}{Lie}
\DeclareMathOperator{\PGL}{PGL}
\DeclareMathOperator{\GL}{GL}
\DeclareMathOperator{\SL}{SL}
\DeclareMathOperator{\fchar}{char}
\DeclareMathOperator{\tr}{tr}
\DeclareMathOperator{\Iso}{Iso}
\DeclareMathOperator{\SB}{SB}
\DeclareMathOperator{\SO}{SO}
\DeclareMathOperator{\Spin}{Spin}
\DeclareMathOperator{\Isom}{Isom}
\DeclareMathOperator{\im}{im}
\DeclareMathOperator{\disc}{disc}
\DeclareMathOperator{\Stab}{Stab}
\DeclareMathOperator{\Nrd}{Nrd}
\DeclareMathOperator{\CH}{CH}
\DeclareMathOperator{\pt}{pt}
\DeclareMathOperator{\codim}{codim}
\DeclareMathOperator{\OGr}{OGr}
\DeclareMathOperator{\an}{an}
\DeclareMathOperator{\Cor}{Cor}
\DeclareMathOperator{\Mor}{Mor}
\DeclareMathOperator{\pr}{pr}
\DeclareMathOperator{\id}{id}
\DeclareMathOperator{\res}{res}
\DeclareMathOperator{\Trd}{Trd}
\DeclareMathOperator{\sing}{sing}
\DeclareMathOperator{\Bl}{Bl}
\DeclareMathOperator{\Frac}{Frac}

%\DeclareFontFamily{OT1}{pzc}{}
%\DeclareFontShape{OT1}{pzc}{m}{it}{<-> s * [1.2] pzcmi7t}{}
%\DeclareMathAlphabet{\mathpzc}{OT1}{pzc}{m}{it}
\newcommand{\categ}{\mathpzc}

\renewcommand{\O}{\mathrm{O}}
\renewcommand{\epsilon}{\varepsilon}
\newcommand{\ph}{\varphi}
\renewcommand{\emptyset}{\varnothing}

\newcommand{\term}{\textbf}
\newcommand{\rdfn}{=\mathrel{\mathop:}}
\newcommand{\dfn}{\mathrel{\mathop:}=}
\newcommand{\isom}{\simeq}

\newtheorem{theorem}{Теорема}[subsection]
\newtheorem{lemma}[theorem]{Лемма}
\newtheorem{proposition}[theorem]{Утверждение}
\newtheorem{corollary}[theorem]{Следствие}
\newtheorem{conjecture}[theorem]{Гипотеза}

\theoremstyle{definition}
\newtheorem{example}[theorem]{Пример}
\newtheorem{fact}[theorem]{Факт}
\newtheorem{remark}[theorem]{Замечание}
\newtheorem{exercise}[theorem]{Упражнение}
\newtheorem{definition}[theorem]{Определение}

\newcommand{\la}{\langle}
\newcommand{\ra}{\rangle}
\newcommand{\lAngle}{\langle\!\langle}
\newcommand{\rAngle}{\rangle\!\rangle}
\newcommand{\trleq}{\trianglelefteq}
\newcommand{\ol}{\overline}
\newcommand{\wt}{\widetilde}

\newcommand{\TBW}{\textbf{TBW}}

\begin{document}

\author{Иван Панин\and Виктор Петров}
\title{Мотивы Воеводского и арифметика линейных алгебраических групп
\footnote{Конспект лекций семинара весны 2012 года; предварительная
версия. Автор \TeX-версии~--- Александр Лузгарев.
Основано на конспекте Алексея Бешенова первых двух лекций.}}
\date{2012}
\maketitle

\section{Введение}

\subsection{Планы}

% 13.02.2012

Работа Панина и Пименова о квадратичных формах.

Простая формулировка. {\it Пусть $K = \mathbb{C} (z_1,\ldots,z_n)$ и
  $R \dfn \{ \frac{g(z)}{h(z)} \mid h(0) \ne 0 \} \subset K$~---
  регулярные функции в окрестности $0$. Пусть $u \in
  R^\times$. Рассмотрим уравнение

\[ T_1^2 + \cdots + T_k^2 = u. \]

\noindent (Предполагаем $k \ge 2$.) Если уравнение имеет решение в
$K$, то оно имеет решение и в $R$.}

\vspace{2em}

Интересующая нас задача: классифицировать простые алгебраические
группы над произвольным полем (или локальным регулярным кольцом). В
каком смысле~--- мы объясним. Что такое простые алгебраические
группы~--- это обсуждается в записках спецкурса. 

\vspace{2em}

Как все знают, над алгебраически замкнутыми полями классификацию
простых алгебраических групп дают диаграммы Дынкина. Среди них~---
четыре бесконечные серии, которым соответствуют следующие
присоединенные группы:

\begin{itemize}
\item $A_n$~--- $\PGL_{n+1}$.

\item $B_n$~--- $\SO_{2n+1}$.

\item $C_n$~--- $\PGSp_{2n}$.

\item $D_n$~--- $\SO_{2n}$.
\end{itemize}

Исключительные группы: $E_6$, $E_7$, $E_8$, $F_4$, $G_2$.

Имеется точная последовательность

\[ 1 \to \mu_n \to \SL_n \to \PGL_n \to 1. \]

<<Теорема типа Спрингера>>: \emph{пусть $G$ и $G^\prime$~--- группы типа $G_2$ над полем $K$. Пусть расширение $[L : K]$ нечетное. Тогда если $G_L \isom G_L^\prime$, то $G \isom G^\prime$.}

С точностью до каких-то тонкостей, имеем

\[ H^1 (K, G_0^{ad}) \approx \{\text{присоед. простые алг. группы над }K\text{ того же типа, что и }G_0\}. \]

Это соответствие функториально в том смысле, что расширение полей $L/K$ индуцирует
морфизм $H^1 (K,G^{ad}) \to H^1 (L,G^{ad})$.

Наша высокая цель~--- построить функтор $F$, сопоставляющий полям
абелевы группы с гомоморфизмом следа, так что конечное расширение
$[L:K]$ давало бы морфизм $F(L) \to F(K)$ и естественное
преобразование

\[ H^1 (K,G_0^{ad}) \to F(K). \]

Например, для $G_0 = \PGL_2 = \Aut (M_2)$ ответ такой:

\[ F\colon K \rightsquigarrow K_2^M (K) / 2, \]

\noindent где $K_2^M (K) = I^2(K) / I^3 (K)$.

\subsection{Теорема Меркурьева--Суслина и гипотеза Блоха--Като}

Пусть $A$~--- центральная простая алгебра над полем $K$ (более общее
понятие~--- \term{алгебра Азумайи}, \term{Azumaya algebra}). Ей
соответствует элемент $[A]$ в группе Брауэра $\Br (K)$.

\begin{itemize}
\item \textbf{Теорема Меркурьева} (1981)~--- изоморфизм ${}_2 \Br(K)
  \isom K_2^M / 2$, а также следствие про $[A] \in {}_2 \Br(K)$.

[\url{http://www.mathunion.org/ICM/ICM1986.1/Main/icm1986.1.0389.0393.ocr.pdf}]

[\url{http://www.math.ethz.ch/~knus/sridharan/merkurjev84.pdf}]

\item \textbf{Теорема Меркурьева--Суслина} (1982)~--- изоморфизм ${}_p
  \Br(K) \isom K_2^M / p$.

[L.H. Rowen, Ring theory, Vol. 2, \S 7.2]

\item \textbf{Гипотеза Блоха--Като} (<<norm residue isomorphism
  theorem>>)~--- $K_n^M/p (-) \isom H^n_\text{ét} (-, \mu_p^{\otimes
    n})$.

[\url{http://arxiv.org/abs/0805.4430}]
\end{itemize}

\subsection{Кольцо Гротендика--Витта}

$H^1 (K, \O_n)$~--- это классы изометрии невырожденных квадратичных
форм ранга $n$.

Имеется функтор в кольцо Витта
\[ H^1 (K, \O_n) \to W(K), \quad f \mapsto [f]. \]

Разберемся, что такое \term{кольцо Витта} $W(K)$. Его образующие~---
классы изометрии квадратичных форм над $K$, а соотношения выглядят так:
\begin{gather*}
[f] + [g] = [f \perp g],\\
[f]\cdot [g] = [f\otimes g],\\
\mathbb{H} = 0,
\end{gather*}
где $\mathbb{H}$~--- класс изометрии двумерной квадратичной формы $f(x,y)=xy$, а
$f \perp g$ имеет следующий смысл. Если $f$~--- квадратичная форма на
$V$, а $g$~--- квадратичная форма на $W$, то на $V\oplus W$ задается
квадратичная форма $(f\perp g) (u\oplus v) \dfn f(u) + g(v)$.

Имеется корректно определенный гомоморфизм
\begin{eqnarray*}
\rk\colon W(K) & \to & \mathbb{Z}/2,\\
{}[f] & \mapsto & \rk f \mod 2.
\end{eqnarray*}
$I \dfn \ker \rk$ называется \term{фундаментальным идеалом}.

\term{Кольцо Гротендика--Витта} $GW (K)$ определяется следующим
образом. В нем те же образующие, но  нет условия $[xy] = 0$. Сначала
определяется сложение
и умножение, делающее $GW (K)$ полукольцом:
\begin{gather*}
[f] + [g] = [f \perp g],\\
[f]\cdot [g] = [f\otimes g].
\end{gather*}
Потом мы берем группу Гротендика и получаем кольцо.

\subsection{$\mathbb{A}^1$-гомотопии и гипотеза Мореля}

[\url{http://mathunion.org/ICM/ICM1998.1/Main/00/Voevodsky.MAN.ocr.pdf}]

$\mathbb{A}^1$-гомотопическая категория пространств с отмеченными
точками над $K$.

Сфера $S^0 = \{ +, \bullet \}$ состоит из двух точек, из которых
$\bullet$~--- отмеченная.

Теорема Мореля (1999?) состоит в вычислении

\[ \pi_0^{stab} (S^0) \isom \GW (K) \]

Fabien Morel, On The Motivic $\pi_0$ of the Sphere Spectrum.\\
\url{http://dx.doi.org/10.1007/978-94-007-0948-5_7}

Желаемый функтор $F$ мог бы давать $H^1 (K,G) \to H_0^{\mathbb{A}^1} (B^\text{èt} G)$.

Аналог этого в топологии следующий. Пусть задано главное
$G$-расслоение $\mathfrak{g} \to X$ для клеточного пространства
$X$. Сопоставим ему отображение в классифицирующее пространство $X
\xrightarrow{f_\mathfrak{g}} BG$.

Существует соответствие между множеством классов изоморфности главных
$G$-расслоений над $X$ и гомотопическими классами $[X,BG]$.

Имеется инъекция

\[ [X,BG] = \pi_0 (\Map (X,BG)) \hookrightarrow H_0 (\Map (X,BG)). \]

Гипотеза Мореля заключается в том, что в алгебраической ситуации тоже получается инъекция

\[ [\Spec K, B^\text{èt} G] = \pi_0^{\mathbb{A}^1} (B^\text{èt} G)
\hookrightarrow H_0^{\mathbb{A}^1} (B^\text{èt} G). \]

\subsection{Формы Пфистера}

Рассмотрим фильтрацию на кольце Витта

\[ W(K) \supset I \supset I^2 \supset I^3 \supset \cdots \]

\begin{theorem}
$\bigcap_n I^n = \{0\}$.
\end{theorem}

Мы уже знаем, что $W(K) / I = \mathbb{Z}/2$.
$I/I^2$ как абелева группа порождается элементами вида $\lAngle a\rAngle \dfn x^2 - a\,y^2$ для некоторого $a\in K^\times$.
Более общо, $I^n/I^{n+1}$ порождается тензорными произведениями элементов
\[ \lAngle a_1, \ldots, a_n\rAngle = \lAngle a_1\rAngle\otimes\cdots\otimes\lAngle a_n\rAngle. \]
$\lAngle a_1, \ldots, a_n\rAngle$ называется \term{$n$-кратной формой Пфистера}.

\begin{example}
При $n = 1$ имеем $a \in K^*/(K^*)^2$; $K (\sqrt{a})$~--- квадратичное расширение.

При $n = 2$ символ $\lAngle a,b\rAngle$ есть норма алгебры кватернионов
$H = (a,b)$ над $K$.

При $n = 3$ символ $\lAngle a,b,c\rAngle$ есть норма алгебры октонионов
$(a,b,c)$ над $K$ (что соответствует группам типа $G_2$ над $K$).
\end{example}

\begin{theorem}[Арасон]
Если $[q] \in I^n$, то $\rk q \ge 2^n$.
Если при этом $\rk q = 2^n$, то $q \isom \alpha \cdot \lAngle
a_1,\ldots,a_n \rAngle$, где $\alpha \in K^\times$.
В частности, $\bigcap I^n = 0$.
\end{theorem}

\begin{itemize}
\item $e_0 (q) \dfn \rk q \mod 2$.

\item Если $e_0 = 0$, то $[q] \in I$. Определим $e_1 (q) \dfn [\![q]\!] \in I/I^2$. Этому соответствует $\lAngle a\rAngle$, где $a$~--- дискриминант $q$ (с точностью до знака?).

\item Если $e_1 = 0$, то $e_2 (q) \dfn [\![ q ]\!] \in I^2/I^3$.
Форме $q$ можно сопоставить $C_0^+ (q)$, четную положительную часть алгебры
Клиффорда, это будет центральная простая алгебра. Имеем $[C_0^+ (q)]
\in {}_2 \Br (K)$. По теореме Меркурьева, это сумма
\[ [(a_1,b_1)]\,[(a_2,b_2)]\cdots [(a_k,b_k)] \]

\[ \lAngle a_1,b_1\rAngle + \lAngle a_2,b_2\rAngle + \cdots +\lAngle a_k,b_k\rAngle. \]

\item Если $e_2 (q) = 0$, то можно определить $e_3 (q)$~---
  \term{инвариант Арасона}.
\end{itemize}

\subsection{Торсоры}

Пусть $G$~--- простая алгебраическая группа над $K$.

\term{$G$-торсором} называется многообразие $X$ над $K$, такое что
\begin{itemize}
\item определено действие $G\times X\to X$;

\item над алгебраическим замыканием $\overline{K}$ имеется изоморфизм
  $X_{\overline{K}} \isom G_{\overline{K}}$ (как многообразий с
  $G$-действием).
\end{itemize}

Раньше торсоры назывались <<главными однородными пространствами>>
(principal homogeneous space).

\begin{example}
Действие $G$ сдвигами на себе дает \term{тривиальный $G$-торсор}.
\end{example}

По определению, $H^1 (K;G)$ есть множество классов изоморфности
$G$-торсоров с отмеченной точкой (тривиальный $G$-торсор).

\begin{example}
Зафиксируем $a\in K$.
Для каждой $K$-алгебры $R$ положим
$\mu_2(R) = \{ x\in R\mid x^2 = 1 \}$, $X(R) \dfn \{ y\in R\mid y^2 = a\}$.
Получаем схемы $\mu_2$ и $X$, причем
$\mu_2$ действует на $X$ умножением:
если $y^2 = a$, $x^2 = 1$, то $(x\,y)^2 = a$.
\end{example}

$X$~--- тривиальный $G$-торсор iff у него есть рациональная точка: $X
(K) \ne \emptyset$.

Если $G$~--- абелева группа, то на торсорах имеется сложение. При этом
$H^1 (K,\mu_2) \isom K^* / (K^*)^2$ как абелева группа. И вообще $H^1
(K,\mu_n) \isom K^* / (K^*)^n$.

\subsection{Скрученные формы}

Напомним, что $\O_{2n} = \Aut (q_{split})$, где $q_{split} = x_1\,y_1 + \cdots +
x_n\,y_n$~--- \term{расщепимая форма} (от переменных $x_1,\dots,x_n,y_1,\dots,y_n$.

$H^1 (K, \O_{2n})$ можно отождествить с множеством классов изометрии
невырожденных квадратичных форм ранга $2n$. 
Действительно, пусть $q$~--- квадратичная форма,
Мы утверждаем, что $\Iso (q_{split},
q)$ есть искомый торсор: на нем действует $\O_{2n}$.
Здесь $\Iso(\ph,\psi)$ обозначает функтор изоморфизмов между квадратичными
формами $\ph$ и $\psi$; более точно,
$\Iso(\ph,\psi)(R) = \{f\colon\ph_R\to\psi_R\mid\mbox{$f$~--- изоморфизм}\}$.
Над алгебраически
замкнутым полем $q$ изоморфна $q_{split}$, и получается $\Iso
(q_{split},q_{split})=\Aut (q_{split})=\O_{2n}$.

Пусть $A$~--- некоторая алгебраическая структура над полем $K$
(например, квадратичное пространство, конечномерная ассоциативная
алгебра, конечномерная неассоциативная алгебра). Тогда
\term{скрученная форма $A^\prime$} для $A$ есть такая структура над
$K$, что при переходе к алгебраическому замыканию
$A^\prime_{\overline{K}} \isom A_{\overline{K}}$.

\noindent\textbf{Теорема}. $H^1 (K, \Aut(A))$ есть множество классов
изоморфности скрученных форм $A$ над $K$.

Изоморфизм такой:
\[ A' \xmapsto{\sim} \Iso (A,A^\prime). \]
На $\Iso (A,A^\prime)$ есть структура алгебраического многообразия.

\vspace{2em}

\noindent\textbf{Замечание}. Пусть $X$~--- проективное многообразие
над $K$. Теорема (Гротендик): \emph{функтор $U \mapsto \Aut_U (X\times
  U)$ представим в схемах}; то есть, существует схема $R$ такая, что
$\Aut_U(X\times U)$ естественно изоморфно $\Hom(U,R)$.

\vspace{2em}

Контрпример: $\Aut (\mathbb{A}^n)$ не конечномерно.

Пример: $A \dfn M_n (K)$. $\Aut (A) = \PGL_n$.\\
$H^1 (K, M_n(K))$~--- это скрученные формы $M_n (K)$, то есть
центральные простые алгебры размерности $n^2$, взятые с точностью до
изоморфизма.

$\Aut (\mathbb{P}^{n-1}) = \PGL_n = \GL_n / \mathbb{G}_m$.

Автоморфизмы сохраняют ранг.

$H^1 (K, \PGL_n)$ есть множество скрученных форм $\mathbb{P}^{n-1}$
над $K$ = \term{многообразия Севери--Брауэра}.

\[ A \mapsto \SB (A) = \{ \text{левые идеалы $I\trleq A$}\mid
\dim_K(I)=n \} \]

Пример при $n=2$: кватернионы $A = (a,b)$.

$\beta\,u + \gamma\,v + \delta\,u\,v$. Имеем векторное пространство
$u,v,uv$. Условие $\{ \text{норма} = 0 \}$ задает проективное
подмногообразие в $\mathbb{P}^2$.

$x^2 - a\,y^2 - b\,z^2 = 0$~--- коника.

\begin{eqnarray*}
\PGL_2 & \isom & \SO_3, \\
\{\text{кватернионы}\} & \isom & \{\text{формы ранга }3\text{ с
  трив. дискриминантом}\}.
\end{eqnarray*}

\subsection{Точные последовательности алгебраических групп}

Точность последовательности алгебраических групп над $K$

\[ 1 \to C \to H \to G \to 1 \]

\noindent означает следующее:

\begin{enumerate}
\item $C$~--- алгебраическая подгруппа в $H$.

\item После расширения скаляров $H (\overline{K})\to G (\overline{K})$
  является сюръекцией
  \emph{над алгебраическим замыканием поля $K$}.

\item $C = \ker (H\to G)$, $C(R) = \ker (H(R) \to G(R))$.
\end{enumerate}

\begin{example}
Следующая последовательность алгебраических групп точна в указанном смысле:

\[ 1 \to \mu_2 \to \mathbb{G}_m \to \mathbb{G}_m \to 1, \]

\noindent где $\mathbb{G}_m (K) \dfn \{ (x,y)\in K^2 \mid x\,y = 1 \}$, и
отображение $\mathbb{G}_m \to \mathbb{G}_m$ есть $x \mapsto x^2$ (это
сюръекция над алгебраическим замыканием).
\end{example}

\begin{example}
Следующая последовательность точна:

\[ \mu_2 (K) \to K^\times \to K^\times \to H^1 (K,\mu_2) \to H^1
(K,\mathbb{G}_m) \to H^1 (K,\mathbb{G}_m). \]
\end{example}

\noindent (Отображение $K^\times \to K^\times$ есть $x \mapsto x^2$.)

\begin{theorem}[Теорема Гильберта 90]

\begin{gather*}
H^1 (K,\mathbb{G}_m) = \{\bullet\},\\
H^1 (K,\GL_n) = \{\bullet\}.
\end{gather*}
\end{theorem}

Из точности последовательности выше и теоремы 90 получается

\[ H^1 (K,\mu_2) \isom K^\times / (K^\times)^2. \]

\begin{example}
Точная последовательность
\[ 1 \to \SL_n \to \GL_n \xrightarrow{\det} \mathbb{G}_m \to 1. \]
приводит к последовательности
\[ \GL_n (K) \xrightarrow{\det} \mathbb{G}_m (K) \to H^1 (K,\SL_n) \to H^1 (K,\GL_n). \]

Здесь $H^1 (K,\SL_n) = \{\bullet\}$ и $H^1 (K,\GL_n) = \{\bullet\}$.
\end{example}

\begin{example}
\[ 1 \to \mu_n \to \SL_n \to \PGL_n \to 1. \]

\[ \mu_n (K) \to \SL_n (K) \to \PGL_n (K) \to K^\times / (K^\times)^2 \to \{\bullet\} \to H^1 (K,\PGL_n). \]

\[ 1 \to \mathbb{G}_m \to \GL_n \to \PGL_n \to 1. \]
\end{example}

\begin{example}
\begin{eqnarray*}
\SL_n & \to & \PGL_n,\\
g & \mapsto & (x \mapsto g\,x\,g^{-1}) \in \Aut (M_n).
\end{eqnarray*}

Это сюръекция алгебраических групп, но не сюръекция на точках.

\[ 1 \to \mu_n \to \SL_n \to \PGL_n \to 1. \]
\end{example}

\begin{theorem}
Если имеется точная последовательность $1 \to C \to H \to G \to 1$, то
возникает точная последовательность множеств с отмеченной точкой

\[ 1 \to C(K) \to H(K) \to G(K) \to H^1 (K,C) \to H^1 (K,H) \to H^1 (K,G). \]
\end{theorem}

См. книгу Серра <<Когомологии Галуа>>.

% 27.02.2012

\subsection{Вторые когомологии}

\begin{itemize}
\item Напомним, что $H^1 (F,G)$~--- множество $G$-торсоров. \emph{Если $G$
    коммутативна}, то это аффинная алгебраическая группа.

(Как в этом случае умножаются торсоры?~--- Что-то типа $E_1 \mathop{*}
E_2 = (E_1 \times E_2) / ((e_1,e_2) = (g\,e_1,g\,e_2))$.)

\item $H^0 (F,G)$~--- это функтор $F \rightsquigarrow G(F)$,
  т.е. функтор точек. \emph{Если $G$ коммутативна}, то $H^i (F,G)$
  можно определить как $i$-й производный функтор. При $i = 1$ это
  совпадает с первым определением.
\end{itemize}

\begin{theorem}
Пусть имеется точная последовательность $1 \to C \to G \to H \to
1$. Предположим, что $C \le \Cent (G)$. Тогда точная
последовательность продолжается до вторых когомологий:

\[ 1 \to C(F) \to G(F) \to H(F) \to H^1 (F,C) \to H^1 (F,G) \to H^1 (F,H) \to H^2 (F,C). \]
\end{theorem}

\begin{example}
$H^2 (F, \mathbb{G}_m)=\Br (F)$~--- \emph{группа Брауэра} поля $F$:
она состоит из классов эквивалентности $[A]$ центральных простых
алгебр $A$ над $F$; умножение выглядит так: $[A]\cdot
[B]=[A\otimes_FB]$.

Пусть $X$~--- квазипроективное многообразие. Тогда $H^2
(X,\mathbb{G}_m)_{tors} = \Br (X)$ (\emph{теорема Габбера} (Gabber)).
(Загадочное замечание:
подразумевается топология fppf, а для этальной топологии в определении
торсора вместо $\overline{F}$ нужно взять $F^{sep}$.)
\end{example}

\begin{example}[Топологический аналог]
Пусть $X$~--- хорошее топологическое пространство (например, область в
$\mathbb R^n$, многообразие или CW-комплекс).

Пусть $G$~--- топологическая группа (например, $S^1$, $S^3$, $\SL_2
(\mathbb{C})$, $\O_n (\mathbb{C})$).

Имеется левое действие $G \times (G\times X) \to (G\times X)$, $g_1
\cdot (g_2,x) \mapsto (g_1\,g_2, x)$.

Левое действие послойно и свободно на скрученной форме $G\times
\mathcal{G} \to \mathcal{G}$.

Для всех $x \in X$ возникает действие $G \times \mathcal{G} (x) \to
\mathcal{G} (x)$. Здесь $\mathcal{G} (x) \isom G$, и этот изоморфизм
зависит от $x$.

$(\mathcal{G}, G\times \mathcal{G} \to \mathcal{G})$ в топологии
называется \term{главным $G$-расслоением} (\term{principal
  $G$-bundle}).

$\mathcal{G}/G = X$.
\end{example}

\begin{example}
$\mathbb{C}^\times = \GL_1 (\mathbb{C}) = \Aut (\mathbb{C}^1)$.

Пусть $L \to X$~--- комплексное линейное расслоение, $z (X)$~---
нулевое сечение.

Рассмотрим отображение $\mathbb{C}^\times \times (L - z(X))\to (L -
z(X))$, $(\lambda, v)\mapsto \lambda v$. Имеем изоморфизм $L(x) -
0\isom \mathbb{C}^\times$, зависящий
от $x$.

\begin{itemize}
\item Тогда $H^1 (X,\mathbb{C}^\times)$~--- классы изоморфизма
  $\mathbb{C}^\times$-торсоров над $X$. Они соответствуют линейным
  расслоениям над $X$: расслоению $L$ соответствует описанный выше
  торсор $L - z(X)$, и по торсору $\mathcal G^\times$ можно построить
  расслоение $\mathcal L$.

\item Таким образом, мы видим, что $H^1 (X, \Aut (\mathbb{C}^1))$~---
  это скрученные формы расслоения
  $\mathbb{C}\times X$ над $X$.

\item Аналогично, $H^1 (X, \Aut (\mathbb{C}^n))$~--- это (1) скрученные формы
  расслоения $\mathbb{C}^n\times X$ над $X$, то есть (2) векторные
  расслоения над $X$ со слоем $\mathbb{C}^n$ (с точностью до изоморфизма).

\item Пусть $\Aut_{\mathbb{C}} (\mathbb{C}^{2n}, \sum u_i\,v_i)$~---
  автоморфизмы, сохраняющие квадратичную форму.
  Тогда \[H^1 (X, \Aut_{\mathbb{C}} (\mathbb{C}^{2n}, \sum
  u_i\,v_i))\]--- это (1) скрученные формы расслоений вида
  $(\mathbb{C}^n\times X, \sum u_i,v_i) \to X$, то есть (2) векторные
  расслоения $E \to X$ со слоем $\mathbb{C}^n$ и с квадратичной формой в
  слоях.

\item Рассмотрим $\Aut (M_n (\mathbb{C})) = \PGL_n
  (\mathbb{C})$. Тогда $H^1 (X, \Aut (M_n (\mathbb{C})))$~--- это
  скрученные формы расслоений вида $M_n (\mathbb{C})\times X \to
  X$. Например, по каждому расслоению $E\to X$ можно построить
  расслоение $\End(E)\to X$, и послойно $\End(E)(x)=\End(E(x))$. Но
  бывают и расслоения, не изоморфные никакому $\End(E)\to X$~--- это
  нетривиальные топологические алгебры Адзумайи.
\end{itemize}
\end{example}

Имеется точная последовательность

\[ 1 \to \mathbb{C}^\times \to \GL_n (\mathbb{C}) \to \PGL_n (\mathbb{C}) \to 1. \]

Отсюда получается точная последовательность

\begin{gather*}
1 \to \Gamma (X, \mathbb{C}^\times) \to \Gamma (X, \GL_n (\mathbb{C}))
\to \Gamma (X, \PGL_n (\mathbb{C})) \to \\ \to H^1 (X,\mathbb{C}^1)
\to H^1 (X,\GL_n (\mathbb{C})) \to H^1 (X, \PGL_n (\mathbb{C})) \to
H^2 (X, \mathbb{C}^\times).
\end{gather*}

\term{Топологическая группа Брауэра} есть $\Br_{top} (X) \dfn
H^2_{top} (X, \mathbb{C}^\times)$.

\begin{example}
Мы утверждаем, что
\[ H^2 (X, S^1) \twoheadrightarrow H^3 (X, \mathbb{Z})_{tors}. \]

Заметим, что
$\mathbb{C}^\times \isom S^1 \times \mathbb{R}$. Имеем точную
последовательность 
\[ 0 \to \mathbb{Z} \to \mathbb{R} \to S^1 \to 0, \]
откуда получаем точную последовательность
\[ H^2(X,\mathbb{Z})\to H^2(X,\mathbb{R}\to H^2(X,S^1)\to
H^3(X,\mathbb Z)\to H^3(X,\mathbb R)=H^3(X,\mathbb Z)\otimes \mathbb R.\]
Обозначим отображение $H^3(X,\mathbb Z)\to H^3(X,\mathbb
Z)\otimes\mathbb R$ через $\alpha$. Тогда связывающий гомоморфизм
дает нам отображение $H^2(X,S^1)\to\ker(\alpha)=H^3(X,\mathbb Z)_{tors}$.

На самом деле,
\[ \Br_{top} (X) = H^3 (X,\mathbb{Z})_{tors}. \]
Нечто такое написано как определение (у Серра? Гротендика?).
\end{example}

А какие нам известны нетривиальные скрученные формы алгебры $M_n (K)$?
Так это и есть центральные простые алгебры.

Имеем точную последовательность $1 \to \mathbb{G}_m \to \GL_n \to
\PGL_n \to 1$, откуда

\[ H^1 (F,\GL)_n \to H^1 (F, \PGL_n) \to H^2 (F,\mathbb{G}_m). \]

При этом
$H^1 (F,\GL)_n = \{ \bullet \}$, и $H^1 (F, \PGL_n)$~--- центральные
простые алгебры степени $n$. Это отображение дает изоморфизм % ???

\begin{eqnarray*}
\Br (F) & \isom & H^2 (F,\mathbb{G}_m),\\
A & \mapsto & [A].
\end{eqnarray*}
% тут пропущен кусок про умножение в H^1???

Имеется точная последовательность
\[ 1 \to \mu_n \to \mathbb{G}_m \to \mathbb{G}_m \to 1, \]
где $\mathbb{G}_m \to \mathbb{G}_m$~--- отображение $x \mapsto x^n$.

Получаем точную последовательность
\[ \xymatrix{ H^1 (F,\mathbb{G}_m) \ar[r]\ar@{=}[d] & H^2 (F,\mu_n)
  \ar[r] & H^2 (F,\mathbb{G}_m) \ar[r]\ar@{=}[d] & H^2
  (F,\mathbb{G}_m)\ar@{=}[d] \\
\{ \bullet \} & & \Br (F) \ar[r]^{\cdot n} & \Br (F) }, \]
откуда
$H^2 (F,\mu_n) = {}_n \Br (F)$.

Еще один пример: пусть $\fchar F \ne 2$. Имеется точная последовательность
\[ 1 \to \SO_n \to \O_n \xrightarrow{\det} \mu_2 \to 1. \]

Тогда
\[ \O_n (F) \twoheadrightarrow \mu_2 (F) \to H^1 (F,\SO_n) \to H^1 (F,\O_n) \xrightarrow{\disc} H^1 (F,\mu_2)=F^*/(F^*)^2. \]

Отсюда $H^1 (F,\SO_n)$ (невырожденные квадратичные формы дискриминанта
1)~--- подмножество в $H^1 (F,\O_n)$ (невырожденные квадратичные формы
ранга $n$ с точностью до изометрии).

Это дает нам инвариант $\disc=e_1\colon I\to I/I^2\isom F^*/(F^*)^2$.

Еще один пример:
\[ 1 \to \mu_2 \to \Spin_n \to \SO_n \to 1. \]

\[ \Spin_n (F) \to \SO_n (F) \xrightarrow{N} H^1 (F,\mu_2) \xrightarrow{0} H^1 (F,\Spin_n) \to H^1 (F,\SO_n) \xrightarrow{e_2} H^2 (F,\mu_2). \]

Здесь
\begin{itemize}
\item $H^2 (F,\mu_2) = {}_2 \Br (F)$.

\item $N$~--- \term{спинорная норма}. А именно, каждый элемент $\SO_n$
  раскладывается в произведение отражений $g = S_{v_1} \cdots
  S_{v_{2k}}$; тогда $N(g) \dfn q (v_1) \cdots q (q_{2k})
  \pmod{(F^\times)^2}$.

\item Отображение $\Spin_n(F)\to\SO_n(F)$ уже не обязательно является
  сюръективным.

\item Мы не знаем, что такое $H^1(F,\Spin_n)$. Отображение
  $H^1(F,\mu_2)\to H^1(F,\Spin_n)$ равно  $0$ по теореме Эйхлера.
\end{itemize}

Самая правая стрелка в этой длинной последовательности дает нам инвариант
$e_2\colon I^2 \to I^2 / I^3 = {}_2 \Br (F)$. Его можно описать так:
по форме $q$ можно построить алгебру Клиффорда $C(q)$ с четной частью
$C_0(q)$. Тогда $e_2$ сопоставляет форме $q\in I^2$ класс $[C^+_0(q)]$
в $\Br(F)$.

Пусть $E$~--- левый $G$-торсор, $G$ действует на $X$ справа.
Рассмотрим скрученную форму $X$

\[ {}_EX \dfn (X\times E)/ \! {(x,e) \sim (x\,g^{-1}, g\,e)}. \]

На ней действует ${}_EG$. Действительно,
${}E_G=\Aut_{G-\text{торс}}(E)$~--- автоморфизмы $E$ как
$G$-торсора. ${}_EG$ является группой (тут $G$ действует сопряжениями
на себе). 

% тут пропущен кусок???

\begin{example}
Рассмотрим $H^1 (F,\O_n)$.

$E \in H^1 (F,\O_n)$ задается квадратичной формой $q$, и $q$ должна
быть формой расщепимой квадратичной формы $q_0$. При этом $E = \Isom
(q_0, q)$.

$O (q_0)$ действует на квадрике $Q_0 \dfn \{ q_0 = x_1\,y_1 + \cdots +
x_n\,y_n = 0 \} \subseteq \mathbb{P}^{2n-1}$.

После подкрутки: ${}_EQ_0 = \{ q = 0 \}$ и на ${}_EQ_0$ действует
группа $O(q)={}_EO(q_0)$.
\end{example}

\subsection{Многообразия Севери--Брауэра}

\begin{example}
Рассмотрим $H^1 (F, \PGL_n)$.

Торсор $E\in h^1(F,\PGL_n)$ задается центральной простой алгеброй $A$
степени $n$:
$E = \Isom_{F\text{-}\categ{Alg}} (M_n, A)$. Напомним, что $\PGL_n=\Aut(M_n)$.

%Что такое $\mathbb{P}^{n-1}$?

Каждому вектору $v \in \mathbb{A}^n - \{ 0 \}$ соответствует правый
идеал $\{x\mid \im x\leq \left<v\right>\}$ в $M_n$. Множество всех
идеалов, получающихся таким образом~--- это в точности множество
правых идеалов размерности $n$.

%$\left<v\right> \rightsquigarrow \text{правый идеал в } M_n \{ x \mid
%\im x \subseteq \left<v\right> \}$.

${}_E\mathbb{P}^{n-1}$~--- Множество правых идеалов размерности $n$ в
$A$~--- \term{многообразие Севери--Брауэра} $\SB (A)$.

Уравнения:
\[ \SB (A) \dfn \{ W \subset A \mid W\cdot A \subseteq A \}. \]

Таким образом,
\[ \SB (A) \hookrightarrow \Gr (n,A) = \Gr (n,n^2). \]

\[ \xymatrix{
A\times \SB(A) & A\times \Gr (n,n^2) \\
\left.\tau\right|_{\SB (A)}\ar[d]\ar@{^(->}[u] & \tau_n\ar[d]\ar@{^(->}[u] & W\ar[d] \\
\SB (A)\ar@{^(->}[r] & \Gr (n,n^2) & \{ w \}
} \]
\end{example}

\begin{lemma}
$\End_{\SB (A)} (J_A) \isom A$ (эндоморфизмы расслоения).
\end{lemma}
Поэтому два описания $H^1(F,\PGL_n)$~--- как алгебры Адзумайи и как
формы $\mathbb P^{n-1}$~--- эквивалентны.

\begin{proposition}
Подрасслоение $\left.\tau\right|_{\SB(A)}$ выдерживает правое
$A$-действие на $A\times \SB(A)$.

$J_A \dfn \left.\tau\right|_{\SB (A)}$.
\end{proposition}

\begin{example}
$\Gr (K,n)$ (линейные $k$-мерные подпространства в $\mathbb{A}^n$).

Если $U$~--- $k$-мерное подпространство в $\mathbb{A}^n$, то
$\{ x \mid \im x \subseteq U \}$~--- правый идеал в $M_n$ размерности
$kn$.

$\SB_k (A)$~--- обобщенное многообразие Севери--Брауэра~--- многообразие
правых идеалов размерности $k$.
\end{example}

$\Gr (k,n) \isom \Gr (n-k, n)$ (напомним, что это не канонический
изоморфизм). Аналог этой двойственности: $\SB_k (A) \isom \SB_{n-k}
(A^{op})$.

\begin{proposition}
$\SB (A) (F) \ne \emptyset \Rightarrow A \isom M_n$.

$\SB_k (A) (F) \ne \emptyset \Rightarrow \ind A \mid k$.
\end{proposition}

(Напомним, что такое $\ind$. Для центральной простой алгебры $A$ имеем
$A \isom M_m (D)$, где $D$~--- тело. $m \cdot \deg D = n$. $\deg D \rdfn
\ind A$, где $\deg D \dfn \sqrt{\dim D}$.)

Скрученные формы $\mathbb{P}^{n-1}$~--- это скрученные формы
$M_n$. Имеем $\Aut (\mathbb{P}^{n-1}) = \PGL_n$.

Предположим $\fchar F \ne 2$.

$\Aut(q_0) = O(q_0)$.

$\Aut (Q_0)^+ = \PGO (q_0)$, где $Q_0$~--- квадрика $\{ q = 0 \}$.

$H^1 (F, \PGO (q_0))$~--- классы $(A,\sigma)$ изоморфности центральных
простых алгебр $A$ с ортогональной инволюцией $\sigma$.


\[ \{ \text{правые идеалы }I\text{ в }(A,\sigma)\text{ размерности }\deg A \mid \sigma (I) \cdot I = 0 \} \rdfn X_{(A,\sigma)} \hookrightarrow \SB(A). \]

При поднятии до $\overline{F}$ получаем:
\[ Q_{\sigma (q_{\overline{F}})} \hookrightarrow \mathbb{P}^{\deg A -
  1}_{\overline{F}} = \SB (A) \otimes_F \overline{F}. \]

Вложение $X_{(A,\sigma)} \hookrightarrow \SB(A)$ есть аналог вложения
квадрики в проективное пространство.

% 05.03.2012

\section{Проективные однородные многообразия}

\subsection{Первые примеры}

Еще раз про аналогию с топологией:

$E\to X$~--- торсор на топологическом пространстве $X$, $G$ действует
на $E$. Существует покрытие $\{U_i\}$ пространства $X$ такое, что
\[\begin{xymatrix}{U_i\times G\isom E|_{U_i}\ar[r]\ar[d] & E \\
U_i\ar@{^(->}[r] & X}\end{xymatrix}\]

У нас: возьмем $X=\Spec K$. Пусть $E\to\Spec K$~--- торсор. Существует
расширение полей $L/K$ такое, что торсор $E_L\to\Spec L$ изоморфен
торсору $G_L\to\Spec L$.

Мы хотим описать $H^1(K,G)$. Стратегия: для торсора $E$ и (гладкого
проективного) $G$-многообразия $X$ мы определили ${}_EX$~---
${}_EG$-многообразие (снова гладкое проективное), которое называется
\emph{скрученной формой $X$},  то есть,
\begin{itemize}
\item $E_{\overline{K}}\isom
G_{\overline{K}}$ как $G_{\overline{K}}$-многообразие,
\item $({}_EG)_{\overline{K}}\isom G_{\overline{K}}$ как алгебраическая
группа,
\item $({}_EX)_{\overline{K}}\isom X_{\overline{K}}$ как
  $G_{\overline{K}}$-многообразие.
\end{itemize}

\begin{example}
Пусть $A\in H^1(K,\PGL_n)$; то есть, $A$~--- центральная простая
алгебра степени $n$. Положим $E=\Isom(M_n,A)$, $X=\mathbb P^{n-1}$,
$G=\PGL_n$. Тогда ${}_EG=\Aut(A)$, ${}_EX=\SB(A)$~--- многообразие
правых идеалов в $A$ размерности $n$. Заметим, что $\SB(A)(K)$ непусто
тогда и только тогда, когда $A\isom M_n$. Вообще, свойства многообразия
$\SB(A)$ отражают свойства исходного торсора.
\end{example}

\begin{example}
Пусть $G=\O_n$, $q\in H_1(K,\O_n)$~--- невырожденная квадратичная
форма ранга $n$.

$E=\Isom(q_0,q)$, где $q_0$ расщепима (то есть, имеет
вид $\langle 1,-1\rangle\perp\dots\perp\langle 1,-1\rangle$ плюс,
возможно, слагаемое $\langle 1\rangle$).

$X=\{q_0=0\}$ в проективном смысле. Тогда $Q={}_EX=\{q=0\}$. Заметим,
что  $Q(K)$ непусто тогда и только тогда, когда форма $q$ изотропна, то
есть, $q\isom\langle 1,-1\rangle\perp q'$. Этот факт остается верным
при любом расширении $L/K$: $Q(L)$ непусто тогда и только тогда, когда
форма $q_L$ изотропна.
\end{example}

\begin{fact}
Пусть $q$ имеет вид $\lAngle a_1,\dots,a_k\rAngle=\la
1,-a_1\ra\otimes\dots\otimes\la 1,-a_k\ra$. Тогда
\begin{multline}
\text{
Форма $q_L$ изотропна тогда и только тогда, когда она расщепима (то
есть,} \\
\text{раскладывается в сумму форм вида $\la 1,-1\ra$).}\tag{*}
\end{multline}
Наоборот, если $\dim q$ четна и (*) выполнено для любого расширения
полей, то $q$ пфистерова с точностью до скаляра. Если же $\dim q$
нечетна, то $q\perp\la 1\ra$ пфистерова с точностью до скаляра.
\end{fact}

Таким образом, от торсора $E$ можно переходить к многообразию ${}_EX$
(и можно варьировать $X$), смотреть на его инварианты (в смысле
алгебраической геометрии) и получать отсюда информацию об инвариантах
торсора.

Пусть $X$~--- гладкое проективное многообразие. Мы ограничимся
случаем, когда $X$ \emph{однородное}, то есть, $G(\overline{K})$
действует на $X(\overline{K})$ транзитивно (заметим, что это означает,
что отображение $G\times X\to X\times X$, $(g,x)\mapsto (gx,x)$
сюръективно как пучок, а не в категорном смысле; категорное понятие
эпиморфизма не подходит для наших целей: например, отображение
$\Spec{\mathbb Q}\to\Spec{\mathbb Z}$ сюръективно в категории схем).
Неформально говоря, у $G$ на $X$ одна орбита.
Тогда $X$ называется \term{проективным однородным многообразием}.

Как строить проективные однородные многообразия? Пусть $G$~---
расщепимая группа, $V$~--- неприводимое представление (в положительной
характеристике нужно действовать осторожене). Рассмотрим $\mathbb
P(V)$~--- многообразие прямых в $V$, проходящих через $0$. У группы
$G$ есть ровно одна замкнутая орбита на $\mathbb P(V)$~--- это и есть
наше $X$. На самом деле, все проективные однородные многообразия так
получаются (но не обязательно единственным образом).

\begin{example}
$G=\SL_n$ действует на $V=K^n$ (имеется в виду обычное, \emph{векторное}
представление). Пусть $u$, $v$~--- два вектора. Можно ли найти $g$
такое, что $\la gu\ra=\la v\ra$ (здесь через $\la
x\ra$ мы обозначаем прямую, натянутую на $x$)? Ответ~--- можно,
если $u$ и $v$ отличны от $0$. Значит, в $K^n$ есть две орбиты
действия группы $G$: $\{0\}$ и $\{v\mid v\neq 0\}$. После
проективизации в $\mathbb P(K^n)=\mathbb P^{n-1}$ остается только одна
орбита.
\end{example}

\begin{example}
$G=\SL_n$ действует на $V=\Lambda^k(K^n)$, $k=1,\dots,n-1$. На
неразложимых поливекторых орбит много, но на разложимых действие
транзитивно. Свойство <<быть разложимым>> определяется уравнениями
Плюккера. Орбита в $\mathbb P(\Lambda^k(K^N))$~--- это $\Gr(k,n)$.

Пусть, к примеру, $n=4$, $k=2$. Диаграмма Хассе весов нашего
представления выглядит так:
\[
\begin{tikzpicture}[scale=1.0,thin,
aln/.style={above left=-2pt},
arn/.style={above right=-2pt},
bln/.style={below left=-2pt},
brn/.style={below right=-2pt},
every label/.style={above=2pt}]
\def\sm{0.7}
\coordinate (1) at ($\sm*(0, 0)$);
\coordinate (2) at ($\sm*(1.4, 0)$);
\coordinate (3) at ($\sm*(2.4, 1)$);
\coordinate (4) at ($\sm*(2.4, -1)$);
\coordinate (5) at ($\sm*(3.4, 0)$);
\coordinate (6) at ($\sm*(4.8, 0)$);
\draw (1)--node[above] {$2$}(2);
\draw (2)--node[aln] {$1$}(3);
\draw (2)--node[bln] {$3$}(4);
\draw (3)--node[arn] {$3$}(5);
\draw (4)--node[brn] {$1$}(5);
\draw (5)--node[above] {$2$}(6);
\foreach \point in
{1,2,3,4,5,6}
  {
     \fill [white] (\point) circle (2.0pt);
     \draw [black] (\point) circle (2.0pt);
  }
\end{tikzpicture}
\]
Разложимый тензор задается двумя векторами. Запишем их в базисе
$(e_1,e_2,e_3,e_3)$:
$(a_1e_1+a_2e_2+a_3e_3+a_4e_4)\wedge(b_1e_1+b_2e_2+b_3e_3+b_4e_4)$. Обозначим
координату тензора $x$ при бивекторе $e_i\wedge e_j$ через
$x_{ij}$. Тогда разложимость $x$ равносильно обращению в $0$ выражения
$x_{12}x_{34}-x_{13}x_{24}+x_{14}x_{23}$. Это следует, например, из
соотношения на миноры матрицы
$\begin{pmatrix}a_1&a_2&a_3&a_4\\b_1&b_2&b_3&b_4\end{pmatrix}$.
\end{example}

\subsection{Параболические подгруппы}
Оказывается, любое $X$, являющееся орбитой в $\mathbb P(V)$, задается
квадратичными уравнениями в проективных координатах.

Пусть $v\in V$, $X=G\cdot\la v\ra$~--- орбита вектора
$v$. Тогда $\Stab_G(\la v\ra)=P$~--- параболическая подгруппа
в $G$. Проективное однородное многообразие задается подгруппой $P$ с
точностью до сопряженности.

Посмотрим, как тор $T$ в $G$ действует на вектор $v$. Из равенства
$T\la v\ra=\la v\ra$ следует, что найдется $\lambda\colon T\to\mathbb
G_m$ (\term{вес} неприводимого представления $V$) такое, что
$tv=\lambda(t)v$ для всех $t\in T$. Представление задается своим старшим весом
(точнее, орбитой веса относительно $W$, но в этой орбите есть
единственный доминантный вес). %???
% $v$~--- вектор старшего веса.
Пусть $\alpha_1,\dots,\alpha_l$~--- простые корни.
Рассмотрим базис $\alpha_1^\vee,\dots,\alpha_l^\vee$
в двойственном пространстве, где $\alpha_i^\vee$ определяется
равенством
$\alpha_i^\vee(\beta) = \frac{2(\alpha_i,\beta)}{(\alpha_i,\alpha_i)}$.
Пусть $\varpi_1,\dots,\varpi_l$~--- двойственный к нему базис.
Таким образом,
$\frac{2(\alpha_i,\varpi)}{(\alpha_i,\alpha_i)} = \delta_{ij}$.
Эти элементы $\varpi_1,\dots,\varpi_l$ называются
\term{фундаментальными весами}.
Вес $\lambda$ раскладывается по этому базису следущим образом:
$\lambda = \sum m_i\omega_i$.
После этого $X$ (и $P$) зависит только от того, какие из $m_i$ не равны $0$.
То есть, проективные однородные многообразия задаются подмножеством
вершин на диаграмме Дынкина, состоящим из тех вершин, для которых
$m_i\neq 0$.
Мы будем их обводить на картинке.

Например, картинка для проективного пространства такая:
\[
\begin{tikzpicture}[scale=1.0,thin]
\def\sm{0.7}
\coordinate (1) at ($\sm*(0, 0)$);
\coordinate (2) at ($\sm*(1.4, 0)$);
\coordinate (3) at ($\sm*(2.8, 0)$);
\coordinate (3p) at ($\sm*(3.3, 0)$);
\coordinate (d1) at ($\sm*(3.5, 0)$);
\coordinate (d2) at ($\sm*(3.7, 0)$);
\coordinate (d3) at ($\sm*(3.9, 0)$);
\coordinate (4m) at ($\sm*(4.1, 0)$);
\coordinate (4) at ($\sm*(4.6, 0)$);
\coordinate (5) at ($\sm*(6.0, 0)$);
\node at (1) [below=3pt,font=\scriptsize] {$1$};
\node at (2) [below=3pt,font=\scriptsize] {$2$};
\node at (3) [below=3pt,font=\scriptsize] {$3$};
\node at (4) [below=3pt,font=\scriptsize] {$n-2$};
\node at (5) [below=3pt,font=\scriptsize] {$n-1$};
\draw (1)--(2);
\draw (2)--(3);
\draw (3)--(3p);
\draw (4m)--(4);
\draw (4)--(5);
\foreach \point in
{1,2,3,4,5}
  {
     \fill [white] (\point) circle (2.0pt);
     \draw [black] (\point) circle (2.0pt);
  }
\draw [black] (1) circle (5.0pt);
\foreach \point in
{d1,d2,d3}
  {
     \fill [black] (\point) circle (0.7pt);
  }
\node (c) at ($\sm*(7.5, 0)$) {$\mathsf{A}_{n-1}$};
\end{tikzpicture}
\]
Это соответствует векторному представлению $V = V(\varpi_1)$ группы
$\SL_n$.
Вообще, если на диаграмме $\mathsf{A}_{n-1}$ обвести вершину с номером $k$,
получится $\Gr(k,n)$.
Есть еще, например, присоединенное представление: $\SL_n$ действует
на своей алгебре Ли $\Lie(\SL_n)$. Картинка для этого представления
такая: 
\[
\begin{tikzpicture}[scale=1.0,thin]
\def\sm{0.7}
\coordinate (1) at ($\sm*(0, 0)$);
\coordinate (2) at ($\sm*(1.4, 0)$);
\coordinate (2p) at ($\sm*(1.9, 0)$);
\coordinate (d1) at ($\sm*(2.1, 0)$);
\coordinate (d2) at ($\sm*(2.3, 0)$);
\coordinate (d3) at ($\sm*(2.5, 0)$);
\coordinate (3m) at ($\sm*(2.7, 0)$);
\coordinate (3) at ($\sm*(3.2, 0)$);
\coordinate (4) at ($\sm*(4.6, 0)$);
\node at (1) [below=3pt,font=\scriptsize] {$1$};
\node at (2) [below=3pt,font=\scriptsize] {$2$};
\node at (3) [below=3pt,font=\scriptsize] {$n-2$};
\node at (4) [below=3pt,font=\scriptsize] {$n-1$};
\draw (1)--(2);
\draw (2)--(2p);
\draw (3m)--(3);
\draw (3)--(4);
\foreach \point in
{1,2,3,4}
  {
     \fill [white] (\point) circle (2.0pt);
     \draw [black] (\point) circle (2.0pt);
  }
\draw [black] (1) circle (5.0pt);
\draw [black] (4) circle (5.0pt);
\foreach \point in
{d1,d2,d3}
  {
     \fill [black] (\point) circle (0.7pt);
  }
\node (c) at ($\sm*(6.1, 0)$) {$\mathsf{A}_{n-1}$};
\end{tikzpicture}
\]
Первая вершина соответствует $V$, последняя~--- $V^*$,
в итоге получаем $V^*\otimes V\isom \End(V)$.

Если обведена одна вершина ($V = V(\varpi_i)$), то $P$ называется
\term{максимальной} параболической.
Если все вершины обведены, то $P$ называется \term{борелевской}
(это минимальная среди параболических).
Любая гладкая замкнутая подгруппа, содержащая $B$, называется
\term{параболической} и получается таким образом: $B\leq P \leq G$.

Пусть теперь на диаграмме Дынкина системы $\mathsf{A}_{n-1}$ обведены
вершины с номерами $k_1,\dots,k_m$.
Полученное многообразие можно описать в терминах стандартного
представления $V = K^n$ группы $\SL_n$.
А именно,
\[
X = \{U_1\leq\dots\leq U_m\mid \dim U_i = k_i\}.
\]
Такое $X$ называется \term{многообразием флагов}.
При этом $\SL_n$ действует на $X$ транзитивно.
Заметим, что тензорное произведение
$V(\varpi_{k_1})\otimes\dots\otimes V(\varpi_{k_m}0$ уже не обязано
быть неприводимым, но можно взять кусок, соответствующий
весу $\varpi_{k_1} + \dots + \varpi_{k_m}$.

Так мы описали все однородные проективные многообразия для группы
$\SL_n$.
В общем случае (для произвольной $G$) иногда однородное проективное
многообразие называют \term{обобщенным флаговым многообразием}.
Его можно описать так:
\[
X = \{P'\leq G\mid P'\mbox{ сопряжена с }P\},
\]
где значок $P'\leq G$ означает, что $P'$~--- гладкая замкнутая подгруппа
в $G$.
Более точно,
\[
X(R) = \{P'\leq G_R \mid\mbox{существуют }S/R, g\in G(S):\; gP'g^{-1} = P\}.
\]
После подкрутки на торсор $E$ получаем
\[
{}_{E}X = \{P'\leq {}_{E}G\mid P'_{\ol{K}}\mbox{ сопряжена с }P_{\ol{K}}
\mbox{ внутри }({}_{E}G)_{\ol{K}} = G_{\ol{K}}\}.
\]
Обратите внимание, что в ${}_{E}G$ никакой $P$ может не оказаться.

Проективное однородное многообразие $X$ \term{изотропно},
если $X(K)\neq\emptyset$.
Сама группа ${}_{E}G$ называется \term{изотропной}, если для какого-то
проективного однородного многообразия ${}_{E}X$, отличного от точки,
${}_{E}X$ изотропно.

\begin{example}
Пусть $G = \PGL_n$.
Ее скрученная форма ${}_{E}G$ имеет вид $\Aut(A)$, а соответствующая
скрученная форма проективного пространства~--- $\SB(A)$.
Заметим, что у $\PGL_n$ (в отличие от $\SL_n$) нет векторного представления.
Почему?
Для начала поймем, откуда берется скрученная форма $\SL_n$.
Отображение определителя $\det\colon\GL_n\to\mathbb{G}_m$
скручивается в \emph{приведенную норму} (\emph{reduced norm})
\[
\Nrd\colon A^* = \GL_1(A) \to \mathbb{G}_m.
\]
Ядро этого отображения обозначается через
$\SL_1(A) = \{g\in A\mid\Nrd(g)=1\}$.
Например, $(\Nrd(x))^n = \det(y\mapsto xy)$.

Решетка корней содержится в решетке весов:
\[
\mathbb{Z}\alpha_1\oplus\dots\oplus\mathbb{Z}\alpha_l
\leq
\mathbb{Z}\varpi\oplus\dots\oplus\mathbb{Z}\varpi_l
\]
Диаграммы Дынкина классифицируют расщепимые полупростые группы с точностью
до изогении, а класс изоморфности внутри класса изогении задается
промежуточной решеткой между этими двумя (с точностью до внешних
автоморфизмов).
Минимальная решетка $\mathbb{Z}\alpha_1\oplus\dots\oplus\mathbb{Z}\alpha_l$
соответствует присоединенной группе (без центра);
максимальная решетка $\mathbb{Z}\varpi_1\oplus\dots\oplus\mathbb{Z}\varpi_l$
соответствует односвязной группе (у нее самый большой центр).
\end{example}

% 12.03.2012

\subsection{$\SO_{2n}$}

Посмотрим на однородные многообразия для $\SO_{2n}$.
Диаграмма Дынкина выглядит так:
\[
\begin{tikzpicture}[scale=1.0,thin]
\def\sm{0.7}
\coordinate (1) at ($\sm*(0, 0)$);
\coordinate (2) at ($\sm*(1.4, 0)$);
\coordinate (2p) at ($\sm*(1.9, 0)$);
\coordinate (d1) at ($\sm*(2.1, 0)$);
\coordinate (d2) at ($\sm*(2.3, 0)$);
\coordinate (d3) at ($\sm*(2.5, 0)$);
\coordinate (3m) at ($\sm*(2.7, 0)$);
\coordinate (3) at ($\sm*(3.2, 0)$);
\coordinate (4) at ($\sm*(4.6, 0)$);
\coordinate (5) at ($\sm*(5.6, 1)$);
\coordinate (6) at ($\sm*(5.6, -1)$);
\node at (1) [below=5pt,font=\scriptsize] {$1$};
\node at (4) [below=5pt,font=\scriptsize] {$n-2$};
\node at (5) [below=5pt,font=\scriptsize] {$n-1$};
\node at (6) [below=5pt,font=\scriptsize] {$n$};
\draw (1)--(2);
\draw (2)--(2p);
\draw (3m)--(3);
\draw (3)--(4);
\draw (4)--(5);
\draw (4)--(6);
\foreach \point in
{1,2,3,4,5,6}
  {
     \fill [white] (\point) circle (2.0pt);
     \draw [black] (\point) circle (2.0pt);
  }
\foreach \point in
{d1,d2,d3}
  {
     \fill [black] (\point) circle (0.7pt);
  }
\node (c) at ($\sm*(7.1, 0)$) {$\mathsf{D}_{n}$};
\end{tikzpicture}
\]
Весу $\varpi_1$ отвечает квадрика $\{q(v)=0\}$, что соответствует естественному
представлению $V$ группы $\SO_{2n}$.
Весу $\varpi_2$~--- представление $\Lambda^2 V$.
Соответствующее многообразие~--- множество вполне изотропных плоскостей
$\la u,v\ra$, то есть, таких, что $q|_{\la u,v\ra}=0$.
Это условие можно описать так: $q(u) = q(v) = f(u,v) = 0$, где $f$~---
поляризация формы $q$: $f(u,v) = q(u+v) - q(u) - q(v)$.
Аналогично (с помощью вполне изотропных подпространств различной размерности)
описываются случаи $\varpi_3,\dots,\varpi_{n-2}$.

Весам $\varpi_{n-1}$ и $\varpi_n$ соответствуют вполне изотропные подпространства
размерности $n$.
Дело в том, что многообразие вполне изотропных подпространств размерности $n$
имеет две компоненты связности. Для того, чтобы объяснить этот эффект,
выберем базис $e_1,\dots,e_n,e_{-n},\dots,e_{-1}$, относительно
которого матрица Грама формы $q$ имеет вид
\[
\begin{pmatrix}
0 & 0 & \dots & 0 & 1\\
0 & 0 & \dots & 1 & 0\\
\vdots & \vdots & \ddots & \vdots & \vdots\\
0 & 1 & \dots & 0 & 0\\
1 & 0 & \dots & 0 & 0
\end{pmatrix}
\]
Оказывается, подпространства $\la e_1,\dots,e_{n-1},e_n\ra$
и $\la e_1,\dots,e_{n-1},e_{-n}\ra$ вполне изотропны, но не переводятся
друг в друга действием $\SO_{2n}$.
Первое соответствует весу $\varpi_{n-1}$, а второе~--- весу $\varpi_n$.
Куда же делось многообразие вполне изотропных подпространств размерности $n-1$?
Оно не максимальное однородное (соответствует не максимальной параболической
подгруппе), и соответствует весу $\varpi_{n-1} + \varpi_n$.
Действительно,
\[
\Stab(\la e_1,\dots,e_{n-1}\ra) = 
\Stab(\la e_1,\dots,e_{n-1},e_n\ra) \cap
\Stab(\la e_1,\dots,e_{n-1},e_{-n}\ra).
\]
Вообще, немаксимальные многообразия соответствуют флагам.
Посмотрим на вес $\varpi_{i_1} + \dots + \varpi_{i_k}$.
Флаг для него~--- это набор подпространств таких размерностей:
\[
\begin{tikzpicture}[scale=1.0,thin]
\def\sm{0.7}
\coordinate (1) at ($\sm*(0, 0)$);
\coordinate (2) at ($\sm*(1.4, 0)$);
\coordinate (2p) at ($\sm*(1.9, 0)$);
\coordinate (d1) at ($\sm*(2.1, 0)$);
\coordinate (d2) at ($\sm*(2.3, 0)$);
\coordinate (d3) at ($\sm*(2.5, 0)$);
\coordinate (3m) at ($\sm*(2.7, 0)$);
\coordinate (3) at ($\sm*(3.2, 0)$);
\coordinate (4) at ($\sm*(4.6, 0)$);
\coordinate (5) at ($\sm*(5.6, 1)$);
\coordinate (6) at ($\sm*(5.6, -1)$);
\node at (1) [below=5pt,font=\scriptsize] {$1$};
\node at (2) [below=5pt,font=\scriptsize] {$2$};
\node at (3) [below=5pt,font=\scriptsize] {$n-3$};
\node at (4) [below=5pt,font=\scriptsize] {$n-2$};
\node at (5) [below=5pt,font=\scriptsize] {$n$};
\node at (6) [below=5pt,font=\scriptsize] {$n$};
\draw (1)--(2);
\draw (2)--(2p);
\draw (3m)--(3);
\draw (3)--(4);
\draw (4)--(5);
\draw (4)--(6);
\foreach \point in
{1,2,3,4,5,6}
  {
     \fill [white] (\point) circle (2.0pt);
     \draw [black] (\point) circle (2.0pt);
  }
\foreach \point in
{d1,d2,d3}
  {
     \fill [black] (\point) circle (0.7pt);
  }
\node (c) at ($\sm*(7.1, 0)$) {$\mathsf{D}_{n}$};
\end{tikzpicture}
\]
с правильной инцидентностью.
А именно, для каждой из двух цепочек от первой вершины до двух последних
инцидентность~--- это включение, а для весов $\varpi_{n-1}$ и $\varpi_n$
инцидентность означает, что пересечение соответствующих подпространств
размерности $n$ имеет размерность $n-1$.

Перед нами пример \emph{геометрии}.
Гораздо более простой пример~--- случай системы $\mathsf{A}_2$.
Там всего два фундаментальных веса:
$\varpi_1$ соответствует точкам (и параболическим подгруппам типа $\varpi_1$),
а $\varpi_2$~--- прямым (и параболическим подгруппам типа $\varpi_2$).
Более подробно, посмотрим на трехмерное векторное пространство $F^3$.
Ненулевой вектор $u$ порождает одномерное подпространство
$\la u\ra\subseteq F^3$, и его стабилизатор
$\Stab_{\SL_3}(\la u\ra)$~--- это параболическая подгруппа типа $\varpi_1$:
\[
\begin{pmatrix}
* & * & * \\ 0 & * & * \\ 0 & * & *
\end{pmatrix}
\mbox{ --- стабилизатор вектора }
\begin{pmatrix} 1 \\ 0 \\ 0 \end{pmatrix}.
\]
Для описания прямых можно воспользоваться двойственностью и перейти
к пространству $(F^3)^*$.
Ненулевой ковектор $\ph\in(F^3)^*$ порождает одномерное подпространство
$\la\ph\ra\subseteq (F^3)^*$, и его стабилизатор
$\Stab_{\SL_3}(\la \ph \ra)$~--- это параболическая подгруппа типа
$\varpi_2$:
\[
\begin{pmatrix}
* & * & * \\ * & * & * \\ 0 & 0 & *
\end{pmatrix}
\mbox{ --- стабилизатор ковектора }
\begin{pmatrix} 0 & 0 & 1\end{pmatrix}.
\]
Отношение инцидентности между ними такое:
точка лежит на прямой тогда и только тогда, когда $\ph(u) = 0$.
В терминах параболических подгрупп:
$\Stab(\la u\la) \cap \Stab(\la \ph \ra)$ содержит борелевскую подгруппу
(то есть, параболическую подгруппу типа $\varpi_1 + \varpi_2$).

Если мы теперь посмотрим на геометрию, заданную абстрактными аксиомами
проективной плоскости (с аксиомой Дезарга, обеспечивающей ассоциативность,
но без аксиомы Паппа, обеспечивающей коммутативность),
мы получим группу $\SL_1(A)$, где $A$~--- центральная простая алгебра
степени $3$.

\subsection{Вычисление колец Чжоу}\label{ssect:chow-map-definition}

Пусть $E \in H^1 (F, G)$, и задано однородное проективное $G$-многообразие $X$.
Рассмотрим скрученное многообразие ${}_E X$; нас интересуют инварианты этого
многообразия в смысле алгебраической геометрии.
Например, $\CH^*({}_E X)$.

Вложение поля $F$ в его алгебраическое замыкание $\ol{F}$ дает морфизм
схем $\Spec\ol{F} \to \Spec F$.
Пулбэком получается многообразие $X_{\ol{F}}$:
\[
\begin{tikzcd}
X_{\ol{F}} \arrow{r} \arrow{d} & X \arrow{d} \\
\Spec\ol{F} \arrow{r} & \Spec F
\end{tikzcd}
\]
Отсюда получаем гомоморфизм
\[
\CH^*({}_{E}X) \to \CH^*(({}_{E}X)_{\ol{F}}) = \CH^*(X_{\ol{F}}).
\]
Нас интересует образ этого гомоморфизма: кручение содержится в его ядре,
за счет чего легче жить.
Первый шаг~--- вычисление $\CH^*(X_{\ol{F}})$.

\subsection{Пример: проективное пространство}\label{ssect:chow-ring-of-pn}

\begin{example}\label{example:projective-space}
Рассмотрим $\mathbb{P}^n \times \mathbb{P}^n$ с диагональным действием
$\SL_{n+1}$.
Это действие не транзитивно: есть диагональ $\mathbb{P}^n$.
Как выглядит дополнение к диагонали?
Мы утверждаем, что оно расслаивается над $\Gr(1,2;n+1)$ со слоем
$\mathbb{A}^1$.
Здесь $\Gr(1,2;n+1)$~--- многообразие флагов, состоящих из прямой и плоскости,
в $(n+1)$-мерном пространстве.
\[
\begin{tikzcd}
\mathbb{P}^n \times \mathbb{P}^n & \arrow{d}{\mathbb{A}^1} &
\mathbb{P}^n \arrow[left hook->]{ll} \\
 & \Gr(1,2;n+1)
\end{tikzcd}
\]
Это расслоение выглядит так: пара
$(\la u\ra, \la v\ra)\in\mathbb{P}^n \times \mathbb{P}^n$
отправляется во флаг $\la u\ra \leq \la u,v\ra$.
Прообраз флага при этом~--- это многообразие способов дополнить прямую
до плоскости, то есть, $\mathbb{P}^1 \setminus \mathbb{P}^0 = \mathbb{A}^1$.
Более строго, нужно говорить про расслоения на $\Gr(1,2;n+1)$:
есть двумерное векторное расслоение $\tau_2$, сопоставляющее
флагу $\la u\ra \leq \la u,v\ra$ плоскость $\la u,v\ra$,
и есть одномерное векторное расслоение $\tau_1$, сопоставляющее
флагу $\la u\ra \leq \la u,v\ra$ прямую $\la u\ra$.

Теперь зафиксируем в этом описании $u$, то есть, возьмем слой всей
картинки над точкой в первом сомножителе $\mathbb{P}^n$.
Получим картинку
\[
\begin{tikzcd}
\mathbb{P}^n & \arrow{d}{\mathbb{A}^1} &
\pt \arrow[left hook->]{ll} \\
 & \Gr(1;n)
\end{tikzcd}
\]
Заметим, что $\Gr(1,n) = \mathbb{P}^{n-1}$.
Поэтому можно написать точную последовательность локализации:
\[
\CH^{*-n}(\pt) \to \CH^*(\mathbb{P}^n) \to \CH^*(\mathbb{P}^{n-1}) \to 0.
\]
Средняя стрелка является гомоморфизмом колец, а первый член почти всегда
равен нулю.
Поэтому
\[
\CH^i(\mathbb{P}^n) = \begin{cases}
\CH^i(\mathbb{P}^{n-1}), & i < n,\\
\mathbb{Z}, & i = n,\\
0, & i > n.
\end{cases}
\]
По индукции получаем, что у $\CH^*(\mathbb{P}^n)$ в каждой размерности
от $0$ до $n$ стоит одна копия $\mathbb{Z}$.
\end{example}

\begin{example}\label{example:projective-space-filtration}
Опишем другой способ.
Пусть $\dim(V) = n+1$.
Рассмотрим действие группы $\SL(V)$ (или $\PGL(V)$)
на $\mathbb{P}(V^*) \times \mathbb{P}(V)$
(соответствующее весу $\varpi_1 + \varpi_n$).
Там имеется подмногообразие $\{\ph(u) = 0\}$:
\[
\begin{tikzcd}
\mathbb{P}(V^*) \times \mathbb{P}(V)
& \arrow{d}{\mathbb{A}^n}
& \{\ph(u) = 0\}\arrow[left hook->]{ll}\\
& \mathbb{P}(V^*)
\end{tikzcd}
\]
Зафиксировав $\ph$, получаем
\[
\begin{tikzcd}
\mathbb{P}^n
& \arrow{d}{\mathbb{A}^n}
& \mathbb{P}^{n-1} \arrow[left hook->]{ll}\\
& \pt
\end{tikzcd}
\]
Значит, имеется следующая точная последовательность локализации:
\[
\CH^{*-1}(\mathbb{P}^{n-1}) \to \CH^*(\mathbb{P}^n) \to \CH^*(\pt) \to 0.
\]
Вычисление по индукции приводит к тому же результату, что и
в предыдущем примере.
\end{example}

\begin{fact}
Если $Z\subseteq X$~--- замкнутое подмногообразие,
и $U = X\setminus Z$, имеется точная последовательность локализации
\[
\CH^{* - \codim_{X}Z} \to \CH^*(X) \to \CH^*(U) \to 0,
\]
где первое отображение~--- push-forward, а второе~--- pull-back
(и является гомоморфизмом колец).
\end{fact}

\begin{example}
Тот же результат можно получить и прямым вычислением:
понять, что компонента кольца Чжоу коразмерности $i$
порождается классом подпространства $[\mathbb{P}^{n-i}]$,
причем $[\mathbb{P}^n] = 1$.
Кроме этого,
\[
[\mathbb{P}^{n-1}]^i = \begin{cases}
[\mathbb{P}^{n-i}], & i \leq n,\\
0, & i > n.
\end{cases}
\]
Например, выбрав на $\mathbb{P}^n$ однородные координаты
$[x_0:\dots:x_n]$, можно взять $\mathbb{P}^{n-1} = \{x_0=0\}$,
другое $\mathbb{P}^{n-1} = \{x_1 = 0\}$ и обнаружить,
что их пересечение равно $\{x_0 = x_1 = 0\} = \mathbb{P}^{n-2}$.
\end{example}

\begin{remark}
По сути, в примере~\ref{example:projective-space-filtration}
мы нарисовали фильтрацию
\[
\begin{tikzcd}
\mathbb{P}^n
& \mathbb{P}^{n-1} \arrow[left hook->]{l}{\mathbb{A}^n}
& \dots \arrow[left hook->]{l}{\mathbb{A}^{n-1}}
& \pt \arrow[left hook->]{l}{\mathbb{A}^1}
\end{tikzcd}
\]
Вообще, если у многообразия $X$ существует фильтрация замкнутыми
подмногообразиями $S\supseteq X_1\supseteq X_2\supseteq\dots$
такая, что $X_i\setminus X_{i+1} = \coprod\mathbb{A}^{k_i}$,
то $X$ называется \term{клеточным}.
В этом случае
\begin{itemize}
\item все $\CH^i$~--- свободные конечно порожденные абелевы группы (их ранг
равен количеству клеток в соответствующей разности);
\item $\CH(X)_i\isom \CH(X_L)_i$ для любого расширения $L/F$.
\end{itemize}
\end{remark}

\subsection{Пример: многообразие Севери--Брауэра}

Перейдем теперь к $\SB(D)$, где $D$~--- тело, $\ind D = n+1$.
Это скрученная форма $\mathbb{P}^n$: $\SB(D) = {}_{E}\mathbb{P}^n$.
В разделе~\ref{ssect:chow-map-definition} мы построили отображение
\[
\CH^*(\SB(D)) \to \CH^*(\mathbb{P}^n_{\ol{F}}).
\]
Циклы из его образа называются \term{рациональными}
(по отношению к скручивающему торсору $E$).
В разделе~\ref{ssect:chow-ring-of-pn} мы вычислили правую часть:
там стоит копия $\mathbb{Z}$ в каждой компоненте с номерами от $0$ до $n$.
Образующая компоненты коразмерности $0$ всегда оказывается в образе.

Предположим, что класс $[\pt]$ оказался рационален.
Это означает, что есть конечные (сепарабельные) расширения
$L_1,\dots,L_k$ такие, что
\begin{itemize}
\item над каждым $L_i$ наше многообразие имеет рациональную точку;
\item $\gcd_i([L_i:F]) = 1$.
\end{itemize}
Заметим, что первое условие равносильно тому, что
$[D_{L_i}]=0$ в $\Br(L_i) = 0$.
Применим отображение трансфера $\Br(L_i) \to \Br(F)$.
Получим, что $[L_i:F]\cdot [D]=0$ в $\Br(F)$ для всех $i$.
Из этого (а также из второго условия)
следует, что $[D] = 0$ в $\Br(F)$.

\subsection{Пример: квадрика}\label{ssect:quadric}

Рассмотрим квадрику $Q = \{q=0\}$.
В $Q\times Q = \{(\la u\ra, \la v\ra)$ есть подмножество $\{f(u,v)=0\}$,
а в нем~--- диагональ $\{\la u\ra = \la v\ra\}\isom Q$.
Получаем фильтрацию
\[
\begin{tikzcd}
Q\times Q & \arrow{d}{\mathbb{A}^{\dim Q}} &
\{f(u,v) = 0\} \arrow[left hook->]{ll} & \arrow{d}{\mathbb{A}^1} &
Q.\arrow[left hook->]{ll}\\
 & Q & & \OGr(1,2;f)
\end{tikzcd}
\]
Здесь $\OGr(1,2;f)$ означает многообразие флагов, состоящих из
вполне изотропных подпространств вида $\la u\ra \leq \la u,v\ra$.

Расслоение $Q\times Q\setminus \{f(u,v)=0\} \to Q$
устроено так: пара $(\la u\ra, \la v\ra)$ отправляется
в $\la u\ra$.
Проверим, что слой изоморфен $\mathbb{A}^{\dim Q}$.
Пусть $u = e_1$.
Тогда наше дополнение имеет вид $\{f(e_1,v)\neq 0\}$.
Условие $f(e_1,v)\neq 0$ равносильно тому, что коэффициент у $v$
при базисном векторе $e_{-1}$ не равен $0$.
Поэтому можно читать, что он равен $1$.
Теперь все коэффициенты $v$, кроме тех, что стоят при $e_{1}$ и $e_{-1}$,
можно брать какими угодно, а коэффициент при $e_1$ определяется
однозначно из условия изотропности $q(v) = 0$.
Иначе говоря, если $\tau$~--- тавтологическое расслоение на $Q$,
рассмотрим $(\tau^{\perp})^*$.
Его слой над точкой $\la u\ra\in Q$ равен $(\la u\ra)^{\perp})^*$.
Вот нужный нам изоморфизм:
\begin{align*}
\mathbb{A}^{\dim Q} & \to
\mathbb{P}((\la u\ra^{\perp})^*) \setminus
\mathbb{P}(\{\ph\in(\la u\ra^{\perp})^*\mid \ph(u) = 0\},\\
v & \mapsto (\ph\colon w\mapsto f(v,w)).
\end{align*}

Расслоение $\{f(u,v)=0\} \setminus Q \to \OGr(1,2;f)$ устроено проще:
его слой равен
$\mathbb{P}(\tau_2) \setminus \mathbb{P}(\tau_1)\isom\mathbb{A}^1$,
как и в примере~\ref{example:projective-space}.

Теперь зафиксируем $u$; получим фильтрацию
\[
\begin{tikzcd}
Q & \arrow{d}{\mathbb{A}^{\dim Q}} &
\{f(u,v) = 0\} \arrow[left hook->]{ll} & \arrow{d}{\mathbb{A}^1} &
\pt,\arrow[left hook->]{ll}\\
 & \pt & & Q'
\end{tikzcd}
\]
где $Q'$~--- квадрика размерности $\dim Q - 2$.
Получаем точные последовательности
\begin{gather*}
\CH^{*-1}(\{f(u,v)=0\}) \to \CH^*(Q) \to \CH^*(\pt) \to 0,\\
\CH^{*-\dim Q + 1}(\pt) \to \CH^*(\{f(u,v)=0\}) \to \CH^*(Q') \to 0.
\end{gather*}
Теперь при помощи индукции можно доказать следующее.

Пусть $\dim Q = n$ четно.
Тогда $\CH^i(Q)$~--- свободная абелева группа ранга $1$
для всех $i=0,\dots,n$, кроме $i= n/2$; $\CH^{n/2}(Q)\isom\mathbb Z^2$.
Обозначим за $h = [Q'']\in\CH^1(Q)$ класс подквадрики коразмерности $1$.
Это гиперплоское сечение $Q$ в общем положении.
Тогда $1$~--- образующая $\CH^0(Q)$
$h$~--- образующая $\CH^1(Q)$,
$h^2$~--- образующая $\CH^2(Q)$,\dots.
С другой стороны, $\pt$~--- образующая $\CH^n(Q)$,
$[\mathbb{P}^1]$~--- образующая $\CH^{n-1}(Q)$,
$[\mathbb{P}^2]$~--- образующая $\CH^{n-2}(Q)$,\dots.
Это классы изотропных подпространств соответствующих размерностей.
Наконец, $h^{n/2}$ является суммой двух образующих; в качестве
одной из них можно взять $[\mathbb{P}^{n/2}$.

Это можно увидеть в координатной записи:
$Q$ задается уравнением $x_1 y_1 + \dots + x_{n/2+1}y_{n/2+1} = 0$.
После этого $Q''$ задается уравнением $x_{n/2+1} - y_{n/2+1} = 0$
(это гиперплоское сечение, как и было обещано),
а следующие образующие задаются последовательным
наложением уравнений $x_{1} = 0$,
$x_{2} = 0$, и так далее.
Когда дойдем до коразмерности $n/2$,
получим два варианта: либо
\[ x_1 = \dots = x_{n/2+1} = 0, \]
либо
\[ x_1 = \dots = x_{n/2} = y_{n/2+1} = 0. \]

\begin{example}
Пусть $n=4$, то есть, мы имеем дело с $\mathsf{D}_3$.
Перед нами четырехмерная квадрика.
Ее уравнение выглядит так: $x_1y_1 + x_2y_2 + x_3y_3 = 0$.
Уравнения двух образующих в коразмерности $4/2=2$ выглядят так:
\begin{gather*}
x_1 = x_2 = x_3 = 0,\\
x_1 = x_2 = y_3 = 0.
\end{gather*}
Их пересечение имеет вид $x_1 = x_2 = x_3y_3 = 0$,
что равносильно $x_1 = x_2 = 0$.
Почему-то это условие равносильно $x_3 = y_3 = 0$.
\end{example}

% 19.03.2012

\subsection{Нерасщепимая квадрика}

Что произойдет, если взять нерасщепимую квадрику?
Возьмем торсор $E\in H^1(F, O_{2n+2})$ и построим ${}_{E}Q$.
Как вычислить $\CH^*({}_{E}Q)$?
Более простой вопрос:
рассмотрим отображение
\[
\CH^*({}_{E}Q) \xrightarrow{\res} \CH^*(({}_{E}Q)_{\ol{F}}).
\]
Что можно сказать про образ этого отображения (то есть,
про рациональные циклы)?
Продолжим считать для простоты, что $n$ четно.
Мы знаем, что стоит в правой части: образующие
$1,h,h^2,\dots$ в коразмерностях $0,1,2,\dots,$ до середины,
образующие $[\pt],[\mathbb{P}^1],\mathbb{P}^2],\dots$
в коразмерностях $n,n-1,n-2,\dots$ до середины,
и две образующие $[\Pi_1],[\Pi_2]$ в коразмерности $n/2$.
При этом $h^{n/2} = [\Pi_1] + [\Pi_2]$.
Умножение выглядит так:
$h\cdot[\mathbb{P}^i] = [\mathbb{P}]^{i-1}$,
$h\cdot[\Pi_1] = h\cdot [\Pi_2] = [\mathbb{P}^{n/2}]$.

Во всяком случае, $h$ рационален: можно взять любую гладкую
подквадрику коразмерности $1$.
Пусть ${}_{E}Q$ задается уравнением $q=0$.
В случае расщепимой [четномерной] квадики это было уравнение
$x_1y_1 + \dots + x_{n/2+1}y_{n/2+1} = 0$,
и подквадрика выделялась дополнительным условием $x_{n/2+1} - y_{n/2+1}=0$.
В общем случае можно взять любой $v$ такой, что $q(v)\neq 0$,
и $q|_{\la v\ra^{\perp}}$ задает гладкую подквадрику коразмерности $1$.

\begin{theorem}[Springer]
Предположим, что $q$ \term{анизотропна}, то есть,
$q(v)\neq 0$ при $v\neq 0$.
Тогда класс $[\pt]$ не рационален.
\end{theorem}

Теорема доказывается так: класс $[\pt]$ рационален тогда и только тогда,
когда найдутся расширения $E_i/F$  такие, что
$\gcd([E_i:F]) = 1$, и над каждым $E_i$ квадрика
$q_{E_i}$ изотропна.
В частности, среди степеней расширений должна быть хотя бы одна нечетная,
и потому для некоторого $E/F$ с нечетным $[E:F]$ квадрика $q_E$
изотропна.
Но из этого следует, что $q$ изотропна (это, собственно, и есть
классическая теорема Спрингера).

Вот ответ на вопрос про образ: если $Q$ анизотропна,
то
\[
\im(\CH^k({}_{E}Q) \to \CH^k(({}_{E}Q)_{\ol{F}}))
= \begin{cases}
\CH^k(({}_{E}Q)_{\ol{F}}), & k < n/2, \\
2\cdot\CH^k(({}_{E}Q)_{\ol{F}}), & k > n/2, \\
2\cdot\mathbb{Z}[\Pi_1] + \mathbb{Z}([\Pi_1] - [\Pi_2]), & k = n/2
\end{cases} % Check the coefficients here!!
\]
Если $q$ \term{изотропна}, то есть существует ненулевой вектор $v$
такой, что $q(v) = 0$,
то можно выделить гиперболическую плоскость:
$q = \la 1,-1\ra \perp q'$.
Проитерируем этот процесс: получим
\[
q = k\cdot \la 1, -1\ra \perp q_{\an},
\]
где $q_{\an}$ и $k$ определены однозначно ($q_{\an}$~--- с точностью
до изометрии).
При этом $k$ называется \term{индексом Витта} формы $q$,
а $q_{\an}$~--- ее \term{анизотропной частью}.

Так вот, если индекс Витта нашей формы $q$ равен $k$,
то циклы $[\pt], [\mathbb{P}^1], \dots, [\mathbb{P}^{k-1}]$
рациональны.
Обратное тоже верно: если эти циклы рациональны, то индекс Витта
не меньше $k$.

Резюме: рациональные циклы на самой квадрике контролируют только
ее индекс Витта.

Посмотрим теперь на другое многообразие, связанное с торсором
$E\in H^1(F, O_{2n+2})$ (мы для удобства изменим нумерацию).
А именно, рассмотрим $\OGr(2,Q)$~--- многообразие вполне изотропных
плоскостей.
Мы реализовали $Q$ как $\{\la v\ra\mid q(v) = 0\}$.
Тогда $\OGr(2,Q) = \{\la u,v\ra\mid q(u) = q(v) = f(u,v) = 0\}$.

Чтобы добраться до этого многообразия,
положим $X = \{f(u,v)=0\}$
и рассмотрим фильтрацию из раздела~\ref{ssect:quadric}:
\[
\begin{tikzcd}
Q\times Q & \arrow{d}{\mathbb{A}^{n}} &
\{f(u,v) = 0\} \arrow[left hook->]{ll} & \arrow{d}{\mathbb{A}^1} &
\pt,\arrow[left hook->]{ll}\\
 & Q & & \OGr(1,2;Q).
\end{tikzcd}
\]
С одной стороны, $\OGr(1,2;Q)$~--- расслоение над $\OGr(2;Q)$ со слоем
$\mathbb{P}^1$.
С другой стороны, написанная фильтрация позволяет нам написать разложение
\[
\CH^*(Q\times Q) = \CH^{*-n}(Q)\oplus\CH^{*-n+1}(\OGr(1,2;Q))\oplus\CH^*(Q).
\]
Более того, морфизмы в левую части из слагаемых в правой части задаются
явным образом (с помощью пулбэков и пушфорвардов), и они
$O(q)$-эквивариантны.

\section{Мотивы Чжоу}

\subsection{Категория соответствий}\label{ssect:corr-category}

До сих пор мы смотрели на $\CH^*$ и на морфизмы вида
$\CH^*(X)\to\CH^*(\ol{X})$.
Посмотрим теперь на \term{мотив Чжоу} многообразия $X$.

Начнем с категории гладких проективных многообразий над $F$.
Что в ней плохо?
Например, то, что морфизмы нельзя складывать: она не аддитивна
(и тем более не абелева).
Каждому морфизму $f\colon X\to Y$ можно сопоставить
его график $\Gamma_f\subseteq X\times Y$
и получить $[\Gamma_f] \in \CH^*(X\times Y)$.
Элементы $\CH^*(X\times Y)$ уже можно складывать!
Поэтому в качестве промежуточного шага можно рассмотреть
\term{категорию соответствий} $\Cor_F$.
Ее объекты~--- гладкие проективные многообразия над $F$.
Морфизмы: $\Mor(X, Y) = \CH^{\dim Y}(X\times Y)$.
Роль тождественного морфизма играет класс диагонали.

\begin{remark}
В этой конструкции можно заменить $\CH$ на что-то другое,
где есть пулбэки и пушфорварды (они понадобятся нам ниже),
например, на другую теорию когомологий.
Если взять $K$-теорию~--- получим \emph{$K$-мотивы},
а не мотивы Чжоу).
\end{remark}

Как определить композицию таких морфизмов?
Пусть $\alpha\in\CH^{\dim Y}(X\times Y)$, $\beta\in\CH^{\dim Z}(Y\times Z)$.
Рассмотрим диаграмму
\[
\begin{tikzcd}
& X\times Y\times Z \arrow{dl}[swap]{\pr_{XY}} \arrow{dr}{\pr_{YZ}}
\arrow{dd}{\pr_{XZ}} \\
X\times Y & & Y\times Z\\
& X\times Z
\end{tikzcd}
\]
Из нее получается следующая диаграмма на уровне Чжоу:
\[
\begin{tikzcd}
& \CH^*(X\times Y\times Z)
\arrow{dd}{(\pr_{XZ})_*} \\
\CH^{\dim Y}(X\times Y)\arrow{ur}{\pr_{XY}^*}
& & \CH^{\dim Z}(Y\times Z)\arrow{ul}{\pr_{YZ}^*}\\
& \CH^{*}(X\times Z)
\end{tikzcd}
\]
Поэтому
$\pr_{XY}^*(\alpha)\cdot\pr_{YZ}^*(\beta)
\in\CH^{\dim Y + \dim Z}(X\times Y\times Z))$, и
мы можем определить
\[
\beta\circ\alpha = (\pr_{XZ})_*(\pr_{XY}^*(\alpha)\cdot\pr_{YZ}^*(\beta))
\]
Это произведение имеет характер свертки.
Например, можно взять в качестве $X,Y,Z$ метрические пространства,
а в качестве морфизмов~--- ядерные операторы,
и получится свертка.
Или в качестве $X,Y,Z$~--- конечные множества, а в качестве морфизмов~---
матрицы, и тогда получится произведение матриц.

\subsection{Карубизация}

Итак, в категории соответствий $\Cor_F$ морфизмы уже можно складывать:
это аддитивная категория.
Заметим, что в ней есть и прямые сумммы
($X\oplus Y = X\coprod Y$), и произведения ($X\otimes Y = X\times Y$).
Но эта категория не абелева (и даже не псевдоабелева).
Напомним, что категория называется \term{псевдоабелевой}, если
у любого проектора есть образ (в категорном смысле).
То есть, если $p\colon X\to X$~--- морфизм, для которого $p^2=p$,
то $X = X_1\oplus X_2$, причем $p$~--- проекция на $X_1$.

Есть стандартная процедура, как из аддитивной категории получить
псевдоабелеву: \term{пополнение по Каруби} (\term{карубизация}).
Таким образом по $\Cor_F$ строится
\term{категория мотивов Гротендика--Чжоу} $\mathcal{M}$.
Ее объекты~--- пары $(X,p)$, где $p\colon X\to X$~--- идемпотент.
Неформально говоря, эта пара символизирует <<образ>> морфизма $p$
(которого может не быть в исходной категории).
Морфизмы определяются так:
\[
\Mor((X,p),(Y,q)) = q\circ\Mor(X,Y)\circ p.
\]
Есть функтор $\Cor\to\mathcal{M}$, $X\mapsto (X,\id_X)$.
Для многообразия $X$ объект $M(X) = (X,\id_X)$ называется
\term{мотивом $X$}.
На самом деле, нужно писать $\Cor_{\operatorname{rat},\operatorname{eff}}$
вместо $\Cor$, и $\operatorname{Chow}^{\operatorname{eff}}$ вместо
$\mathcal{M}$.
Мы получили функторы
\[
\begin{tikzcd}
\operatorname{SmProj}/F \arrow{r} \arrow[bend right=15, swap]{rr}{M}
& \Cor_{\operatorname{rat},\operatorname{eff}}(F) \arrow{r}
& \operatorname{Chow}^{\operatorname{eff}}(F),
\end{tikzcd}
\]
где $M$~--- функтор <<взятия мотива>>.
При этом $M(X\coprod Y) = M(X) \oplus M(Y)$,
$M(X\times Y) = M(X)\otimes M(Y)$.

\subsection{Мотив проективной прямой}

Попробуем <<посчитать>> мотив проективной прямой $M(\mathbb{P}^1)$.
Напомним, что у $\CH^(\mathbb{P}^1)$ стоит $\mathbb{Z}$ в коразмерностях
$0$ (с образующей $1$) и $1$ (с образующей $[\pt]$).
Рассмотрим вложение $i\colon \pt\to \mathbb{P}^1$
и проекцию $\pi\colon\mathbb{P}^1\to\pt$.
Композиция $\pi\circ i \colon \pt \to \mathbb{P}^1 \to \pt$
тождественна, поэтому $p = i\circ\pi$ является проектором на $\mathbb{P}^1$.
Это идемпотент, отправляющий все в точку.
Поэтому в категории мотивов
$M(\mathbb{P}^1) = M(\pt) \oplus (\mathbb{P}^1, 1 - [p])$.
Слагаемое $(\mathbb{P}^1, 1 - [p])$ обозначается через $\mathbb{L}$
и называется \term{мотивом Лефшеца}.
Это аналог аффинной прямой в категории мотивов.
Оказывается, мотив Лефшеца неразложим.

Мотив точки часто обозначается через $\mathbb Z = M(\pt)$;
он играет роль нейтрального объекта относительно $\otimes$.
При этом мотив Лефшеца $\mathbb{L}$ обозначается
через $\mathbb{Z}(1)[2] = \mathbb{Z}\{1\}$.
Тензорные степени мотива Лефшеца обозначаются так:
$L^{\otimes k} = \mathbb{Z}(k)[2k] = \mathbb{Z}\{k\}$.
Это в некотором смысле <<мотив>> $k$-мерного аффинного пространства.

\subsection{Представимость функтора Чжоу}

Часто удается разложить мотив многообразия $X$ в прямую сумму вида
$M(X) = \bigoplus M(Y_i)\otimes\mathbb{L}^{\otimes k_i}$,
где $Y_i$~--- какие-то другие многообразия.
Поэтому удобно обозначение
$M(Y)\otimes\mathbb{L}^{\otimes k} = M(Y)(k)[2k] = M(Y)\{k\}$.

Что дает такого рода разложение?
\begin{fact}
Пусть
$M(X) = \bigoplus M(Y_i)\otimes\mathbb{L}^{\otimes k_i}$.
Тогда
$\CH^n(X) = \bigoplus\CH^{n-k_i}(Y_i)$.
\end{fact}
Например, из разложения $M(\mathbb{P}^1) = M(\pt)\otimes M(\pt)\{1\}$
следует, что
\begin{align*}
\CH^0(\mathbb{P}^1) &= \CH^0(\pt) = \mathbb{Z},\\
\CH^1(\mathbb{P}^1) &= \CH^1(\pt)\oplus\CH^0(\pt) = \mathbb{Z}.
\end{align*}
Вообще, $\CH^n(X) = \Mor(X,\mathbb{L}^{\otimes n})$
и $\CH_n(X) = \Mor(\mathbb{L}^{\otimes n}, X)$,
то есть, $\CH$~--- представимый функтор в категории мотивов,
а $\mathbb{L}$ играет роль пространства Эйленберга--Маклейна.

Умножение в $\CH^*$ тоже происходит из категории мотивов.
Пусть $\alpha\in\CH^k(X)$, $\beta\in\CH^n(X)$, то есть
$\alpha\colon M(X) \to \mathbb{L}^{\otimes k}$,
$\beta\colon M(X) \to \mathbb{L}^{\otimes n}$.
Перемножая эти отображения, получаем
\[
\alpha\otimes\beta \colon M(X)\otimes M(X) \to
\mathbb{L}^{\otimes k} \otimes \mathbb{L}^{\otimes n}.
\]
Правая часть изоморфна $\mathbb{L}^{\otimes(k+n)}$.
Взяв композицию с морфизмом $M(\Delta)\colon M(X) \to M(X\times X)$,
получаем $\alpha\cup\beta\colon M(X) \to \mathbb{L}^{\otimes(k+n)}$.

\subsection{Теорема Карпенко}
\begin{theorem}[Карпенко, 2000]\label{thm:karpenko}
Пусть дана фильтрация многообразия $X$ замкнутыми (не обязательно гладкими)
подмножествами
\[
\begin{tikzcd}[column sep=1.2em]
X = X_0 & \arrow{d}{\mathbb{A}^{k_0}} &
X_1 \arrow[left hook->]{ll} & \arrow{d}{\mathbb{A}^{k_1}} &
X_2 \arrow[left hook->]{ll} & & \dots \arrow[left hook->]{ll} & &
X_n \arrow[left hook->]{ll} & \arrow{d}{\mathbb{A}^{k_n}} &
X_{n+1} = \emptyset. \arrow[left hook->]{ll} \\
& Y_0 & & Y_1 & & & \dots & & & Y_n
\end{tikzcd}
\]
Вертикальные стрелки означают, что для каждого $i=0,\dots,n$
задан плоский морфизм $X_i\setminus X_{i+1} \to Y_i$,
слои которого~--- аффинные пространства $\mathbb{A}^{k_i}$.

Тогда $M(X) = \bigoplus M(Y_i)\{k_i\}$
и, кроме того,
$M(X) = \bigoplus M(Y_i)\{\dim X - \dim Y_i - k_i\}$.
В частности, имеется функториальный (по $Z$)
изоморфизм $\CH^*(X\times Z) \isom \bigoplus\CH^{*-k_i}(Y_i\times Z)$.
\end{theorem}

Фильтрация из теоремы~\ref{thm:karpenko} называется
\term{относительным клеточным разложением}.

\begin{example}
Фильтрация
\[
\begin{tikzcd}
\mathbb{P}^1 & \arrow{d}{\mathbb{A}^1} & \pt\arrow[left hook->]{ll}\\
& \pt
\end{tikzcd}
\]
приводит к разложению $M(\pt) \oplus M(\pt)\{1\}$.
\end{example}

\begin{example}
Фильтрация
\[
\begin{tikzcd}
Q\times Q & \arrow{d}{\mathbb{A}^{\dim Q}} &
\{f(u,v) = 0\} \arrow[left hook->]{ll} & \arrow{d}{\mathbb{A}^1} &
Q\arrow[left hook->]{ll}\\
 & Q & & \OGr(1,2;f)
\end{tikzcd}
\]
из раздела~\ref{ssect:quadric}
приводит к разложению
\[
M(Q\times Q) = M(Q) \oplus M(\OGr(1,2;Q))\{1\}\oplus M(Q)\{\dim Q\}.
\]
\end{example}

\begin{example}
Пусть на квадрике $Q$ есть рациональная точка (то есть, форма $q$
изотропна).
Тогда $q = \la 1,-1\ra \perp q'$, и есть фильтрация
\[
\begin{tikzcd}
Q & \arrow{d}{\mathbb{A}^{\dim Q}} &
X' \arrow[left hook->]{ll} & \arrow{d}{\mathbb{A}^1} &
\pt,\arrow[left hook->]{ll}\\
 & \pt & & Q'
\end{tikzcd}
\]
где $Q' = \{q'=0\}$.
Получаем разложение
\[
M(Q) = M(\pt) \oplus M(Q')\{1\} \oplus M(\pt)\{\dim Q\}.
\]
На картинке это выглядит так:
\[
\begin{tikzpicture}[scale=1.0,thin]
\def\sm{0.7}
\coordinate (1) at ($\sm*(0, 0)$);
\coordinate (2) at ($\sm*(1.4, 0)$);
\coordinate (2p) at ($\sm*(1.9, 0)$);
\coordinate (d1) at ($\sm*(2.1, 0)$);
\coordinate (d2) at ($\sm*(2.3, 0)$);
\coordinate (d3) at ($\sm*(2.5, 0)$);
\coordinate (3m) at ($\sm*(2.7, 0)$);
\coordinate (3) at ($\sm*(3.2, 0)$);
\coordinate (4) at ($\sm*(4.2, 1)$);
\coordinate (5) at ($\sm*(4.2, -1)$);
\coordinate (6) at ($\sm*(5.2, 0)$);
\coordinate (6p) at ($\sm*(5.7, 0)$);
\coordinate (e1) at ($\sm*(5.9, 0)$);
\coordinate (e2) at ($\sm*(6.1, 0)$);
\coordinate (e3) at ($\sm*(6.3, 0)$);
\coordinate (7m) at ($\sm*(6.5, 0)$);
\coordinate (7) at ($\sm*(7.0, 0)$);
\coordinate (8) at ($\sm*(8.4, 0)$);
\draw (1)--(2);
\draw (2)--(2p);
\draw (3m)--(3);
\draw (3)--(4);
\draw (3)--(5);
\draw (4)--(6);
\draw (5)--(6);
\draw (6)--(6p);
\draw (7m)--(7);
\draw (7)--(8);
\draw[dotted] ($\sm*(1, 1.3)$)--($\sm*(7.4, 1.3)$)--($\sm*(7.4,-1.3)$)
--($\sm*(1, -1.3)$)--($\sm*(1, 1.3)$);
\foreach \point in
{1,2,3,4,5,6,7,8}
  {
     \fill [white] (\point) circle (2.0pt);
     \draw [black] (\point) circle (2.0pt);
  }
\draw [black] (1) circle (5.0pt);
\draw [black] (8) circle (5.0pt);
\foreach \point in
{d1,d2,d3,e1,e2,e3}
  {
     \fill [black] (\point) circle (0.7pt);
  }
\node at (1) [below=3pt,font=\scriptsize] {$1$};
\node at (8) [below=3pt,font=\scriptsize] {$[\pt]$};
\draw [|->] ($\sm*(0, -1)$) --node[below=3pt, font=\scriptsize] {$1$}
($\sm*(1.4, -1)$);
\draw [|->] ($\sm*(0, -2)$) --node[below=3pt, font=\scriptsize] {$\dim Q$}
($\sm*(8.4, -2)$);
\end{tikzpicture}
\]
Обратите внимание, что на картинке выделен мотив подквадрики $Q'$,
который сдвигается на $1$.
Кроме того, мотив точки (справа) сдвигается на $\dim Q$.
Иными словами, у нас появились проекторы
$1\times [\pt]$, $[\pt]\times 1$, $\Delta_Q - 1\times[\pt] - [\pt]\times 1$.
\end{example}

\begin{remark}
Обозначение $\mathbb{Z}(1)[2]$ для мотива Лефшеца может показаться странным.
Здесь второй сдвиг соответствует сдвигу в триангулированной категории
Воеводского.
При желании можно представлять это как композицию двух сдвигов:
$(1)[1]$~--- сдвиг на $\mathbb{G}_m$, $(0)[1]$~--- сдвиг на $S^1$.
\end{remark}

\begin{remark}
В общем случае, разложение Брюа показывает, что если $G$~--- расщепимая
группа, $P$~--- ее параболическая подгруппа,
то мотив однородного многообразия $G/P$ равен прямой сумме
сдвигов $\mathbb{Z}$:
$M(G/P) = \bigoplus\mathbb{Z}\{\dots\}$.
При этом $\mathbb{Z}\{i\}$ встречается столько раз, каково количество
минимальных представителей классов смежности из $W/W_P$ длины $i$.
Поэтому сдвиги считываются из диаграммы Хассе.
\end{remark}
\begin{example}
Например, для $\mathbb{P}^n$ диаграмма Хассе выглядит так:
\[
\begin{tikzpicture}[scale=1.0,thin]
\def\sm{0.7}
\coordinate (1) at ($\sm*(0, 0)$);
\coordinate (2) at ($\sm*(1.4, 0)$);
\coordinate (2p) at ($\sm*(1.9, 0)$);
\coordinate (d1) at ($\sm*(2.1, 0)$);
\coordinate (d2) at ($\sm*(2.3, 0)$);
\coordinate (d3) at ($\sm*(2.5, 0)$);
\coordinate (3m) at ($\sm*(2.7, 0)$);
\coordinate (3) at ($\sm*(3.2, 0)$);
\coordinate (4) at ($\sm*(4.6, 0)$);
\node at (1) [below=3pt,font=\scriptsize] {$0$};
\node at (2) [below=3pt,font=\scriptsize] {$1$};
\node at (3) [below=3pt,font=\scriptsize] {$n-1$};
\node at (4) [below=3pt,font=\scriptsize] {$n$};
\draw (1)--(2);
\draw (2)--(2p);
\draw (3m)--(3);
\draw (3)--(4);
\foreach \point in
{1,2,3,4}
  {
     \fill [white] (\point) circle (2.0pt);
     \draw [black] (\point) circle (2.0pt);
  }
\foreach \point in
{d1,d2,d3}
  {
     \fill [black] (\point) circle (0.7pt);
  }
\end{tikzpicture}
\]
Поэтому мотив $\mathbb{P}^n$ равен
$\mathbb{Z}\oplus\mathbb{Z}\{1\}\oplus\dots\oplus\mathbb{Z}\{n\}$.
Для [расщепимой] квадрики $Q$ четной размерности диаграмма такая:
\[
\begin{tikzpicture}[scale=1.0,thin]
\def\sm{0.7}
\coordinate (1) at ($\sm*(0, 0)$);
\coordinate (2) at ($\sm*(1.4, 0)$);
\coordinate (2p) at ($\sm*(1.9, 0)$);
\coordinate (d1) at ($\sm*(2.1, 0)$);
\coordinate (d2) at ($\sm*(2.3, 0)$);
\coordinate (d3) at ($\sm*(2.5, 0)$);
\coordinate (3m) at ($\sm*(2.7, 0)$);
\coordinate (3) at ($\sm*(3.2, 0)$);
\coordinate (4) at ($\sm*(4.2, 1)$);
\coordinate (5) at ($\sm*(4.2, -1)$);
\coordinate (6) at ($\sm*(5.2, 0)$);
\coordinate (6p) at ($\sm*(5.7, 0)$);
\coordinate (e1) at ($\sm*(5.9, 0)$);
\coordinate (e2) at ($\sm*(6.1, 0)$);
\coordinate (e3) at ($\sm*(6.3, 0)$);
\coordinate (7m) at ($\sm*(6.5, 0)$);
\coordinate (7) at ($\sm*(7.0, 0)$);
\coordinate (8) at ($\sm*(8.4, 0)$);
\node at (1) [below=3pt,font=\scriptsize] {$0$};
\node at (2) [below=3pt,font=\scriptsize] {$1$};
\node at (5) [below=3pt,font=\scriptsize] {$n/2$};
\node at (7) [below=3pt,font=\scriptsize] {$n-1$};
\node at (8) [below=3pt,font=\scriptsize] {$n$};
\draw (1)--(2);
\draw (2)--(2p);
\draw (3m)--(3);
\draw (3)--(4);
\draw (3)--(5);
\draw (4)--(6);
\draw (5)--(6);
\draw (6)--(6p);
\draw (7m)--(7);
\draw (7)--(8);
\foreach \point in
{1,2,3,4,5,6,7,8}
  {
     \fill [white] (\point) circle (2.0pt);
     \draw [black] (\point) circle (2.0pt);
  }
\foreach \point in
{d1,d2,d3,e1,e2,e3}
  {
     \fill [black] (\point) circle (0.7pt);
  }
\end{tikzpicture}
\]
Поэтому $M(Q) = \mathbb{Z}\oplus\mathbb{Z}\{1\}\oplus\dots
\oplus\mathbb{Z}\{n/2\}^{\otimes 2} \oplus\mathbb{Z}\{n/2+1\}\oplus\dots
\oplus\mathbb{Z}\{n\}$.
\end{example}

\subsection{Метод общей точки}

Пока что мы получали мотивные разложения только с помощью фильтраций
и теоремы Карпенко.
Сейчас мы узнаем еще один прием~--- \emph{метод общей точки}.
Пусть $X,Y$~--- многообразия, причем $Y$ неприводимо.
Рассмотрим $X_{F(Y)}$ и его кольцо Чжоу $\CH^*(X_{F(Y)}$
(напомним, что $F(Y)$~--- поле рациональных функций на $Y$).
\begin{lemma}
Отображение $\CH^*(X\times Y)\to\CH^*(X_{F(Y)})$,
полученное из декартова квадрата
\[
\begin{tikzcd}
X_{F(Y)} \arrow{r}\arrow{d} & X\times Y \arrow{d} \\
\Spec F(Y) \arrow[right hook->]{r} & Y,
\end{tikzcd}
\]
сюръективно.
\end{lemma}
\begin{proof}
Указанную диаграмму можно представлять себе как индуктивный предел
диаграмм вида
\[
\begin{tikzcd}
X\times U \arrow{r}\arrow{d} & X\times Y \arrow{d} \\
U \arrow[right hook->]{r} & Y,
\end{tikzcd}
\]
где $U$~--- открытое непустое в $Y$ (поскольку $\Spec F(Y) = \varinjlim U$).
Но каждое полученное таким образом отображение
$\CH^*(X\times Y) \to \CH^*(X\times U)$ сюръективно в силу точной
последовательности локализации.
\end{proof}

Как этим пользоваться?
Чтобы выделить прямое слагаемое в мотиве $X$, нам нужно
найти проектор $p\in\CH^{\dim X}(X\times X)$.
Для этого есть два варианта:
\begin{enumerate}
\item взять $Y:= X$, выбрать какой-то элемент из $\CH^i(X_{F(X)})$
и поднять его в $\CH^i(X\times X)$;
\item взять какой-нибудь $Y$, построить элементы из $\CH^i(X\times Y)$,
$\CH^i(Y\times X)$, взять их композицию, и дальше как в первом пункте.
\end{enumerate}
Пусть $X$~--- гладкое проективное над $F$.
Напомним, что цикл $\alpha\in\CH^*(X_{\ol{F}})$ называется
рациональным, если он лежит в образе отображения
\[
\res\colon\CH^*(X) \to \CH^*(X_{\ol{F}}).
\]
Аналогично, можно рассмотреть отображение
\[
\res\colon\CH^*(X\times X) \to \CH^*(X_{\ol{F}}\times X_{\ol{F}}).
\]
Правая часть гораздо лучше левой.

Предположим, что мы нашли цикл $p\in\CH^*(X_{\ol{F}}\times X_{\ol{F}})$
такой, что
\begin{enumerate}
\item $p$~--- проектор (в смысле композиции, определенной в
разделе~\ref{ssect:corr-category});
\item $p$ рационален, то есть, $p$ поднимается до какого-то
$\wt{p}\in\CH^*(X\times X)$.
\end{enumerate}
Следует ли из этого, что $\wt{p}$ является проектором?
Вообще говоря~--- нет, но для однородных многообразий есть такая теорема.

\begin{theorem}[Теорема нильпотентности Роста]
Если $X$~--- проективное однородное многообразие,
$p$~--- рациональный проектор на $X_{\ol{F}}$, то он
поднимается до проектора $\wt{p}$ на $X$.
Более сильное утверждение:
\[
\ker(\CH^*(X\times X) \to \CH^*(X_{\ol{F}}\times X_{\ol{F}}))
\]
состоит из нильпотентных (в смысле композиции) элементов.
\end{theorem}

Как выглядят очевидные элементы $\CH^*(X\times X)$?
Можно взять $\alpha\in\CH^*(X)$, $b\in\CH^*(X)$, и получить
$a\times b\in\CH^*(X\times X)$.

\begin{exercise}
В этом случае $(a\times b)\circ (c\times d) = \deg(ad)(c\times b)$,
где $\deg\colon\CH^*(Y) \to \CH^*(\pt)$ происходит из
морфизма $Y\to\pt$.
\end{exercise}

\begin{exercise}
Пусть $R$~--- коммутативное кольцо, $I\trleq R$~--- идеал, состоящий
из нильпотентных элементов.
Тогда любой идемпотентв $R/I$ поднимается до идемпотента в $R$.
\end{exercise}

\begin{definition}
Многообразие $X$ называется \term{клеточным}, если существует
фильтрация вида
\[
\begin{tikzcd}[column sep=1.2em]
X = X_0 & \arrow{d}{\mathbb{A}^{k_0}} &
X_1 \arrow[left hook->]{ll} & \arrow{d}{\mathbb{A}^{k_1}} &
X_2 \arrow[left hook->]{ll} & & \dots \arrow[left hook->]{ll} & &
X_n \arrow[left hook->]{ll} & \arrow{d}{\mathbb{A}^{k_n}} &
X_{n+1} = \emptyset, \arrow[left hook->]{ll} \\
& \pt & & \pt & & & \dots & & & pt
\end{tikzcd}
\]
в которой все базы~--- точки.
\end{definition}
\begin{example}
Разложение Брюа говорит, что если группа $G$ расщепима,
$P$~--- параболическая подгруппа в $G$,
то многообразия $G/P$ клеточное.
\end{example}
Из теоремы Карпенко~\ref{thm:karpenko} следует, что для клеточного
многообразия $M(X) = \bigoplus \mathbb{Z}\{r_i\}$.

\begin{definition}
Многообразие $X$ называется \term{клеточным над общей точкой}
(\term{generically cellular}), если $X_{F(X)}$ клеточное.
\end{definition}
\begin{example}
Пусть $Q$~--- \term{Квадрика Пфистера}, то есть, $Q = \{q=0\}$,
где $q = \lAngle a_1,\dots,a_k\rAngle = \la 1,-a_1\ra \otimes\dots
\otimes\la 1,-a_k\ra$~--- $k$-кратная форма Пфистера
(размерность $Q$ равна $2^k-2$).
Тогда $Q$ клеточная над общей точкой.
\end{example}
Верно и обратное: все анизотропные четномерные квадрики,
клеточные над общей точкой, так выглядят.

\subsection{Мотив квадрики Пфистера}

Пусть $Q$~--- квадрика Пфистера размерности $2^k - 2$.
Мы знаем, что $Q_{F(Q)}$~--- клеточное многообразие.
Обозначим образующие компонент $\CH^*(Q_{F(Q)})$:
\[
\begin{tikzpicture}[scale=1.0,thin]
\def\sm{0.7}
\coordinate (1) at ($\sm*(0, 0)$);
\coordinate (2) at ($\sm*(1.4, 0)$);
\coordinate (2p) at ($\sm*(1.9, 0)$);
\coordinate (d1) at ($\sm*(2.1, 0)$);
\coordinate (d2) at ($\sm*(2.3, 0)$);
\coordinate (d3) at ($\sm*(2.5, 0)$);
\coordinate (3m) at ($\sm*(2.7, 0)$);
\coordinate (3) at ($\sm*(3.2, 0)$);
\coordinate (4) at ($\sm*(4.2, 1)$);
\coordinate (5) at ($\sm*(4.2, -1)$);
\coordinate (6) at ($\sm*(5.2, 0)$);
\coordinate (6p) at ($\sm*(5.7, 0)$);
\coordinate (e1) at ($\sm*(5.9, 0)$);
\coordinate (e2) at ($\sm*(6.1, 0)$);
\coordinate (e3) at ($\sm*(6.3, 0)$);
\coordinate (7m) at ($\sm*(6.5, 0)$);
\coordinate (7) at ($\sm*(7.0, 0)$);
\coordinate (8) at ($\sm*(8.4, 0)$);
\node at (1) [below=3pt,font=\scriptsize] {$0$};
\node at (2) [below=3pt,font=\scriptsize] {$1$};
\node at (5) [below=3pt,font=\scriptsize] {$2^{k-1}-1$};
\node at (7) [below=3pt,font=\scriptsize] {$2^k-3$};
\node at (8) [below=3pt,font=\scriptsize] {$2^k-2$};
\node at (1) [above=3pt,font=\scriptsize] {$1$};
\node at (2) [above=3pt,font=\scriptsize] {$h$};
\node at (4) [above=3pt,font=\scriptsize] {$\rho$};
\node at (6) [above=3pt,font=\scriptsize] {$h\rho$};
\node at (8) [above=3pt,font=\scriptsize] {$h^{2^{k-1}}\rho$};
\draw (1)--(2);
\draw (2)--(2p);
\draw (3m)--(3);
\draw (3)--(4);
\draw (3)--(5);
\draw (4)--(6);
\draw (5)--(6);
\draw (6)--(6p);
\draw (7m)--(7);
\draw (7)--(8);
\foreach \point in
{1,2,3,4,5,6,7,8}
  {
     \fill [white] (\point) circle (2.0pt);
     \draw [black] (\point) circle (2.0pt);
  }
\foreach \point in
{d1,d2,d3,e1,e2,e3}
  {
     \fill [black] (\point) circle (0.7pt);
  }
\end{tikzpicture}
\]
Возьмем $\rho\in\CH^*(Q_{F(Q)})$ и поднимем его до какого-то элемента
$\alpha\in\CH^*(Q\times Q)$.
Рассмотрим коммутативную диаграмму
\[
\begin{tikzcd}
\alpha\in\CH^*(Q\times Q) \arrow[->>]{r} \arrow{d}{\res}
& \CH^*(Q_{F(Q)})\ni\rho \arrow{d}{\isom} \\
\ol{\alpha}\in\CH^*(Q_{\ol{F}}\times Q_{\ol{F}}) \arrow[->>]{r}
& \CH^*(Q_{\ol{F}(Q)})\ni\ol{\rho}
\end{tikzcd}
\]
Мы не умеем следить за $\alpha$, но знаем, что $\ol{\alpha} = \res(\alpha)$
переходит в $\ol{\rho}$ (который отождествляется с $\rho$
при помощи изоморфизма), и знаем, как выглядит нижняя
горизонтальная стрелка.

Итак, $\ol{\alpha}$ является прообразом $\rho$, поэтому обязан иметь вид
\[
\ol{\alpha} = \ol{\rho}\times 1
+ c_1\cdot h^{2^{k-1}-2} \times h
+ c_2\cdot h^{2^{k-1}-3} \times h^2
+ \dots
+ c_{2^{k-1}-1}\cdot 1\times h^{2^{k-1}-1}
+ c\cdot 1\times\ol{\rho}.
\]
Заметим, что все слагаемые в правой части, кроме первого и последнего,
содержатся в образе $\res$, посольку $h$ рационален.
Поэтому (подправив $\alpha$) можно считать,
что $\ol\alpha = \ol\rho\times 1 + c\cdot 1 \times \ol\rho$.

Кроме того, цикл $2\ol\rho$ рационален, поскольку квадрика $Q$
ращепляется квадратичным расширением.
Действительно, если $q = \lAngle a_1,\dots,a_k\rAngle$, достаточно взять
поле $F(\sqrt{a_1})$.
Над этим полем $q$ изотропна, а потому гиперболична.
Рассуждение заканчивается рассмотрением диаграммы
\[
\begin{tikzcd}
\CH^*(Q) \arrow[yshift=-2pt]{r} \arrow{d} & \CH^*(Q_{F(\sqrt{a_1})}) \arrow{d}
\arrow[yshift=2pt]{l} \\
\CH^*(Q_{\ol{F}}) \arrow[yshift=-2pt]{r} & \CH^*(Q_{\ol{F}(\sqrt{a_1})}).
\arrow[yshift=2pt]{l}
\end{tikzcd}
\]

Стало быть, либо $\ol\rho\times 1$ рационален, либо
$\ol\rho\times 1 + 1\times\ol\rho$ рационален (в зависимости от четности $c$).

Предположим для начала, что $\ol\rho\times 1$ рационален.
Докажем, что в этом случае $Q$ изотропна.
Действительно, циклы $\ol\rho\times 1$ и $1\times\ol\rho$ рациональны,
а потому и $(\ol\rho\times 1)\cdot(1\times\ol\rho) = \ol\rho\times\ol\rho$
рационален.
Кроме того, $\ol{h}\times 1$ и $1\times\ol{h}$ рациональны,
а потому и $\ol{\pt}\times\ol{\pt}$ рационален.
Рассмотрим пушфорвард относительно проекции $Q\times Q$ на первый сомножитель:
\begin{align*}
\CH^*(\ol{Q}\times\ol{Q}) &\to \CH^*(\ol{Q}),
\ol{\pt}\times\ol{\pt} &\mapsto \ol{\pt}.
\end{align*}
Поэтому и цикл $\ol{\pt}$ рационален.
Значит, на $Q$ есть $0$-цикл степени $1$.
По теореме Спрингера из этого следует, что на $Q$ есть рациональная точка,
то есть, $Q$ изотропна~--- и это неинтересный случай.

Значит, на самом деле цикл $\ol\rho\times 1 + 1\times\ol\rho$
рационален.
Из него можно постараться изготовить проектор.
Заметим, что для любых $i,j$ цикл
$(\ol{h}^i\ol\rho)\times \ol{j}^j + \ol{h}^i\times (\ol{h}^j\ol\rho)$
тоже рационален.
Как подобрать $i,j$, чтобы это был проектор?
Заметим, что $\ol\rho$ лежит в коразмерности $2^{k-1} - 1$,
поэтому нужно, чтобы $j = 2^{k-1} - 1 - i$.
Оказывается, этого достаточно: нужно вспомнить формулу
$(a\times b)(c\times d) = \deg(ad)c\times b$
и равенство $\ol{h}^{2^{k-1}-1}\ol\rho = \ol{\pt}$.
После этого прямое вычисление показывает, что
мы получили проектор.
Варьируя $i$, получаем $2^{k-1} - 1$ проекторов.

Соответствующее разложение мотива $Q$ можно нарисовать так:
\[
\begin{tikzpicture}[scale=1.0,thin]
\def\sm{0.7}
\coordinate (1) at ($\sm*(0, 0)$);
\coordinate (2) at ($\sm*(1.4, 0)$);
\coordinate (3) at ($\sm*(2.8, 0)$);
\coordinate (3p) at ($\sm*(3.3, 0)$);
\coordinate (d1) at ($\sm*(3.5, 0)$);
\coordinate (d2) at ($\sm*(3.7, 0)$);
\coordinate (d3) at ($\sm*(3.9, 0)$);
\coordinate (4m) at ($\sm*(4.1, 0)$);
\coordinate (4) at ($\sm*(4.6, 0)$);
\coordinate (5) at ($\sm*(6.0, 0)$);
\coordinate (6) at ($\sm*(7.0, 1)$);
\coordinate (7) at ($\sm*(7.0, -1)$);
\coordinate (8) at ($\sm*(8.0, 0)$);
\coordinate (9) at ($\sm*(9.4, 0)$);
\coordinate (9p) at ($\sm*(9.9, 0)$);
\coordinate (e1) at ($\sm*(10.1, 0)$);
\coordinate (e2) at ($\sm*(10.3, 0)$);
\coordinate (e3) at ($\sm*(10.5, 0)$);
\coordinate (10m) at ($\sm*(10.7, 0)$);
\coordinate (10) at ($\sm*(11.2, 0)$);
\coordinate (11) at ($\sm*(12.6, 0)$);
\coordinate (12) at ($\sm*(14.0, 0)$);
\draw (1)--(2)--(3)--(3p);
\draw (4m)--(4)--(5)--(6);
\draw (5)--(7);
\draw (6)--(8);
\draw (7)--(8)--(9)--(9p);
\draw (10m)--(10)--(11)--(12);
\draw[dotted] (1) edge[bend left=45](6);
\draw[dotted] (2) edge[bend left=35](8);
\draw[dotted] (3) edge[bend left=45](9);
\draw[dotted] (4) edge[bend right=45](10);
\draw[dotted] (5) edge[bend right=35](11);
\draw[dotted] (7) edge[bend right=45](12);
\foreach \point in
{1,2,3,4,5,6,7,8,9,10,11,12}
  {
     \fill [white] (\point) circle (2.0pt);
     \draw [black] (\point) circle (2.0pt);
  }
\foreach \point in
{d1,d2,d3,e1,e2,e3}
  {
     \fill [black] (\point) circle (0.7pt);
  }
\end{tikzpicture}
\]
Иными словами, над базой слагаемые в разложении мотива квадрики
объединяются в пары.

\begin{exercise}
Пусть $R$~--- коммутативное кольцо, $I\trleq R$~--- идеал, состоящий
из нильпотентных элементов.
Тогда любой обратимый элемент $R/I$ поднимается до обратимого элемента в $R$.
\end{exercise}

Слагаемое, которое дает первый проектор из этих,
называется \term{мотивом Роста} и обозначается через $R$.
Таким образом, над замыканием
$\ol{R} = \mathbb{Z} \oplus \mathbb{Z}\{2^{k-1}-1\}$.
Мотив квадрики Пфистера, таким образом, составлен
из сдвигов мотива Роста:
\[
M(Q) = R \oplus R\{1\} \oplus \dots \oplus R\{2^{k-1}-1\}.
\]
\begin{fact}
Пусть $\lAngle a_1,\dots,a_k\rAngle$,
$\lAngle a'_1,\dots,a'_k\rAngle$~--- две $k$-формы Пфистера.
Соответствующие этим квадрикам мотивы Роста изоморфны (в категории мотивов)
тогда и только тогда, когда сами формы изоморфны,
что в свою очередь равносильно равенству чашечных произведений
$(a_1)\cup\dots\cup(a_k) = (a'_1)\cup\dots\cup(a'_k)$
в $H^k(F, \mathbb{Z}/2\mathbb{Z})$.
\end{fact}

\begin{remark}
Мотив Роста над замыканием превращается в
$\mathbb{Z} \oplus \mathbb{Z}\{2^{k-1}-1\}$.
Верно и обратное:
если мотив над замыканием выглядит так, то это мотив Роста
(теорема Никиты Семенова).
\end{remark}

% 02.04.2012

\subsection{Пример: $\mathsf{F}_4$}

Над замкнутым полем группа типа $\mathsf{F}_4$~--- это автоморфизмы
эрмитовых матриц $3\times 3$ над октонионами:
$\mathsf{F}_4 = \Aut(H_3(\mathbb{O}))$.

В общем случае приведем сначала <<конструкцию по модулю $2$>>.
Вместо $\mathbb{O}$ нужно взять другие октонионы
(они задатся формой Пфистера $\lAngle a,b,c\rAngle$),
и диагональную эрмитову форму вида $\la 1, -d, -e\ra$.
Здесь $a,b,c,d,e\in F^* / (F^*)^2$.
Если у поля $F$ нет расширений нечетной степени, то любая группа
типа $\mathsf{F}_4$ так выглядит.

Например, над $\mathbb{R}$ есть три группы типа $\mathsf{F}_4$:
\begin{enumerate}
\item построенная по расщепимым октонионам (и тогда неважно, каковы $d,e$);
\[
\begin{tikzpicture}[scale=1.0,thin]
\def\sm{0.7}
\coordinate (1) at ($\sm*(0, 0)$);
\coordinate (2) at ($\sm*(1.4, 0)$);
\coordinate (3) at ($\sm*(2.8, 0)$);
\coordinate (4) at ($\sm*(4.2, 0)$);
\draw (1)--(2);
\draw (3)--(4);
\draw (2) edge[transform canvas={yshift=2pt}] (3);
\draw (2) edge[transform canvas={yshift=-2pt}] (3);
\draw ($\sm*(2.0, 0.2)$) -- ($\sm*(2.2, 0)$) -- ($\sm*(2.0, -0.2)$);
\foreach \point in
{1,2,3,4}
  {
     \fill [white] (\point) circle (2.0pt);
     \draw [black] (\point) circle (2.0pt);
     \draw [black] (\point) circle (5.0pt);
  }
\end{tikzpicture}
\]
\item построенная по компактным октонионам (\term{октавам}) и $d=e=1$;
\[
\begin{tikzpicture}[scale=1.0,thin]
\def\sm{0.7}
\coordinate (1) at ($\sm*(0, 0)$);
\coordinate (2) at ($\sm*(1.4, 0)$);
\coordinate (3) at ($\sm*(2.8, 0)$);
\coordinate (4) at ($\sm*(4.2, 0)$);
\draw (1)--(2);
\draw (3)--(4);
\draw (2) edge[transform canvas={yshift=2pt}] (3);
\draw (2) edge[transform canvas={yshift=-2pt}] (3);
\draw ($\sm*(2.0, 0.2)$) -- ($\sm*(2.2, 0)$) -- ($\sm*(2.0, -0.2)$);
\foreach \point in
{1,2,3,4}
  {
     \fill [white] (\point) circle (2.0pt);
     \draw [black] (\point) circle (2.0pt);
  }
\draw [black] (4) circle (5.0pt);
\end{tikzpicture}
\]
\item построенная по октавам и $d = e = -1$~--- она анизотропна
(над $\mathbb{R}$ это равносильно компактности).
\[
\begin{tikzpicture}[scale=1.0,thin]
\def\sm{0.7}
\coordinate (1) at ($\sm*(0, 0)$);
\coordinate (2) at ($\sm*(1.4, 0)$);
\coordinate (3) at ($\sm*(2.8, 0)$);
\coordinate (4) at ($\sm*(4.2, 0)$);
\draw (1)--(2);
\draw (3)--(4);
\draw (2) edge[transform canvas={yshift=2pt}] (3);
\draw (2) edge[transform canvas={yshift=-2pt}] (3);
\draw ($\sm*(2.0, 0.2)$) -- ($\sm*(2.2, 0)$) -- ($\sm*(2.0, -0.2)$);
\foreach \point in
{1,2,3,4}
  {
     \fill [white] (\point) circle (2.0pt);
     \draw [black] (\point) circle (2.0pt);
  }
\end{tikzpicture}
\]
\end{enumerate}

По общей теории скрученных форм над любым полем
группа типа $\mathsf{F}_4$~--- это группа автоморфизмов
алгебры $J$, где $J$~--- скрученная
форма йордановой алгебры $H_3(\mathbb{O})$.

Приведем теперь <<конструкцию по модулю $3$>>.
Пусть $D$~--- центральная простая алгебра степени $3$.
Они все циклические, поэтому
$D = (a,b)_3 = \la x,y\mid x^3 = a,\; y^3 = b,\; xy = \zeta yx\ra$
для некоторых $a,b\in F^*/(F^*)^3$.
Здесь $\zeta^3 = 1$.
Возьмем еще $c\in F^*/(F^*)^2$.
тогда на $D\oplus D\oplus D$ можно завести структуру йордановой
алгебры $J(a,b,c)$ (с помощью скаляра $c$).
Ее норма выглядит так:
\[
N(\alpha\oplus\beta\oplus\gamma) = \Nrd(\alpha) + c\Nrd(\beta)
+ c^{-1}\Nrd(\gamma) - \Trd(\alpha\beta\gamma),
\]
где $\Trd$~--- приведенный след.
Автоморфизмы этой нормы образуют группу типа $\mathsf{E}_6$,
а подгруппа в ней, сохраняющая
единицу (то есть, $1\oplus 0 \oplus 0$)~--- это группа типа
$\mathsf{F}_4$.

\begin{conjecture}[Ослабленный вариант гипотезы Серра--Роста]
Полученная группа типа $\mathsf{F}_4$ зависит только
от $\{a,b,c\} \in K_3^M(F)/3$
(или, что то же самое, от
$(a)\cup (b)\cup (c) \in H^3(F,\mathbb{Z}/3\mathbb{Z})$).
\end{conjecture}
Известно, что если $J(a,b,c) = J(a',b',c')$, то
$\{a,b,c\} = \{a',b',c'\}$.

Если у поля нет квадратичных расширений, то любая скрученная
форма $H_3(\mathbb O)$ имеет вид
$J(a,b,c)$.

Для любого поля $F$ определен инвариант
\[
g_3\colon H^1(F, \mathsf{F}_4) \to H^3(F, \mathbb{Z}/3\mathbb{Z}).
\]
\begin{itemize}
\item Если у $F$ нет квадратичных расширений, то образ $g_3$~--- это
в точности чистые символы $(a)\cup (b)\cup (c)$.
\item В этом случае ядро тривиально.
\item Гипотеза состоит в том, что $g_3$ инъективно.
\end{itemize}
Мы построили алгебру $J(a,b,c)$ и группу $G = \Aut(J(a,b,c))$.
\begin{fact}
Группа $G$ или расщепима, или анизотопна (как и в случае группы
изометрий пфистеровых форм).
\end{fact}
В частности, если $X$~--- $G$-однородное проективное многообрази,
то оно является клеточным над общей точкой,
то есть, $X_{F(X)}$ клеточное.
Пусть $X$~--- скрученная форма $\mathsf{F}_4/P_4$.

Отступление: J.-P. Bonnet показал, что
$M({}_{\xi}(\mathsf{G}_2/P_1)) \isom M({}_{\xi}(\mathsf{G}_2/P_2))$.
Многообразия $\mathsf{G}_2/P_1$ и $\mathsf{G}_2/P_2$
оба имеют размерность $5$.
На самом деле, $\mathsf{G}_2/P_1$~--- квадрика.
Более того, это максимальный сосед квадрики Пфистера, и поэтому
ее мотив раскладывается на мотивы Роста.
Напомним, что $\mathsf{G}_2 = \Aut(\mathbb{O})$,
где $\mathbb{O}$ задается формой $\lAngle a, b, c\rAngle$.
Возникающий мотив Роста отвечает как раз квадрике
$\lAngle a, b, c\rAngle$.
Вопрос: что если взять $\mathsf{F}_4/P_1$ и $\mathsf{F}_4/P_4$?
У них тоже одинаковая размерность и многочлен Пуанкаре.
Теорема Зайнуллина--Николенко--Семенова гласит,
что
$M({}_{\xi}(\mathsf{F}_4/P_1)) \isom M({}_{\xi}(\mathsf{F}_4/P_2))$,
если рассматриваемая группа типа $\mathsf{F}_4$ имеет вид
$\Aut(J(a,b,c))$.

Размерность $\mathsf{F}_4/P_4$ равна $15$.
Нарисуем диаграмму Хассе для этого многообразия.
С точностью до каких-то ребер внутри она выглядит так:
\[
\begin{tikzpicture}[scale=1.0,thin,font=\scriptsize]
\def\sm{0.7}
\coordinate (1) at ($\sm*(0, 0)$);
\coordinate (2) at ($\sm*(1.4, 0)$);
\coordinate (3) at ($\sm*(2.8, 0)$);
\coordinate (4) at ($\sm*(4.2, 0)$);
\coordinate (5) at ($\sm*(5.2, 1)$);
\coordinate (6) at ($\sm*(6.6, 1)$);
\coordinate (7) at ($\sm*(8.0, 1)$);
\coordinate (8) at ($\sm*(9.4, 1)$);
\coordinate (9) at ($\sm*(10.8, 1)$);
\coordinate (10) at ($\sm*(12.2, 1)$);
\coordinate (11) at ($\sm*(13.6, 1)$);
\coordinate (12) at ($\sm*(15.0, 1)$);
\coordinate (5x) at ($\sm*(5.2, -1)$);
\coordinate (6x) at ($\sm*(6.6, -1)$);
\coordinate (7x) at ($\sm*(8.0, -1)$);
\coordinate (8x) at ($\sm*(9.4, -1)$);
\coordinate (9x) at ($\sm*(10.8, -1)$);
\coordinate (10x) at ($\sm*(12.2, -1)$);
\coordinate (11x) at ($\sm*(13.6, -1)$);
\coordinate (12x) at ($\sm*(15.0, -1)$);
\coordinate (13) at ($\sm*(16.0, 0)$);
\coordinate (14) at ($\sm*(17.4, 0)$);
\coordinate (15) at ($\sm*(18.8, 0)$);
\coordinate (16) at ($\sm*(20.2, 0)$);
\node at ($\sm*(0, -2)$) {$0$};
\node at ($\sm*(1.4, -2)$) {$1$};
\node at ($\sm*(2.8, -2)$) {$2$};
\node at ($\sm*(4.2, -2)$) {$3$};
\node at ($\sm*(5.2, -2)$) {$4$};
\node at ($\sm*(6.6, -2)$) {$5$};
\node at ($\sm*(8.0, -2)$) {$6$};
\node at ($\sm*(9.4, -2)$) {$7$};
\node at ($\sm*(10.8, -2)$) {$8$};
\node at ($\sm*(12.2, -2)$) {$9$};
\node at ($\sm*(13.6, -2)$) {$10$};
\node at ($\sm*(15.0, -2)$) {$11$};
\node at ($\sm*(16.0, -2)$) {$12$};
\node at ($\sm*(17.4, -2)$) {$13$};
\node at ($\sm*(18.8, -2)$) {$14$};
\node at ($\sm*(20.2, -2)$) {$15$};
\node at (2) [above=3pt] {$h$};
\node at (3) [above=3pt] {$h^2$};
\node at (4) [above=3pt] {$h^3$};
\node at (5) [above=3pt] {$\rho$};
\node at (9) [above=3pt] {$\rho^2$};
\draw (1)--(2)--(3)--(4)--(5)--(6)--(7)--(8)--(9)--(10)--(11)--(12)--(13)--(14)--(15)--(16);
\draw (4)--(5x)--(6x)--(7x)--(8x)--(9x)--(10x)--(11x)--(12x)--(13);
\foreach \point in
{1,2,3,4,5,6,7,8,9,10,11,12,5x,6x,7x,8x,9x,10x,11x,12x,13,14,15,16}
  {
     \fill [white] (\point) circle (2.0pt);
     \draw [black] (\point) circle (2.0pt);
  }
\end{tikzpicture}
\]

\begin{enumerate}
\item Берем образующую в $\CH^1(\mathsf{F}_4/P_4)$.
Она рациональная, то есть, лежит в образе отображения
\[
\CH^1({}_{\xi}(\mathsf{F}_4/P_4)) \xrightarrow{\res}
\CH^1(\mathsf{F}_4/P_4).
\]
В $\mathsf{F}_4$ решетка весов совпадает с решеткой корней.
В частности, $\varpi_4$~--- корень.
\item Тогда $h^2$, $h^3$~--- образующие $\CH^2$ и $\CH^3$.
\item Далее, $h^4$ и $\rho$~--- базис для $\CH^4$.
\item Кроме того, $\rho^2 h^7 = [\pt]\pmod{3}$.
\item $X_{F(X)}$ клеточное.
\item Если $\Aut(J(a,b,c))$ расщепляется над расширением
степени, взаимно простой с $3$, то она и была расщепимой
(это теорема типа Спрингера).
\end{enumerate}

Начинаем применять метод общей точки:
\[
\begin{tikzcd}
\CH^*(X\times X) \arrow[->>]{r} \arrow{d}{\res} &
\CH^*(X_{F(X)}) \arrow{d}{\isom} \\
\CH^*(X_{\ol{F}}\times X_{\ol{F}}) \arrow[->>]{r} &
\CH^*(X_{\ol{F}(X)}).
\end{tikzcd}
\]
Элемент $\ol{\rho}\in\CH^*(X_{\ol{F}(X)})$ поднимается
до элемента в $\CH^*(X\times X)$.
Поэтому в $\CH^*(X_{\ol{F}}\times X_{\ol{F}})$ есть
рациональный элемент вида
\[
\alpha = \rho\times 1 + ? \cdot h^3\times h
+ ?\cdot h^2\times h^2 + ?\cdot h\times h^3
+ ?\cdot 1\times h^4 + c\cdot 1\times \rho.
\]
Подправив $\alpha$ на рациональные элементы (степени $h$),
можно добиться, что останется только
$\rho\times 1 + c\cdot 1\times\rho$.
Поскольку $X$ над кубическим расширением становится клеточным,
можно считать, что $c = 0$ или $c = \pm 1$.
Почему $c\neq 0$?
Если $\rho\times 1$ рационален, то и $1\times\rho$ рационален,
откуда $\rho^2\times\rho^2$ и $h^7\rho^2\times h^7\rho^2$
рационален, а потому и $\pt\times\pt$ рационален.
Применяя пушфорвард, видим, что класс
$\pt\in\CH^*(X)$ рационален~--- противоречие.

Поэтому на самом деле класс $\rho\times 1\pm 1\times\rho$
рационален.
Значит, $\rho^2\times 1 \pm 2\rho\times\rho + 1\times\rho^2$
рационален, а потому и
$\beta = \rho^2 + 1 \mp \rho\times\rho + 1\times\rho^2$
рационален.
Будем умножать полученный цикл на $h^i\times h^j$,
где $i+j=7$.
Получится цикл
\[
h^i\rho^2\times h^j \mp h^i\rho\times h^j\rho
+h^i\times h^j\rho^2\in\CH^{15}(X_{\ol{F}}\times X_{\ol{F}}).
\]
Возьмем его композицию с самим собой, и воспользуемся сравнением.
$h^{i+j}\rho^2 \equiv \pt \pmod{3}$.
Вообще, будем считать все по модулю $3$:
\[
h^i\rho^2\times h^j + h^i\rho\times h^j\rho + h^i\times h^j\rho^2
\]
Значит, этот цикл уже является проектором по модулю $3$.
Над замыканием каждое его слагаемое, конечно, является проектором~---
но не рациональным.

Итак, мы получили восемь ортогональных проекторов
$p_0,\dots,p_7$.
Правые части этих проекторов (24 штуки) образуют базис в группе Чжоу,
но не тот, который у нас был раньше (клетки Шуберта).
\[
\begin{tikzpicture}[scale=1.0,thin,font=\scriptsize]
\def\sm{0.7}
\coordinate (1) at ($\sm*(0, 0)$);
\coordinate (2) at ($\sm*(1.4, 0)$);
\coordinate (3) at ($\sm*(2.8, 0)$);
\coordinate (4) at ($\sm*(4.2, 0)$);
\coordinate (5) at ($\sm*(5.2, 1)$);
\coordinate (6) at ($\sm*(6.6, 1)$);
\coordinate (7) at ($\sm*(8.0, 1)$);
\coordinate (8) at ($\sm*(9.4, 1)$);
\coordinate (9) at ($\sm*(10.8, 1)$);
\coordinate (10) at ($\sm*(12.2, 1)$);
\coordinate (11) at ($\sm*(13.6, 1)$);
\coordinate (12) at ($\sm*(15.0, 1)$);
\coordinate (5x) at ($\sm*(5.2, -1)$);
\coordinate (6x) at ($\sm*(6.6, -1)$);
\coordinate (7x) at ($\sm*(8.0, -1)$);
\coordinate (8x) at ($\sm*(9.4, -1)$);
\coordinate (9x) at ($\sm*(10.8, -1)$);
\coordinate (10x) at ($\sm*(12.2, -1)$);
\coordinate (11x) at ($\sm*(13.6, -1)$);
\coordinate (12x) at ($\sm*(15.0, -1)$);
\coordinate (13) at ($\sm*(16.0, 0)$);
\coordinate (14) at ($\sm*(17.4, 0)$);
\coordinate (15) at ($\sm*(18.8, 0)$);
\coordinate (16) at ($\sm*(20.2, 0)$);
\draw (1)--(2)--(3)--(4)--(5)--(6)--(7)--(8)--(9)--(10)--(11)--(12)--(13)--(14)--(15)--(16);
\draw (4)--(5x)--(6x)--(7x)--(8x)--(9x)--(10x)--(11x)--(12x)--(13);
\draw[dotted] (1) edge[bend left=45] (5);
\draw[dotted] (5) edge[bend left=45] (9);
\draw[dotted] (2) edge[bend left=45] (6);
\draw[dotted] (6) edge[bend left=45] (10);
\node at ($\sm*(1,1)$) {$p_0$};
\node at ($\sm*(2.4,1)$) {$p_1$};
\foreach \point in
{1,2,3,4,5,6,7,8,9,10,11,12,5x,6x,7x,8x,9x,10x,11x,12x,13,14,15,16}
  {
     \fill [white] (\point) circle (2.0pt);
     \draw [black] (\point) circle (2.0pt);
  }

\end{tikzpicture}
\]
Можно также рассмотреть $\beta(h^i\times h^j)$ для произвольных
$i,j$.
Такие рациональные циклы индуцируют изоморфизмы
между $(X,p_0)\{i\}$ и $(X,p_i)$.

Мы посчитали все по модулю $3$, но можно добиться и
целочисленных проекторов.
На самом деле, при каких-то разумных условиях всегда можно
поднять проекторы (и изоморфизмы) по модулю $2$ и $3$
до целочисленных.

Мотив Роста квадрики Пфистера $\lAngle a_1,\dots,a_k\rAngle$
обозначается так: $R_{2,k}(\{a_1,\dots,a_k\})$.
Полученный выше мотив $(X,p_0)$ будем обозначать через
$R_{3,3}(\{a,b,c\})$.

\begin{corollary}
\[
M(X) = \bigoplus_{i=0}^{7} R_{3,3}(\{a,b,c\})\{i\}.
\]
\end{corollary}

\begin{remark}
Точно так же доказывается, что
\[
M({}_{\xi}(\mathsf{F}_4/P_1)) = \bigoplus_{i=0}^{7}
R_{3,3}(\{a,b,c\})\{i\}.
\]
Более того, мотив любого $G$-однородного проективного
многообразия $Y$ раскладывается в сумму сдвигов
мотива $R_{3,3}(\{a,b,c\})$.
Для квадрики Пфистера верно аналогичное замечание.
\end{remark}

\begin{theorem}[Зайнуллин--Петров--Семенов]\label{thm:ZPS}
Пусть $G$~--- полупростая алгебраическая группа над $F$,
$X$~--- $G$-однородное проективное многообразие, клеточное
над общей точкой, $p$~--- простое число.
Тогда мотив $M(X)\otimes\mathbb{Z}/p\mathbb{Z}$ изоморфен
сумме $R_p(G)$ с какими-то сдвигами,
где $R_p(G)$ не зависит от $X$ и неразложим по модулю $p$.
\end{theorem}
Нужно пояснить, что такое $M(X)\otimes\mathbb{Z}/p\mathbb{Z}$:
в конструкции категории соответствий нужно взять
$\Mor(X,Y) = \CH^{\dim Y}(X\times Y) / p$.

\begin{remark}
Условие клеточности $X_{F(X)}$ можно заменить
на требование расщепимости $G_{F(X)}$.
\end{remark}
\begin{remark}
Мотив $R_p(G)$ над алгебраическим замыканием раскладывается
в сумму мотивов вида $\mathbb{Z}/p = M(\pt)\otimes\mathbb{Z}/p$
со сдвигами, и для этих сдвигов есть некоторая формула.
\end{remark}
\begin{remark}
Все ситуации, описанные в теореме, перечислены в работе
Петрова--Семенова.
\end{remark}
\begin{remark}
Сами проекторы можно поднять в $\mathbb{Z}$, но не изоморфизмы
между ними.
\end{remark}
\begin{remark}
Если $G$ не содержит сомножителей типа $\mathsf{A}$
и $p\neq 2,3,5$, то $R_p(G) = \mathbb{Z}/p$.
Случай $p=5$ возникает только для $\mathsf{E}_8$.
\end{remark}
\begin{remark}
Мотив $M(X)\otimes\mathbb{Q}$ раскладывается в
прямую сумму мотивов $M(\pt)\otimes\mathbb{Q}$ со сдвигами.
\end{remark}
\begin{remark}
Если $G=\mathsf{E}_8$, $p=2$, и инвариант Роста
тривиален, то все $G$-однородные проективные многообразия $X$
подходят.
В этом случае $R_2(G)$~--- мотив Роста,
отвечающий $5$-символу (см. работу Никиты Семенова
про конечные подгруппы $\mathsf{E}_8$).
\end{remark}

% 09.04.2012

\subsection{$J$-инвариант}

Остается открытым вопрос, каков размер $R_p({}_{\xi}G)$.
Мы знаем, что над замыканием получается
$R_p({}_{\xi}G)_{\ol{F}} = \bigoplus\mathbb{Z}/p\{\dots\}$.
Как посчитать эти сдвиги?
Можно закодировать их многочленом: по прямой сумме
$\bigoplus_i\mathbb{Z}/p\{i\}^{\oplus a_i}$ построим
\term{полином Пуанкаре} $P(R_p(G),t) = \sum_i a_i t^i$.
Понятно, что этот многочлен контролирует
образ отображения
$\CH^*(X)/p \to \CH^*(X_{\ol{F}})/p$.

Сейчас мы построим набор целых чисел $J_p({}_{\xi}G)$~---
\emph{$J$-инвариант}~--- со следующими свойствами:
\begin{enumerate}
\item $P(R_p({}_{\xi}G),t)$ выражается через $J_p({}_{\xi}G)$;
\item $J_p({}_{\xi}G)$ контролирует, какие ${}_{\xi}G$-однородные
проективные многообразия действительно являются клеточными
над общей точкой;
\item для многих исключительных групп $J_p({}_{\xi}G)$ выражается
через индекс Титса.
\end{enumerate}
Пусть $B\leq G$~--- борелевская подгруппа.
Рассмотрим <<последовательность>> \emph{чего-то}
\[
B \to G \to G/B \to BB = \pt//B.
\]
Переходя к кольцам Чжоу, получаем точную последовательность
градуированных колец вида
\[
\CH^*(BB) \to \CH^*(G/B) \to \CH^*(G) \to \CH^*(B).
\]
При этом $\CH^*(B) = \mathbb{Z}$.
Заметим, что $\CH^*(G/B) \to \CH^*(G)$~--- сюръекция.
Поэтому есть точная последовательность
\[
\CH^*(BB) \to \CH^*(G/B) \to \CH^*(G) \to 0.
\]
Борелевская подгруппа гомотопически эквивалентна тору:
$B \sim \mathbb{G}_m\times \dots \times \mathbb{G}_m$.
Покажем, что $\CH^*(BB) = S^*(X^*(T))$.

Что такое $B\mathbb{G}_m$?
Первое приближение:
$\mathbb{P}^n = (\mathbb{A}^{n+1}\setminus\{0\})/\mathbb{G}_m$.
Правильный ответ: $B\mathbb{G}_m = \varinjlim_n\mathbb{P}^n$.
Мы уже знаем, что $\CH^*(\mathbb{P}^n) = \mathbb{Z}[x] / (x^{n+1})$.
Поэтому $\CH^*(B\mathbb{G}_m) = \mathbb{Z}[x]$~---
проективный предел в категории градуированных колец
(но не в категории колец).

Пусть $\chi_1,\dots,\chi_l$~--- базис решетки $X^*(T)$.
Если $G$ односвязна, можно взять $\varpi_1,\dots,\varpi_l$.
Если $G$ присоединенная, можно взять $\alpha_1,\dots,\alpha_l$.
Тогда $S^*(X^*(T)) \isom \mathbb{Z}[\chi_1,\dots,\chi_l]$.

Чтобы описать отображение $\CH^*(BB) \to \CH^*(G/B)$,
достаточно задать образы элементов $\chi_i$.
Положим $\chi_i\mapsto [L_{\chi_i}]$ (класс линейного
расслоения $L_{\chi_i}$ в группе Пикара).
Здесь $L_{\chi_i}$~--- линейное расслоение на $G/B$,
построенное следующим образом.
Характер $\chi_i$ является отображением
$B\to \mathbb{G}_m$.
Тогда $L_{\chi_i} = G\times_{B}\mathbb{A}^1$,
где на $\mathbb{A}^1$ задано действие с помощью $\chi_i$.
Каноническое отображение
$G\times_{B}\mathbb{A}^1 \to G\times_{B}\pt = G/B$
превращает $L_{\chi_i}$ в линейное расслоение.

Пока что $G$ была расщепима.
Оказывается, если подкрутить все на $\xi$, все
$[L_{\chi_i}]$ останутся рациональными.

\begin{example}
Рассмотрим $\mathbb{P}^1 = \SL_2/B$ и обозначим
характер тора через $\xi$.
Нас интересует действие $B$ на $\SL_2\times_{B}\mathbb{A}^1$,
при котором
матрица $\begin{pmatrix}\alpha & * \\ 0 & \alpha^{-1}\end{pmatrix}$
действует на $\mathbb{A}^1$ умножением на $\alpha$.
Упражнение: $L_{\xi} = \mathcal{O}(-1)$,
$L_{\xi^{-1}} = \mathcal{O}(1)$.

В общем случае в $G/B$ есть клетки Шуберта коразмерности $1$.
Пусть $G$ односвязна.
Эти клетки соответствуют фундаментальным характерам:
клетке $\chi_i$ сопоставляется $i$-ая клетка Шуберта
коразмерности $1$, равная $c_1(L(\chi_i))$.
\end{example}

Теперь мы взяли $\xi\in H^1(F,G)$.
Тогда ${}_{\xi}(G/B)$~--- многообразие, клеточное над общей точкой
(при переходе к его полю функций у $G$ появляется борелевская
подгруппа, поэтому можно написать фильтрацию).
Нас интересует образ отображения
\[
\CH^*({}_{\xi}(G/B)) \xrightarrow{\res} \CH^*(G/B).
\]
Все, что приходит с $\CH^*(BB)$, лежит в образе $\res$,
поскольку линейные расслоения можно скрутить:
$\res([{}_{\xi}L_{\chi}]) = [L_{\chi}]$.
Точность сохранится при факторизации по $p$:
\[
\CH^*(BB)/p \to \CH^*(G/B)/p \to \CH^*(G)/p \to 0.
\]
При этом $\CH^*(BB)/p \isom \mathbb{Z}/p[x_1,\dots,x_l]$.
Умножение $G\times G\to G$ дает нам отображение
$\CH^*(G) \to \CH^*(G) \otimes \CH^*(G)$,
которое превращает $\CH^*(G)$ в алгебру Хопфа.
Таким образом, $\CH^*(G)/p$~--- алгебра Хопфа над $\mathbb{Z}/p$,
градуированная, конечномерная, связная, коммутативная.
\begin{theorem}
Все такие алгебры Хопфа изоморфны (как алгебры)
$\mathbb{Z} / p [x_1,\dots,x_r] / (x_i^{p^{k_i}})$.
\end{theorem}
Если $X$~--- клеточное над $\mathbb{C}$, можно сравнить
$\CH^*(X)$ и $H^*_{\sing}(X)$.
Оказывается, $\CH^i(X) = H^{2i}_{\sing}(X)$.
Для $G$ над $\mathbb{C}$ есть нетривиальные элементы
в $H^{2i+1}_{\sing}(G)$
Например, $G = \SL_2(\mathbb{C})$, и $\SL_2(\mathbb{C})$
гомотопически эквивалентно $S^3$.
\begin{remark}
Если $p\neq 2,3,5$, и $G$ не изогенична $\SL_n$,
то $\CH^*(G)/p = \mathbb{Z}/p$.
Случай $p=5$ возникает только для $\mathsf{E}_8$,
а случай $p=3$~--- только для $\mathsf{F}_4$,
$\mathsf{E}_6$, $\mathsf{E}_7$, $\mathsf{E}_8$
(это делители числа Кокстера).
\end{remark}
\emph{Таблица Каца} дает для каждой $G$ и для каждого $p$
значения $k_i$ и степени элементов $x_i$.
Обозначим $d_i = \deg(x_i)$.

Заметим, что ${}_{\xi}(G/B) = ({}_{\xi}G)/B$; это дает короткую
точную последовательность
\[
{}_{\xi}G \to ({}_{\xi}G)/B \to BB,
\]
из которой получаем стрелку $\CH^*(BB)/p \to \CH({}_{\xi}(G/B))/p$.
Рассмотрим коммутативную диаграму
\[
\begin{tikzcd}
& \CH({}_{\xi}(G/B))/p \arrow{d} \arrow{dr}{\ph} \\
\CH^*(BB)/p \arrow{ur} \arrow{r}
& \CH^*(G/B) \arrow[->>]{r}
& \CH^*(G)/p \arrow[equal]{d} \\
& & (\mathbb{Z}/p)[x_i] / (x^{p^{k_i}}).
\end{tikzcd}
\]
Нас интересует образ вертикальной стрелки
в $\CH^*(G/B)/p$.
Обозначим через $j_i$ наименьшее целое число такое, что
$x_i^{p^{j_i}} + \mbox{члены меньшего порядка} \in \im(\ph)$.
Порядок мы понимаем в смысле Deglex;
$x_1\leq x_2\leq\dots\leq x_r$, если
$d_1\leq d_2\leq\dots\leq d_r$.
Можно рассмотреть $\CH^(G)/p$ как комодуль над собой,
и тогда $\im(\ph)$ будет подкомодулем.

Заметим, что $0\leq j_i \leq k_i$, так как $x_i^{p^{k_i}} = 0$.
Равенство $j_i = 0$ равносильно тому, что
$x_i + \mbox{члены меньшего порядка}\in\im(\ph)$.
Набор чисел $(j_i)$ обозначим через $J_p(\xi)$
(он действительно зависит только от $\xi$, но не от ${}_{\xi}G$).
Тогда полином Пуанкаре выглядит так:
\[
P(R_p(G), t) = \prod_{i=1}^r\frac{1-t^{p^{j_i}\cdot d_i}}{1-t^{d_i}}.
\]
\begin{example}
Рассмотрим группу типа $\mathsf{F}_4$, $p=3$.
Тогда $\CH^*(G) = (\mathbb{Z}/3)[x_1]/(x_1^3)$,
где $d_1 = \deg x_1 = 4$.
Таким образом, $k_1=1$, и для $j_1$ есть два варианта: $0$ и $1$.
\begin{itemize}
\item случай $J_p(\xi) = (0)$ неинтересен
(см. замечание: полином Пуанкаре равен $1$;
\item в случае $J_p(\xi) = (1)$ получаем полином Пуанкаре
\[\frac{1-t^{3\cdot 4}}{1-t^4} = 1 + t^4 + t^8.\]
\end{itemize}
\end{example}
\begin{remark}
$J_p(\xi) = 0$ тогда и только тогда, когда ${}_{\xi}G$ расщепляется
расширением степени, взаимно простой с $p$.
\end{remark}

\begin{example}
Пусть $G$~--- группа типа $\mathsf{G}_2$, $p=2$.
Тогда $\CH^*(G) = (\mathbb{Z}/2)[x_1]/(x_1^2)$, и $\deg x_1 = 3$.
Снова два случая:
\begin{itemize}
\item неинтересный: $J_p(\xi) = (0)$;
\item $J_p(\xi) = (1)$; полином Пуанкаре равен
\[\frac{1-t^{2\cdot 3}}{1-t^3} = 1 + t^3.\]
В этом случае $R_2({}_{\xi}(G))$~--- мотив Роста.
\end{itemize}
\end{example}
\begin{example}
В случаях $\mathsf{F}_4$, $\mathsf{E}_6$ при $p=2$ ответ тот же,
что и для $\mathsf{G}_2$.
\end{example}
\begin{example}
В случаях $\mathsf{E}_6^{\operatorname{sc}}$, $\mathsf{E}_7$
при $p=3$ ответ тот же, что и для $\mathsf{F}_4$.
\end{example}
\begin{remark}
Степень полинома $P(R_p(G), t)$ равна $\sum (p^{j_i}-1)d_i$.
Оказывается, это равно $\operatorname{cd}_p(X)$ (\emph{каноническая размерность}).
Попросту говоря, это наименьшая из размерностей рациональных циклов.
\end{remark}
\begin{example}
Пусть $G$~--- группа типа $\mathsf{E}_8$, $p=5$.
Тогда $\CH^*(G) = (\mathbb{Z}/5)[x_1]/(x_1^5)$,
и $\deg x_1 = 5+1 = 6$ (вообще, $\deg x_i = p+1$, если $r=1$).
В этом случае любое ${}_{\xi}G$-однородное проективное многообразие
является клеточным над общей точкой.
В нетривиальном случае полином Пуанкаре равен
\[
\frac{1-t^{5\cdot 6}}{1-t^6}.
\]
\end{example}

\begin{example}
Рассмотрим группу типа $\mathsf{E}_8$, $p=2$.
Тогда
\[
\CH^*(G) = (\mathbb Z/2)[x_1,x_2,x_3,x_4]/(x_1^8,x_2^4,x_3^2,x_4^2),
\]
$\deg x_1=3, \deg x_2 = 5, \deg x_3 = 9, \deg x_4 = 15$.
Это можно увидеть так:
элементы $x_2$ и $x_3$ получаются операцией Стинрода из $x^1$б
а именно, $x_2 = S^2(x_1)$ и $x_3 = S^4(x_2)$.
При этом $\deg S^m(x) = \deg x + m\cdot(p-1)$.
Если $m=\deg x$, то $S^m$~--- возведение в степень $p$;
если же $m > \deg x$, то $S^m(x) = 0$.
Операции Стинрода удовлетворяют следующим тождествам:
\begin{itemize}
\item $S^m$ линейны;
\item $S^m(xy) = \sum_n S^n(x)S^{m-n}(y)$;
\item Adem relations.
\end{itemize}
Из первых двух соотношений следует, что $\sum_m S^m$~--- гомоморфизм
колец.

Что это означает для $J_p(\xi)$?
Мы знаем, что $J_p(\xi) = (j_1, j_2, j_3, j_4)$,
причем $0\leq j_1\leq 3$, $0\leq j_2\leq 2$, $0\leq j_3,j_4\leq 1$.
Из свойств $S^m$ следует, что $j_1\geq j_2\geq j_3$.
Кроме того, $j_1\leq j_2+1$ и $j_2\leq j_3+1$.
Если $j_1 = 3$, то $j_2=2$ и $j_3=1$.
Если же $j_1 = 0$, то $j_2=0$, $j_3=0$, и возникает
интересный случай, когда при этом $j_4=1$
(заметим, что из равенства $j_1=0$ следует,
что инвариант Роста по модулю $2$ тривиален).
Тогда
\[
P(R_2({}_{\xi}G), t) = \frac{1-t^{2\cdot 15}}{1-t^{15}} = 1 + t^{15},
\]
как у мотива Роста.
Как мы уже упоминали, Никита Семенов доказал, что из этого следует,
что это и есть мотив Роста.
Получается некоторый инвариант в $H^5(F,\mathbb{Z}/2)$.
\end{example}

% 16.04.2012

\subsection{Набросок доказательства теоремы~\ref{thm:ZPS}}

Во-первых, нам понадобится <<теорема Крулля--Шмидта>>.
О какой категории идет речь?
Зафиксируем ${}_{\xi}G$ и рассмотрим категорию мотивов
по модулю простого числа $p$, а в ней~--- псевдоабелеву
подкатегорию, порожденную мотивами ${}_{\xi}G$-однородных
многообразий.
То есть, мы берем замыкание относительно взятия прямых слагаемых и
свдигов (и тензорных произведений, хотя это не важно).

\begin{theorem}[Черноусов--Меркурьев]\label{thm:ChM}
В этой категории выполнена теорема Крулля--Шмидта:
если есть разложение объекта в прямую сумму неразложимых,
то оно единственно.
\end{theorem}
Пусть теперь $X,Y$~--- проективные однородные многообразия,
$X\xrightarrow{f} Y$~--- локально (по Зарискому) тривиальное
расслоение с клеточным слоем $F$.
В силу клеточности $F$, его мотив раскладывается в
сумму сдвигов мотивов Тейта:
\[
M(F) = \bigoplus\mathbb{Z}\{i\}.
\]
Тогда $M(X) = M(Y)\otimes M(F) = \bigoplus M(Y)\{i\}$.
Поэтому, в частности,
$\CH(X) = \CH(Y)\otimes\CH(F)$ как $\CH(Y)$-модуль.
Это можно понять явно:
если наше расслоение тривиализуется над $U\subseteq Y$,
то морфизм $U\times F \to U$ дает нам возможность
по элементу $a\in\CH^*(F)$ построить
цикл $1\times a$ и рассмотреть его замыкание в $X$.
Такие замыкания как раз и образуют $\CH(Y)$-базис в $\CH(X)$,
если в качестве $a$ брать элементы базиса $\CH(F)$.

Пусть теперь $X$~--- ${}_{\xi}G$-однородное проективное
многообразие: $X = {}_{\xi}(G/P)$.
Рассмотрим расслоение
${}_{\xi}(G/b) \to X$ со слоем $P/B$.
Над $F(X)$ наш торсор становится тривиальным,
поэтому расслоение выглядит как $G/B \to G/P$.
Оно локально тривиально по Зарискому, поэтому
это верно и над некоторым открытым $U\subseteq X$.
Теперь мы можем подействовать группой $G$ и покрыть все $X$
такими открытыми.
Значит, все отображение локально тривиально по Зарискому.
По теореме~\ref{thm:ChM} $M({}_{\xi}(G/B)) = M(X)\otimes M(P/B)$.
Поэтому, если $M({}_{\xi}(G/B))\otimes(\mathbb{Z}/p)$
раскладывается в сумму сдвигов $R_p(\xi)$, то и $M(X)$
тоже (по теореме Крулля--Шмидта).

Значит, достаточно для случая $X = {}_{\xi}(G/B)$ доказать,
что $M(X)\otimes(\mathbb{Z}/p)$ раскладывается в сумму
неразложимых кусков $R_p(\xi)$ со сдвигами.
Рассмотрим коммутативную диаграмму
\[
\begin{tikzcd}
\CH^*(BB)/p \arrow{r} \arrow[equal]{d}
& \CH^*({}_{\xi}(G/B))/p \arrow[->>]{r} \arrow{d}
& \CH^*(\xi)/p \arrow{d} \\
\CH^*(BB)/p \arrow{r}
& \CH^*(G/B)/p \arrow[->>]{r}
& \CH^*(G)/p.
\end{tikzcd}
\]
\begin{enumerate}
\item
Мы знаем, что над замыканием
$\CH^*(G)/p = (\mathbb{Z}/p)[x_1,\dots,x_r] / (x_i^{p^{k_i}})$,
и образ $\CH^*(BB)/p$ в $\CH^*(G/B)/p$
состоит из рациональных циклов.
Обозначим через $e_i$ любой прообраз элемента $x_i$.
\item Из определения $J_p(\xi) = (j_1,\dots,j_r)$ мы знаем,
что $e_i^{p^{j_i}} + \mbox{добавка}$~--- тоже рациональный цикл.
У $\CH^*(G/B)/p$ как $\CH^*(BB)/p$-модуля есть система образующих
(точнее, базис над образом $\CH^*(BB)/p$), состоящая
из элементов вида $e^I=e_1^{j_1}\dots e_r^{i_r}$,
где $I = (i_1,\dots,i_r)$, и $0\leq I \leq p^K-1$
(то есть, $0\leq i_m \leq p^{k_m}-1$ для всех $m$).
\item Ищем рациональные циклы на $X\times X$ при помощи
метода общей точки:
\[
\begin{tikzcd}
\CH^*(X\times X) \arrow[->>]{r} \arrow{d}{\res}
& \CH^*(X_{F(X)}) \arrow{d}{\isom} \\
\CH^*(G/B\times G/B) \arrow[->>]{r}
& \CH^*((G/B)_{F(X)}).
\end{tikzcd}
\]
Возьмем какой-нибудь прообраз $\alpha_i\in\CH^*(X\times X)$
элемента $e_i\in\CH^*((G/B)_{F(X)})$,
и пусть $\ol{\alpha_i} = \res(\alpha_i)$~--- его образ
в $\CH^*(G/B\times G/B)$.
Тогда $\ol{\alpha_i} = e_i \times 1 + \mbox{добавка}$.
Более аккуратные рассуждения показывают, что
$\ol{\alpha_i} = e_i \times 1 + (\mbox{добавка}) - 1\times e_i$.
\item Пусть цикл $a\in\im(\CH^*(BB)/p)$ рационален.
Можно найти <<двойственный>> к нему
цикл $\alpha^{\vee} \in\im(\CH^*(BB)/p)$ такой, что
$\deg(e^{p^k-1}\cdot a \cdot a^{\vee}) = 1$.
То есть, на  $\im(\CH^*(BB)/p)$ есть невырожденная билинейная
форма со значениями в $\mathbb Z/p$.
Тогда цикл
\[
q = \alpha^{p^J-1} \cdot
\left(e^{p^JM}a\times e^{p^J(p^{K-J}-1-M)}a^\vee\right)
\]
рационален, поскольку $\alpha^{p^J-1}$ рационален,
и степени $e^{p^J}$ рациональны (по определению
$J$-инварианта).
Нетрудно проверить, что $q$~--- проектор
(если проигнорировать добавки).
Построение $q$ зависит от произвольных $M$ и $a$ из образа.
Если брать все возможные $M$ такие, что $0\leq M \leq p^{K-J}-1$,
и $a$ из базиса образа, получается
набор попарно ортогональных проекторов, которые в сумме
дают диагональ.

Упражнение: некоторая степень любого элемента конечного моноида
является идемпотентом
(на самом деле, нужно брать не такую степень, а еще большую
степень этого, чтобы убить добавки).

Посчитаем $q\cdot q$.
Можно считать, что $\alpha = e\times 1 - 1\times e$.
Тогда
$\alpha^{p^J-1} = \sum_{I} e^I \times e^{p^J-1-I}$,
и поэтому
\[
q = \sum_I e^{p^J M + I} a
\times e^{p^K-p^J M - 1 - I}a^{\vee}
= \deg(e^{p^K-1} a a^{\vee})(\dots).
\]
\end{enumerate}

\subsection{Группы типа $\mathsf{F}_4$}

Приведем пример многообразия, не расщепимого над общей точкой.
Будем обозначать расщепимые октонионы через $\mathbb{O}$,
а их компактную форму (над $\mathbb{R}$) через $C$.
Напомним, что над $\mathbb{R}$
бывают такие группы типа $\mathsf{F}_4$:
\begin{enumerate}
\item $H_3(\mbox{компактная форма},C)$:
$\begin{tikzpicture}[scale=1.0,thin]
\def\sm{0.7}
\coordinate (1) at ($\sm*(0, 0)$);
\coordinate (2) at ($\sm*(1.4, 0)$);
\coordinate (3) at ($\sm*(2.8, 0)$);
\coordinate (4) at ($\sm*(4.2, 0)$);
\draw (1)--(2);
\draw (3)--(4);
\draw (2) edge[transform canvas={yshift=2pt}] (3);
\draw (2) edge[transform canvas={yshift=-2pt}] (3);
\draw ($\sm*(2.0, 0.2)$) -- ($\sm*(2.2, 0)$) -- ($\sm*(2.0, -0.2)$);
\foreach \point in
{1,2,3,4}
  {
     \fill [white] (\point) circle (2.0pt);
     \draw [black] (\point) circle (2.0pt);
  }
\end{tikzpicture}$

\item $H_3(\mathbb{O})$:
$\begin{tikzpicture}[scale=1.0,thin]
\def\sm{0.7}
\coordinate (1) at ($\sm*(0, 0)$);
\coordinate (2) at ($\sm*(1.4, 0)$);
\coordinate (3) at ($\sm*(2.8, 0)$);
\coordinate (4) at ($\sm*(4.2, 0)$);
\draw (1)--(2);
\draw (3)--(4);
\draw (2) edge[transform canvas={yshift=2pt}] (3);
\draw (2) edge[transform canvas={yshift=-2pt}] (3);
\draw ($\sm*(2.0, 0.2)$) -- ($\sm*(2.2, 0)$) -- ($\sm*(2.0, -0.2)$);
\foreach \point in
{1,2,3,4}
  {
     \fill [white] (\point) circle (2.0pt);
     \draw [black] (\point) circle (2.0pt);
     \draw [black] (\point) circle (5.0pt);
  }
\end{tikzpicture}$

\item $H_3(\mbox{гиперболическая форма}, C)$:
$\begin{tikzpicture}[scale=1.0,thin]
\def\sm{0.7}
\coordinate (1) at ($\sm*(0, 0)$);
\coordinate (2) at ($\sm*(1.4, 0)$);
\coordinate (3) at ($\sm*(2.8, 0)$);
\coordinate (4) at ($\sm*(4.2, 0)$);
\draw (1)--(2);
\draw (3)--(4);
\draw (2) edge[transform canvas={yshift=2pt}] (3);
\draw (2) edge[transform canvas={yshift=-2pt}] (3);
\draw ($\sm*(2.0, 0.2)$) -- ($\sm*(2.2, 0)$) -- ($\sm*(2.0, -0.2)$);
\foreach \point in
{1,2,3,4}
  {
     \fill [white] (\point) circle (2.0pt);
     \draw [black] (\point) circle (2.0pt);
  }
\draw [black] (4) circle (5.0pt);
\end{tikzpicture}
$
\end{enumerate}
Что нарисовано в последнем случае?
Существует поле $F$ (например, $F=\mathbb{R}$) и группа типа
$\mathsf{F}_4$ над ним такие, что параболическая подгруппа
типа $P_4$ определена над $\mathbb{R}$, а остальные~--- нет.
Пусть $\xi\in H^1(F,\mathsf{F}_4)$~--- коцикл,
$X = {}_{\xi}(\mathsf{F}_4/P_4)$.
Тогда $\xi_{F(X)}$ может дать группу с индексом
$\begin{tikzpicture}[scale=1.0,thin]
\def\sm{0.7}
\coordinate (1) at ($\sm*(0, 0)$);
\coordinate (2) at ($\sm*(1.4, 0)$);
\coordinate (3) at ($\sm*(2.8, 0)$);
\coordinate (4) at ($\sm*(4.2, 0)$);
\draw (1)--(2);
\draw (3)--(4);
\draw (2) edge[transform canvas={yshift=2pt}] (3);
\draw (2) edge[transform canvas={yshift=-2pt}] (3);
\draw ($\sm*(2.0, 0.2)$) -- ($\sm*(2.2, 0)$) -- ($\sm*(2.0, -0.2)$);
\foreach \point in
{1,2,3,4}
  {
     \fill [white] (\point) circle (2.0pt);
     \draw [black] (\point) circle (2.0pt);
  }
\draw [black] (4) circle (5.0pt);
\end{tikzpicture}$.
По модулю $2$ есть два инварианта:
\begin{itemize}
\item $f_3\colon H^1(F, \mathsf{F}_4) \to H^3(F, \mathbb{Z}/2)$;
\item $f_5\colon H^1(F, \mathsf{F}_4) \to H^5(F, \mathbb{Z}/2)$.
\end{itemize}
Если у поля нет расширений нечетной степени, то это чистые формы.
Поэтому есть $3$-кратная и $5$-кратная формы Пфистера,
ассоциированные с $\xi$.

Пусть $J$~--- 27-мерная йорданова алгебра.
Предполагаем, что она приведенная (reduced).
Наш коцикл $\xi\in H^1(F,\mathsf{F}_4)$
удовлетворяет следующим равносильным условиям:
\begin{enumerate}
\item 
Каноническое отображение $H^1(F,\mathsf{F}_4) \to H^1(F,\mathsf{E}_6)$
переводит наш коцикл $\xi$
в изотропную группу с индексом Титса
\[
\begin{tikzpicture}[scale=1.0,thin]
\def\sm{0.7}
\coordinate (1) at ($\sm*(0, 0)$);
\coordinate (3) at ($\sm*(1.4, 0)$);
\coordinate (4) at ($\sm*(2.8, 0)$);
\coordinate (2) at ($\sm*(2.8, -1.4)$);
\coordinate (5) at ($\sm*(4.2, 0)$);
\coordinate (6) at ($\sm*(5.6, 0)$);
\draw (1)--(3)--(4)--(5)--(6);
\draw (2)--(4);
\foreach \point in
{1,2,3,4,5,6}
  {
     \fill [white] (\point) circle (2.0pt);
     \draw [black] (\point) circle (2.0pt);
  }
\draw [black] (1) circle (5.0pt);
\draw [black] (6) circle (5.0pt);
\end{tikzpicture}
\]
\item $g_3(\xi)=0$, где
$g_3\colon H^1(F,\mathsf{F}_4) \to H^3(F,\mu_3^{\otimes 2})$.
Этого всегда можно добиться кубическим расширением;
а если нет расширений нечетной степени, то это заведомо так.
\item Пусть $Q_x(y) = xyx$~--- квадратичная операция в $J$.
Элемент $e\in J$ называется \term{идемпотентом ранга $1$},
если $Q_e(J)\leq Fe$. Это равносильно тому, что
$N(e)=0$ и $e^{\sharp}=0$ (или тому, что $N(e,e,x)=0$ для всех
$x\in J$).
Рассмотрим многообразие
\[
\{\la e\ra \mid e\mbox{ --- идемпотент ранга $1$}\} =
{}_{\xi}(\mathsf{E}_6/P_6).
\]
Условие состоит в том, что у него есть рациональная точка.
\end{enumerate}
Поэтому внутри ${}_{\xi}(\mathsf{E}_6/P_6)$ есть $\mathbb{G}_m$,
которая дает разложение вида
$J = Fe\oplus U\oplus V$ по весовым пространствам $\mathbb{G}_m$.
При этом $\dim U = 16$, $\dim V = 10$.
Будем записывать элементы $J$ как тройки
вида $(\alpha,b,c)$ в соответствии с этим разложением.
Заметим, что ${}_{\xi}(\mathsf{E}_6/P_6)$~--- замкнутое
подмногообразие в $\mathbb{P}(J)$.
Следующая фильтрация будет $\Spin_{10}$-инвариантной:
\[
\begin{tikzcd}[column sep=0.4em]
{}_{\xi}(\mathsf{E}_6/P_6)
& \arrow{d}{\mathbb{A}^{16}}
& \{\alpha=0\} \arrow[left hook->]{ll}
& \arrow{d}{\mathbb{A}^{5}}
& \{\alpha=0,b=0\} = \Spin_{10}/P_1\arrow[left hook->]{ll} \\
& \{b=0,c=0\} = \pt
&
& \{\alpha=0,c=0\} = \Spin_{10}/P_5
\end{tikzcd}
\]
Можно начать с другого конца и получить такую фильтрацию:
\[
\begin{tikzcd}[column sep=0.4em]
{}_{\xi}(\mathsf{E}_6/P_6)
& \arrow{d}\mathbb{A}^{8}
& \{c=0\} \arrow[left hook->]{ll}
& \arrow{d}\mathbb{A}^{1}
& \{c=0,b=0\} = \pt \arrow[left hook->]{ll} \\
& \{b=0,\alpha=0\} = \Spin_{10}/P_1
&
& \{c=0,\alpha=0\} = \Spin_{10}/P_5
\end{tikzcd}
\]
Кстати, $\Spin_{10}/P_1$~--- это восьмимерная пфистерова квадрика,
которая и дает отображение $f_3$.

Если у ${}_{\xi}(\mathsf{E}_6/P_6)$ есть рациональная точка,
то есть большая клетка, и оно бирационально эквивалентно
проективному пространству:
\begin{eqnarray*}
{}_{\xi}(\mathsf{E}_6/P_6) & \dashrightarrow
& \mathbb{P}(Fe\oplus U),\\
{}[\alpha:b:c] & \mapsto & [\alpha:b],\\
{}[\alpha^2:\alpha b:Q_b(e)] & \mapsfrom & [\alpha:b].
\end{eqnarray*}
Локус стрелки $[\alpha:b:c]\mapsto [\alpha:b]$~--- восьмимерная
квадрика.
Локус стрелки $[\alpha:b]\mapsto [\alpha^2:\alpha b:Q_b(e)]$~---
это $\{\alpha=0, Q_b(e)=0\} = \Spin_{10}/P_5$~--- скрученная
форма максимального ортогонального грассманиана.

Многообразие ${}_{\xi}(\mathsf{F}_4/P_4)$ живет внутри
${}_{\xi}(\mathsf{E}_6/P_6)$ как гиперплоское сечение
вида $\alpha + t(c) = 0$, где $t$~--- линейная форма на $V$
(напомним, что $\dim V = 10$).
Давайте пересечем все с этим уравнением.
Получим бирациональное отображение
\[
{}_{\xi}(\mathsf{F}_4/P_4) \dashrightarrow Q.
\]
Локус слева~--- семимерная квадрика $\Spin_9/P_1$.
Локус справа~--- пятнадцатимерная квадрика
$\{\alpha^2 + t(Q_b(e)) = 0\} = \Spin_{10}/P_5 = \Spin_9 / P_4$.
Здесь $Q$~--- норменная квадрика:
мы начали с формы Пфистера $\lAngle a,b,c,d,e\rAngle$.
Семнадцатимерная
квадратичная форма $\la a\ra\perp \lAngle b,c,d,e\rAngle$~---
инвариант этой формы.
Ассоциированная с ней пятнадцатимерная квадрика и есть $Q$.

Что дальше?
Оказывается, $\Bl_{\Spin_9/P_1}({}_{\xi}(\mathsf{F}_4/P_4)) \isom
\Bl_{\Spin_9/P_4}(Q)$.
Для мотива раздутия есть формула вида
\[
M(\Bl_X(Y)) = M(Y) \oplus ( M(X)\otimes M(\mathbb{P}(N_XY))\{\dots\}).
\]
Мотив норменной квадрики раскладывается в один экземпляр
мотива Роста $R(\lAngle a,b,c,d,e\rAngle)$ (он как раз
пятнадцатимерный) и $R(\lAngle b,c,d,e\rAngle)$
в нужном количестве (7 штук):
\[
\begin{tikzpicture}[scale=1.0,thin,font=\scriptsize]
\def\sm{0.7}
\coordinate (0) at ($\sm*(0, 0)$);
\coordinate (1) at ($\sm*(1.4, 0)$);
\coordinate (2) at ($\sm*(2.8, 0)$);
\coordinate (3) at ($\sm*(4.2, 0)$);
\coordinate (4) at ($\sm*(5.6, 0)$);
\coordinate (5) at ($\sm*(7.0, 0)$);
\coordinate (6) at ($\sm*(8.4, 0)$);
\coordinate (7) at ($\sm*(9.8, 0)$);
\coordinate (8) at ($\sm*(11.2, 0)$);
\coordinate (9) at ($\sm*(12.6, 0)$);
\coordinate (10) at ($\sm*(14.0, 0)$);
\coordinate (11) at ($\sm*(15.4, 0)$);
\coordinate (12) at ($\sm*(16.8, 0)$);
\coordinate (13) at ($\sm*(18.2, 0)$);
\coordinate (14) at ($\sm*(19.6, 0)$);
\coordinate (15) at ($\sm*(21.0, 0)$);
\coordinate (p1) at ($\sm*(6.8, 1)$);
\coordinate (p2) at ($\sm*(7.0, 1)$);
\coordinate (p3) at ($\sm*(7.2, 1)$);
\coordinate (p4) at ($\sm*(15.2, 1)$);
\coordinate (p5) at ($\sm*(15.4, 1)$);
\coordinate (p6) at ($\sm*(15.6, 1)$);
\draw (0)--(1)--(2)--(3)--(4)--(5)--(6)--(7)--(8)--(9)--(10)--(11)--(12)--(13)--(14)--(15);
\draw[dashed] (0) edge[bend left=45](15);
\draw[dotted] (1) edge[bend left=35](8);
\draw[dotted] (2) edge[bend left=35](9);
\draw[dotted] (7) edge[bend left=35](14);
\node at (0) [below=5pt] {$0$};
\node at (1) [below=5pt] {$1$};
\node at (2) [below=5pt] {$2$};
\node at (3) [below=5pt] {$3$};
\node at (4) [below=5pt] {$4$};
\node at (5) [below=5pt] {$5$};
\node at (6) [below=5pt] {$6$};
\node at (7) [below=5pt] {$7$};
\node at (8) [below=5pt] {$8$};
\node at (9) [below=5pt] {$9$};
\node at (10) [below=5pt] {$10$};
\node at (11) [below=5pt] {$11$};
\node at (12) [below=5pt] {$12$};
\node at (13) [below=5pt] {$13$};
\node at (14) [below=5pt] {$14$};
\node at (15) [below=5pt] {$15$};
\foreach \point in
{0,1,2,3,4,5,6,7,8,9,10,11,12,13,14,15}
  {
     \fill [white] (\point) circle (2.0pt);
     \draw [black] (\point) circle (2.0pt);
  }
\foreach \point in
{p1,p2,p3,p4,p5,p6}
  {
     \fill [black] (\point) circle (0.7pt);
  }
\end{tikzpicture}
\]
Мотив норменной квадрики $\Spin_9/P_1$ раскладывается
в $R(\lAngle b,c,d,e\rAngle)$ (1 штука) и мотивы Роста,
соответствующие $f_3$.
Многообразие $\Spin_9/P_4$ расщепимо над общей точкой.
Мотив ${}_{\xi}(\mathsf{F}_4/P_4)$ выглядит так:
\[
\begin{tikzpicture}[scale=1.0,thin,font=\scriptsize]
\def\sm{0.7}
\coordinate (1) at ($\sm*(0, 0)$);
\coordinate (2) at ($\sm*(1.4, 0)$);
\coordinate (3) at ($\sm*(2.8, 0)$);
\coordinate (4) at ($\sm*(4.2, 0)$);
\coordinate (5) at ($\sm*(5.2, 1)$);
\coordinate (6) at ($\sm*(6.6, 1)$);
\coordinate (7) at ($\sm*(8.0, 1)$);
\coordinate (8) at ($\sm*(9.4, 1)$);
\coordinate (9) at ($\sm*(10.8, 1)$);
\coordinate (10) at ($\sm*(12.2, 1)$);
\coordinate (11) at ($\sm*(13.6, 1)$);
\coordinate (12) at ($\sm*(15.0, 1)$);
\coordinate (5x) at ($\sm*(5.2, -1)$);
\coordinate (6x) at ($\sm*(6.6, -1)$);
\coordinate (7x) at ($\sm*(8.0, -1)$);
\coordinate (8x) at ($\sm*(9.4, -1)$);
\coordinate (9x) at ($\sm*(10.8, -1)$);
\coordinate (10x) at ($\sm*(12.2, -1)$);
\coordinate (11x) at ($\sm*(13.6, -1)$);
\coordinate (12x) at ($\sm*(15.0, -1)$);
\coordinate (13) at ($\sm*(16.0, 0)$);
\coordinate (14) at ($\sm*(17.4, 0)$);
\coordinate (15) at ($\sm*(18.8, 0)$);
\coordinate (16) at ($\sm*(20.2, 0)$);
\node at ($\sm*(0, -2)$) {$0$};
\node at ($\sm*(1.4, -2)$) {$1$};
\node at ($\sm*(2.8, -2)$) {$2$};
\node at ($\sm*(4.2, -2)$) {$3$};
\node at ($\sm*(5.2, -2)$) {$4$};
\node at ($\sm*(6.6, -2)$) {$5$};
\node at ($\sm*(8.0, -2)$) {$6$};
\node at ($\sm*(9.4, -2)$) {$7$};
\node at ($\sm*(10.8, -2)$) {$8$};
\node at ($\sm*(12.2, -2)$) {$9$};
\node at ($\sm*(13.6, -2)$) {$10$};
\node at ($\sm*(15.0, -2)$) {$11$};
\node at ($\sm*(16.0, -2)$) {$12$};
\node at ($\sm*(17.4, -2)$) {$13$};
\node at ($\sm*(18.8, -2)$) {$14$};
\node at ($\sm*(20.2, -2)$) {$15$};
\draw (1)--(2)--(3)--(4)--(5)--(6)--(7)--(8)--(9)--(10)--(11)--(12)--(13)--(14)--(15)--(16);
\draw (4)--(5x)--(6x)--(7x)--(8x)--(9x)--(10x)--(11x)--(12x)--(13);
\draw[dashed] (1) edge[bend left=45](16);
\draw[dotted] (2) edge[bend left=45](5);
\draw[dotted] (3) edge[bend left=45](6);
\draw[dotted] (4) edge[bend right=35](7);
\draw[dotted] (8) edge[bend left=45](11);
\draw[dotted] (9) edge[bend left=45](12);
\draw[dotted] (10) edge[bend right=35](13);
\draw[dotted] (5x) edge[bend right=45](8x);
\draw[dotted] (6x) edge[bend right=45](9x);
\draw[dotted] (7x) edge[bend left=45](10x);
\draw[dotted] (11x) edge[bend right=45](14);
\draw[dotted] (12x) edge[bend right=45](15);
\foreach \point in
{1,2,3,4,5,6,7,8,9,10,11,12,5x,6x,7x,8x,9x,10x,11x,12x,13,14,15,16}
  {
     \fill [white] (\point) circle (2.0pt);
     \draw [black] (\point) circle (2.0pt);
  }
\end{tikzpicture}
\]
Предположим, что у ${}_{\xi}(\mathsf{F}_4/P_4)$ есть точка
над $F'$, где $F'/F$~--- расширение нечетной степени.
Верно ли, что у него есть точка над $F$?
Если есть два проективных гладких многообразия $Y_1,Y_2$,
которые бирационально эквивалентны, то
$Y_1(F)\neq\varnothing$ тогда и только тогда,
когда $Y_2(F)\neq\varnothing$.
Теперь из ${}_{\xi}(\mathsf{F}_4/P_4)(F')\neq\varnothing$
следует, что $Q(F')\neq\varnothing$,
и по теореме Спрингера $Q(F)\neq\varnothing$,
откуда ${}_{\xi}(\mathsf{F}_4/P_4)(F)\neq\varnothing$.

Мы получили картинку
\[
\begin{tikzcd}
\Bl_{\Spin_9/P_1}({}_{\xi}(\mathsf{F}_4/P_4)) \arrow{r}{\isom}\arrow{d}
& \Bl_{\Spin_9/P_4}Q\arrow{d} \\
{}_{\xi}(\mathsf{F}_4/P_4) & Q
\end{tikzcd}
\]

% 23.04.2012

\subsection{Многообразия, клеточные над общей точкой}

Вернемся к общей ситуации: по коциклу $\xi\in H^1(F,G)$, где
$G$~--- расщепимая группа, и простому числу $p$ мы построили
набор чисел $J_p(\xi) = (j_1,\dots,j_r)$.
Теорема Зайнуллина--Петрова--Семенова~\ref{thm:ZPS} утверждает,
что если $X$~--- ${}_{\xi}G$-однородное многообразие, клеточное
над $F(X)$, то $M(X)\otimes\mathbb{Z}/p = \bigoplus R_p(G)\{\dots\}$,
и
\[
P(R_p(G),t) = \prod\frac{1-t^{p^{j_i}d_i}}{1-t^{d_i}}.
\]
\begin{remark}
При помощи знания $J_p(\xi)$ для всех $p$ можно описать все такие $X$.
\end{remark}
Выше мы встречали точную последовательность вида
\[
\CH^*(BB) \to \CH^*(G/B) \to CH^*(G) \to 0.
\]
Пусть $X = {}_{\xi}(G/P)$.
Попробуем нарисовать аналогичную последовательность для $P$:
\[
\begin{tikzcd}
\CH^*(BP) \arrow{d} \\
\CH^*(G/P) \arrow{r} & \CH^*(G)\arrow{r} & \CH^*(P)\arrow{r} & 0
\end{tikzcd}
\]
Нижняя строка этой диаграммы является точною последовательностью
градуированных колец.
Однако, точность всей последовательности в члене $\CH^*(G/P)$
неизвестна.
При этом $\CH^*(P) \isom \CH^*(L)$ (поскольку $P/U\isom L$,
а $U$ аффинно).
Более того, мы утверждаем, что $\CH^*(L) \isom \CH^*([L,L])$.
Почему это так?
Рассмотрим коммутативную диаграмму
\[
\begin{tikzcd}
\CH^*(BB_L) \arrow{r}\arrow{d}
& \CH^*(L/B) \arrow{r}\arrow{d}{\isom}
& \CH^*(L) \arrow{d}\arrow{r}
& 0 \\
\CH^*(BB_{[L,L]}) \arrow{r}
& \CH^*(L/B) \arrow{r}
& \CH^*([L,L]) \arrow{r}
& 0
\end{tikzcd}
\]
Здесь $B_L$ обозначает борелевскую подгруппу в $L$, а
$B_{[L,L]}$~--- борелевскую в $[L,L]$.
Левая вертикальная стрелка не обязана быть изоморфизмом,
однако образы $\CH^*(BB_L)$ и $\CH^*(BB_{[L,L]})$ в
$\CH^*(L/B)$ совпадают.
Это утверждение достаточно проверить на $\CH^1$.

Обозначим $H = [L,L]$.
Мы получили точную последовательность
\[
\CH^*(G/P)/p \to \CH^*(G)/p \xrightarrow{\ph} \CH^*(H)/p \to 0.
\]
Нам известно, что
$\CH^*(G)/p = (\mathbb{Z}/p)[x_1,\dots,x_r]/(x_i)^{p^{k_i}}$
и
$\CH^*(H)/p = (\mathbb{Z}/p)[y_1,\dots,y_s]/(y_j)^{p^{l_j}}$.

Оказывается, есть отображение $\sigma\colon\{1,\dots,s\}\to\{1,\dots,r\}$
такое, что
\[
\ph(x_{\sigma(m)}) = c\cdot y_m + \mbox{члены меньшего порядка},
\]
где $c\in (\mathbb{Z}/p)^*$~--- некоторая константа.

\begin{theorem}\label{thm:cellularity}
Пусть $G$~--- простая расщепимая группа,
$\xi\in H^1(F,G)$, $X = {}_{\xi}(G/P)$, $Y = {}_{\xi}(G/B)$.
Тогда следующие условия эквивалентны.
\begin{enumerate}
\item $X$~--- клеточное над общей точкой.
\item Композиция отображений
\[
\CH^*({}_{\xi}(G/B)) \xrightarrow{\res} \CH^*(G/B) \to \CH^*(G) \to
\CH^*(P)
\]
сюръективна.
\item Для любого простого $p$ выполняются условия
\begin{enumerate}
\item $j_{\sigma(m)}(\xi) = 0$, если $m$ такое, что $\deg(y_m) > 1$;
\item\label{item:thm-mod-p} композиция отображений
\[
\CH^1({}_{\xi}(G/B))/p \xrightarrow{\res} \CH^1(G/B)/p \to \CH^1(G)/p \to
\CH^1(P)/p
\]
сюръективна.
\end{enumerate}
\end{enumerate}
\end{theorem}
\begin{example}
Пусть $\xi\in H^1(F,\mathsf{E}_7)$, $X = {}_{\xi}(\mathsf{E}_7/P_7)$.
Когда $X$ расщепимо над общей точкой?
\begin{itemize}
\item $p=2$. В этом случае
\[
\begin{tikzpicture}[remember picture]
\node (a1) {$\CH^*(\mathsf{E}_7)/2 = (\underset{\deg:}{\mathbb{Z}/2})
[\underset{1}{x_1}, \subnode{x1}{$\underset{3}{x_2}$}, \underset{5}{x_3},
\underset{9}{x_4}]/(\dots)$};
\node (b1) [below=of a1] {$\CH^*(\mathsf{E}_6)/2 =
(\underset{\deg:}{\mathbb{Z}/2})[\subnode{y1}{$\underset{3}{y_1}$}]/(y_1^2)$};
\draw[|->] (x1.south) -- (y1.north);
\end{tikzpicture}
\]
Условие~(\ref{item:thm-mod-p}) теоремы~\ref{thm:cellularity}
превращается в $j_2 = 0$.
Элементы $x_3$, $x_4$ получаются операцией Стинрода, и потому
$j_3 = j_4 = 0$ автоматически.
\item $p=3$. В этом случае
\end{itemize}
\[
\begin{tikzpicture}[remember picture]
\node (a2) {$\CH^*(\mathsf{E}_7)/3 = (\underset{\deg:}{\mathbb{Z}/3})
[\subnode{x2}{$\underset{4}{x_1}$}]/(x_1^3)$};
\node (b2) [below=of a2] {$\CH^*(\mathsf{E}_6)/3 =
(\underset{\deg:}{\mathbb{Z}/3})[\subnode{y2}{$\underset{4}{y_1}$}]/(y_1^3)$};
\draw[|->] (x2.south) -- (y2.north);
\end{tikzpicture}
\]
Условие состоит в том, что $j_1=0$; и тогда весь $j$-инвариант для модуля $3$
равен нулю.
\end{example}
Задача: посчитать $j_2$ для компактной формы $\mathsf{E}_7$ над $\mathbb{R}$.
Эквивалентно:
\begin{equation}\label{eqn:condition-e7}
\begin{array}{p{0.8\textwidth}}
верно ли, что над полем $\mathbb{R}(\SB(\mathbb{H}))$
(которое изоморфно $\Frac(\mathbb{R}[x,y]/(x^2+y^2-1))$)
компактная форма $\mathsf{E}_7$ расщепляется?
\end{array}
\end{equation}
Еще одна эквивалентная формулировка: пусть $R_2(\xi)$~--- мотив Роста,
отвечающий кватернионам.
Сравним ${}_{\xi}(\mathsf{E}_7/P_7)$ и $\SB(\mathbb{H})$.
Над $F({}_{\xi}(\mathsf{E}_7/P_7))$ алгебра Титса ($\mathbb{H}$)
тривиальна, и потому $\SB(\mathbb{H})$ приобретает рациональную точку.
Если условие~\ref{eqn:condition-e7} выполняется, то верно и обратное:
над $F(\SB(\mathbb{H}))$ многообразие ${}_{\xi}(\mathsf{E}_7/P_7)$
имеет рациональную точку. Почему это так?
\begin{theorem}
Пусть $X,Y$~--- проективные однородные многообразия (возможно,
относительно разных групп), $p$~--- простое число.
Предположим, что $X_{F(Y)}$ и $Y_{F(X)}$ имеют рациональные точки.
Тогда $M(X)\otimes(\mathbb{Z}/p)$ и $M(Y)\otimes(\mathbb{Z}/p)$
имеют общее слагаемое, <<задевающее>> точку
(то есть, $q\in\CH^{\dim X}(X\times X)$~--- проектор,
и $\im(q)$ над замыканием содержит $\pt\in\CH^{\dim X}(X)$).
\end{theorem}

\begin{proposition}
Если есть изотропная ${}_{\xi}G$, и анизотропное ядро типа $H$,
то $\xi$ приходит из $\zeta\in H^1(F,H)$.
Тогда $J_p(\xi)$ выражается через $J_p(\zeta)$:
\[
j_i(\xi) = \begin{cases} j_m(\zeta), & \text{если $i=\sigma(m)$
для некоторого $m$};\\
0, & \text{иначе}.
\end{cases}
\]
Поэтому имеет смысл считать $J_p(\xi)$ только для анизотропных
групп (то есть, над $\mathbb{R}$~--- только для компактных форм).
\end{proposition}
\begin{itemize}
\item $\mathsf{G}_2$, $\mathsf{F}_4$:
$(\underset{\deg:}{\mathbb{Z}/2})[\underset{3}{x_1}]/(x_1^2)$,
$J_2(\xi) = (1)$~--- а не $0$, ибо у $\mathbb{R}$ нет расширений нечетной
степени. При этом $R_2(\xi)$~--- мотив Роста $\lAngle -1,-1,-1\rAngle$.
\item $\mathsf{E}_6$: компактная форма есть, но она внешняя.
\item $\mathsf{E}_7$: $(\underset{\deg:}{\mathbb{Z}/2})
[\underset{1}{x_1}, \underset{3}{x_2}, \underset{5}{x_3}, \underset{9}{x_4}]$.
\end{itemize}
\begin{proposition}\label{prop:j-inv-for-e7}
В этом случае $J_2(\xi) = (1, 0, 0, 0)$, $R_2(\xi)$~--- это мотив Роста
квадрики $\lAngle -1, -1\rAngle$.
\end{proposition}
Для доказательства этого утверждения нужен \term{инвариант Роста}.

\subsection{Инвариант Роста}
Пусть $G$~--- простая, односвязная, но не обязательно расщепимая группа.
Тогда есть инвариант
\[
r\colon H^1(F,G) \to H^3(F, (\mathbb{Q}/\mathbb{Z})(2))
= \varinjlim_{N} H^3(F, \mu_N^{\otimes 2}),
\]
который порождает всю абелеву группу таких инвариантов.
Вместо предела в правой части можно взять одно достаточно большое $N$.
Заметим, что $r(\xi)$ всегда лежит в кручении.
Посмотрим на наименьшее $N_G$ такое, что всегда $N_G\cdot r(\xi) = 0$
(для всех расширений $E/F$ и для всех $\xi$).
Тогда $N_G$ зависит только от типа $G$.
А именно, для исключительных групп
$N_{\mathsf{G}_2} = 2$, $N_{\mathsf{F}_4} = N_{\mathsf{E}_6} = 6$,
$N_{\mathsf{E}_7} = 12$, $N_{\mathsf{E}_8} = 60$.
Для $\mathsf{G}_2$ мы это видели:
$\mathsf{G}_2 = \Aut(\mathbb{O})$, и $\mathbb{O}$ определяется формой
$\lAngle a, b, c\rAngle \in H^3(F,\mathbb{Z}/2)$.
Для $\mathsf{F}_4$, $\mathsf{E}_6$: $f_3 = 3r\in H^3(F, \mathbb{Z}/2)$,
$g_3 = 2r \in H^3(F,\mathbb{Z}/3)$, где
$r\in H^3(F,\mathbb{Z}/6)$.

\begin{exercise}[Исследовательская задача]
Придумать формулу для $r\colon H^1(F,G) \to H^3(f, \mathbb{Z}/4)$,
где $G$~--- односвязная группа типа $\mathsf{D}_6$
с индексом Титса
\[
\begin{tikzpicture}[scale=1.0,thin]
\def\sm{0.7}
\coordinate (1) at ($\sm*(0, 0)$);
\coordinate (2) at ($\sm*(1.4, 0)$);
\coordinate (3) at ($\sm*(2.8, 0)$);
\coordinate (4) at ($\sm*(4.2, 0)$);
\coordinate (5) at ($\sm*(5.2, 1)$);
\coordinate (6) at ($\sm*(5.2, -1)$);
\draw (1)--(2)--(3)--(4)--(5);
\draw (4)--(6);
\foreach \point in
{1,2,3,4,5,6}
  {
     \fill [white] (\point) circle (2.0pt);
     \draw [black] (\point) circle (2.0pt);
  }
\foreach \point in
{2,4,5}
  {
     \draw [black] (\point) circle (5.0pt);
  }
\end{tikzpicture}
\]
Такая группа задается кватернионами.
\end{exercise}
\begin{theorem}[Черноусов--Гарибальди]
Если $G$~--- расщепимая группа типа $\mathsf{G}_2$, $\mathsf{F}_4$,
$\mathsf{E}_6$ или $\mathsf{E}_7$, то ядро инварианта Роста тривиально.
\end{theorem}
Пусть $\xi\in H^1(F,G)$.
Хотим узнать, тривиален ли $\xi$.
Теорема Черноусова--Гарибальди говорит, что (оказывается!)
достаточно посчитать $r(\xi)$.
Для $\mathsf{E}_8$ это неверно: можно взять
$\xi\in H^1(\mathbb{R},\mathsf{E}_8)$, задающий компактную форму.

Воспользуемся этой теоремой для доказательства
утверждения~\ref{prop:j-inv-for-e7}.
Достаточно доказать, что $\xi_{F(SB(\mathbb{H}))} = *$.
Действительно, если это так, то ${}_{\xi}(\mathsf{E}_7/B)$
и $\SB(\mathbb{H})$ имеют рациональные точки над полями функций
друг друга.
Тогда у них есть общее мотивное слагаемое
$R_2(\xi) = R_2(\lAngle -1, -1\rAngle)$, и потому $J_2(\xi) = (1,0,0,0)$
из формулы для полинома Пуанкаре.

Мы должны представить $\xi$ как элемент $H^1(F,G')$,
где $G'$ односвязна.
Расщепимая $G'$ не подойдет: $\xi$ приходит из
$H^1(F, \mathsf{E}_7^{\operatorname{ad}})$,
но не из $H^1(F, \mathsf{E}_7^{\operatorname{sc}})$
(поскольку алгебра Титса нетривиальна).
Возьмем в качестве $G'$ группу с индексом Титса
\[
\begin{tikzpicture}[scale=1.0,thin]
\def\sm{0.7}
\coordinate (1) at ($\sm*(0, 0)$);
\coordinate (2) at ($\sm*(2.8, -1.4)$);
\coordinate (3) at ($\sm*(1.4, 0)$);
\coordinate (4) at ($\sm*(2.8, 0)$);
\coordinate (5) at ($\sm*(4.2, 0)$);
\coordinate (6) at ($\sm*(5.6, 0)$);
\coordinate (7) at ($\sm*(7.0, 0)$);
\draw (1)--(3)--(4)--(5)--(6)--(7);
\draw (2)--(4);
\foreach \point in
{1,2,3,4,5,6,7}
  {
     \fill [white] (\point) circle (2.0pt);
     \draw [black] (\point) circle (2.0pt);
  }
\foreach \point in
{1,3,4,6}
  {
     \draw [black] (\point) circle (5.0pt);
  }
\end{tikzpicture}
\]
Тогда $\xi\in H^1(F,G')$.
Что можно сказать про $r(\xi)$?
Мы знаем, что $\xi$ расщепим над $\mathbb{C}$, и потому
$2r(\xi) = 0$.
Стало быть, $r(\xi) \in H^3(\mathbb{R}, \mathbb{Z}/2)$.
Значит, это либо $0$, либо $\lAngle -1, -1, -1\rAngle$.

Переходим на $F(\SB(\mathbb{H}))$.
Тогда $r(\xi)$ в любом случае становится $0$,
а $G'_{F(\SB(\mathbb{H}))}$ расщепима.
Поэтому можно применить теорему Черноусова--Гарибальди,
и заключить, что $\xi_{F(\SB(\mathbb{H}))} = *$,
чего мы и добивались.

Посмотрим теперь на компактную форму группы типа $\mathsf{E}_8$.
Мы знаем, что
\[
\CH^*(\mathsf{E}_8)/2 = (\underset{\deg:}{\mathbb{Z}/2})
[\underset{3}{x_1},\underset{5}{x_2},\underset{9}{x_3},\underset{15}{x_4}]
/(\dots).
\]
Инвариант Роста равен нулю; $j_1=0$, и потому $j_2 = j_3 = 0$.
Получаем, что $J_2(\xi) = (0, 0, 0, 1)$.
Отсюда следует, что $P(R_2(\xi), t) = 1 + t^{15}$
и $R_2(\xi) = R_2(\lAngle -1, -1, -1, -1, -1\rAngle)$.

Мораль: над $\mathbb{R}$ мотивы однородных проективных многообразий
(расщепимых над общей точкой) раскладываются в мотивы Роста
от пфистеровых форм с каким-то количеством $-1$;
количество зависит от типа группы.

Пусть $\xi\in H^1(F, G)$, где $G$~--- расщепимая группа типа $\mathsf{E}_6$.
Свяжем $J$-инвариант с индексом Титса.
Например, $J_3(\xi)$ определяется индексом Титса над ко-$3$-замыканием
поля $F$.

Что это значит?
Посмотрим на абсолютную группу Галуа $\operatorname{Gal}(\ol{F}/F)$
и возьмем в ней $3$-силовскую подгруппу.
Ей соответствует расширение $E/F$, которое называется
\term{ко-$3$-замыканием $F$}.
Неформально говоря, мы игнорируем расширения степеней, не делящихся на $3$.

\begin{tabular}{|c|c|c|c|c|c|}
\hline
$J_3(\mathsf{E}_6^{\operatorname{ad}})$ & $(0,0)$ & $(1,0)$ & $(0,1)$ & $(1,1)$ & $(2,1)$ \\
\hline
индекс Титса &
\begin{tikzpicture}[scale=0.5,thin]\def\sm{0.7}
\coordinate (1) at ($\sm*(0, 0)$);
\coordinate (2) at ($\sm*(2.8, -1.4)$);
\coordinate (3) at ($\sm*(1.4, 0)$);
\coordinate (4) at ($\sm*(2.8, 0)$);
\coordinate (5) at ($\sm*(4.2, 0)$);
\coordinate (6) at ($\sm*(5.6, 0)$);
\draw (1)--(3)--(4)--(5)--(6);
\draw (2)--(4);
\foreach \point in
{1,2,3,4,5,6}
  {
     \fill [white] (\point) circle (2.0pt);
     \draw [black] (\point) circle (2.0pt);
  }
\foreach \point in
{1,2,3,4,5,6}
  {
     \draw [black] (\point) circle (5.0pt);
  }
\end{tikzpicture}
&
\begin{tikzpicture}[scale=0.5,thin]\def\sm{0.7}
\coordinate (1) at ($\sm*(0, 0)$);
\coordinate (2) at ($\sm*(2.8, -1.4)$);
\coordinate (3) at ($\sm*(1.4, 0)$);
\coordinate (4) at ($\sm*(2.8, 0)$);
\coordinate (5) at ($\sm*(4.2, 0)$);
\coordinate (6) at ($\sm*(5.6, 0)$);
\draw (1)--(3)--(4)--(5)--(6);
\draw (2)--(4);
\foreach \point in
{1,2,3,4,5,6}
  {
     \fill [white] (\point) circle (2.0pt);
     \draw [black] (\point) circle (2.0pt);
  }
\foreach \point in
{2,4}
  {
     \draw [black] (\point) circle (5.0pt);
  }
\end{tikzpicture}
& \multicolumn{3}{|c|}{
  \begin{tikzpicture}[scale=0.5,thin]\def\sm{0.7}
\coordinate (1) at ($\sm*(0, 0)$);
\coordinate (2) at ($\sm*(2.8, -1.4)$);
\coordinate (3) at ($\sm*(1.4, 0)$);
\coordinate (4) at ($\sm*(2.8, 0)$);
\coordinate (5) at ($\sm*(4.2, 0)$);
\coordinate (6) at ($\sm*(5.6, 0)$);
\draw (1)--(3)--(4)--(5)--(6);
\draw (2)--(4);
\foreach \point in
{1,2,3,4,5,6}
  {
     \fill [white] (\point) circle (2.0pt);
     \draw [black] (\point) circle (2.0pt);
  }
\end{tikzpicture}
}\\
\hline
индекс алгебры Титса & $1$ & $3$ & $1$ & $3$ & $9$ или $27$\\
\hline
\end{tabular}

Приведем аналогичную таблицу для $J_2(\mathsf{E}_7^{\operatorname{sc}})$:

\begin{tabular}{|p{2.9cm}|p{2.5cm}|p{2.6cm}|p{2.9cm}|p{2.5cm}|}
\hline
\begin{center} $J_2(\mathsf{E}_7^{\operatorname{sc}})$ \end{center}
& \begin{center} $(0,0,0)$ \end{center}
& \begin{center} $(1,0,0)$ \end{center}
& \begin{center} $(1,1,0)$ \end{center}
& \begin{center} $(1,1,1)$ \end{center} \\
\hline
\begin{center}индекс Титса
над ко-$2$-замыканием\end{center} &
\begin{tikzpicture}[scale=0.5,thin]
\def\sm{0.7}
\coordinate (1) at ($\sm*(0, 0)$);
\coordinate (2) at ($\sm*(2.8, -1.4)$);
\coordinate (3) at ($\sm*(1.4, 0)$);
\coordinate (4) at ($\sm*(2.8, 0)$);
\coordinate (5) at ($\sm*(4.2, 0)$);
\coordinate (6) at ($\sm*(5.6, 0)$);
\coordinate (7) at ($\sm*(7.0, 0)$);
\draw (1)--(3)--(4)--(5)--(6)--(7);
\draw (2)--(4);
\foreach \point in
{1,2,3,4,5,6,7}
  {
     \fill [white] (\point) circle (2.0pt);
     \draw [black] (\point) circle (2.0pt);
  }
\foreach \point in
{1,2,3,4,5,6,7}
  {
     \draw [black] (\point) circle (5.0pt);
  }
\end{tikzpicture}
&
\begin{tikzpicture}[scale=0.5,thin]
\def\sm{0.7}
\coordinate (1) at ($\sm*(0, 0)$);
\coordinate (2) at ($\sm*(2.8, -1.4)$);
\coordinate (3) at ($\sm*(1.4, 0)$);
\coordinate (4) at ($\sm*(2.8, 0)$);
\coordinate (5) at ($\sm*(4.2, 0)$);
\coordinate (6) at ($\sm*(5.6, 0)$);
\coordinate (7) at ($\sm*(7.0, 0)$);
\draw (1)--(3)--(4)--(5)--(6)--(7);
\draw (2)--(4);
\foreach \point in
{1,2,3,4,5,6,7}
  {
     \fill [white] (\point) circle (2.0pt);
     \draw [black] (\point) circle (2.0pt);
  }
\foreach \point in
{1,6,7}
  {
     \draw [black] (\point) circle (5.0pt);
  }
\end{tikzpicture}
&
\begin{tikzpicture}[scale=0.5,thin]
\def\sm{0.7}
\coordinate (1) at ($\sm*(0, 0)$);
\coordinate (2) at ($\sm*(2.8, -1.4)$);
\coordinate (3) at ($\sm*(1.4, 0)$);
\coordinate (4) at ($\sm*(2.8, 0)$);
\coordinate (5) at ($\sm*(4.2, 0)$);
\coordinate (6) at ($\sm*(5.6, 0)$);
\coordinate (7) at ($\sm*(7.0, 0)$);
\draw (1)--(3)--(4)--(5)--(6)--(7);
\draw (2)--(4);
\foreach \point in
{1,2,3,4,5,6,7}
  {
     \fill [white] (\point) circle (2.0pt);
     \draw [black] (\point) circle (2.0pt);
  }
\foreach \point in
{1}
  {
     \draw [black] (\point) circle (5.0pt);
  }
\end{tikzpicture}
&
\begin{tikzpicture}[scale=0.5,thin]
\def\sm{0.7}
\coordinate (1) at ($\sm*(0, 0)$);
\coordinate (2) at ($\sm*(2.8, -1.4)$);
\coordinate (3) at ($\sm*(1.4, 0)$);
\coordinate (4) at ($\sm*(2.8, 0)$);
\coordinate (5) at ($\sm*(4.2, 0)$);
\coordinate (6) at ($\sm*(5.6, 0)$);
\coordinate (7) at ($\sm*(7.0, 0)$);
\draw (1)--(3)--(4)--(5)--(6)--(7);
\draw (2)--(4);
\foreach \point in
{1,2,3,4,5,6,7}
  {
     \fill [white] (\point) circle (2.0pt);
     \draw [black] (\point) circle (2.0pt);
  }
\end{tikzpicture}
\\
\hline
\begin{center}$r(\xi)$\end{center}
& \begin{center}$0$\end{center}
& \begin{center}чистый символ $\neq 0$ из $H^3(F,\mathbb{Z}/2)$:
$\lAngle a,b,c\rAngle$\end{center}
& \begin{center}сумма двух символов из $H^3(F,\mathbb{Z}/2)$
с общим слотом: $\lAngle a,b,c\rAngle + \lAngle a,d,e\rAngle$\end{center}
& \begin{center}иначе\end{center} \\
\hline
\end{tabular}

Случаи $\mathsf{E}_6\pmod{2}$, $\mathsf{E}_7\pmod{3}$~--- легкие упражнения.
Случаи $\mathsf{E}_8\pmod{2}$, $\mathsf{E}_7^{\operatorname{ad}}\pmod{2}$~---
исследовательское упражнение.
Случай $\mathsf{E}_8\pmod{5}$~--- легкое упражнение.
Вот ответ для $\mathsf{E}_8\pmod{3}$:

\begin{tabular}{|p{2.9cm}|p{3.0cm}|p{3.0cm}|p{3.0cm}|}
\hline
\begin{center} $J_3(\mathsf{E}_8)$ \end{center}
& \begin{center} $(0,0)$ \end{center}
& \begin{center} $(1,0)$ \end{center}
& \begin{center} $(1,1)$ \end{center} \\
\hline
\begin{center}индекс Титса\end{center}
&
\begin{tikzpicture}[scale=0.5,thin]
\def\sm{0.7}
\coordinate (1) at ($\sm*(0, 0)$);
\coordinate (2) at ($\sm*(2.8, -1.4)$);
\coordinate (3) at ($\sm*(1.4, 0)$);
\coordinate (4) at ($\sm*(2.8, 0)$);
\coordinate (5) at ($\sm*(4.2, 0)$);
\coordinate (6) at ($\sm*(5.6, 0)$);
\coordinate (7) at ($\sm*(7.0, 0)$);
\coordinate (8) at ($\sm*(8.4, 0)$);
\draw (1)--(3)--(4)--(5)--(6)--(7)--(8);
\draw (2)--(4);
\foreach \point in
{1,2,3,4,5,6,7,8}
  {
     \fill [white] (\point) circle (2.0pt);
     \draw [black] (\point) circle (2.0pt);
  }
\foreach \point in
{1,2,3,4,5,6,7,8}
  {
     \draw [black] (\point) circle (5.0pt);
  }
\end{tikzpicture}
&
\begin{tikzpicture}[scale=0.5,thin]
\def\sm{0.7}
\coordinate (1) at ($\sm*(0, 0)$);
\coordinate (2) at ($\sm*(2.8, -1.4)$);
\coordinate (3) at ($\sm*(1.4, 0)$);
\coordinate (4) at ($\sm*(2.8, 0)$);
\coordinate (5) at ($\sm*(4.2, 0)$);
\coordinate (6) at ($\sm*(5.6, 0)$);
\coordinate (7) at ($\sm*(7.0, 0)$);
\coordinate (8) at ($\sm*(8.4, 0)$);
\draw (1)--(3)--(4)--(5)--(6)--(7)--(8);
\draw (2)--(4);
\foreach \point in
{1,2,3,4,5,6,7,8}
  {
     \fill [white] (\point) circle (2.0pt);
     \draw [black] (\point) circle (2.0pt);
  }
\foreach \point in
{7,8}
  {
     \draw [black] (\point) circle (5.0pt);
  }
\end{tikzpicture}
&
\begin{tikzpicture}[scale=0.5,thin]
\def\sm{0.7}
\coordinate (1) at ($\sm*(0, 0)$);
\coordinate (2) at ($\sm*(2.8, -1.4)$);
\coordinate (3) at ($\sm*(1.4, 0)$);
\coordinate (4) at ($\sm*(2.8, 0)$);
\coordinate (5) at ($\sm*(4.2, 0)$);
\coordinate (6) at ($\sm*(5.6, 0)$);
\coordinate (7) at ($\sm*(7.0, 0)$);
\coordinate (8) at ($\sm*(8.4, 0)$);
\draw (1)--(3)--(4)--(5)--(6)--(7)--(8);
\draw (2)--(4);
\foreach \point in
{1,2,3,4,5,6,7,8}
  {
     \fill [white] (\point) circle (2.0pt);
     \draw [black] (\point) circle (2.0pt);
  }
\end{tikzpicture}
\\
\hline
\begin{center} $r(\xi)$ \end{center}
& \begin{center} $0$ \end{center}
& \begin{center} чистый символ $\neq 0$ из $H^3(F,\mathbb{Z}/3)$\end{center}
& \begin{center} иначе \end{center} \\
\hline
\end{tabular}

\end{document}

